\input{preamble}

% OK, start here.
%
\begin{document}

\title{Propriétés des espaces algébriques}


\maketitle

\phantomsection
\label{section-phantom}

\tableofcontents

\section{Introduction}
\label{section-introduction}

\noindent
Veuillez consulter
Spaces, section \ref{spaces-section-introduction},
pour une brève introduction aux espaces algébriques, et lire aussi
une partie de ce chapitre pour nos définitions et conventions de base
concernant les espaces algébriques. Dans le présent chapitre, nous introduisons
quelques notions et propriétés élémentaires des espaces algébriques. Une référence
fondamentale pour le cas des espaces algébriques quasi-séparés est
\cite{Kn}.

\medskip\noindent
L'exposé est parfois quelque peu malaisé, car nous avons choisi
de traiter d'abord des propriétés des espaces algébriques
en eux-mêmes, et seulement plus tard de celles des morphismes d'espaces
algébriques. Nous dérogeons à cette règle pour les
{\it morphismes étales} d'espaces algébriques, que nous introduisons à la
section \ref{section-etale-morphisms}.
Mais, jusqu'à cette section, lorsque nous disons qu'un morphisme possède une propriété donnée,
cela signifie automatiquement que sa source est un schéma (ou peut-être que
le morphisme est représentable).

\medskip\noindent
Certains résultats de ce chapitre (notamment ceux qui concernent les points)
seront améliorés dans le chapitre consacré aux espaces algébriques décents.


\section{Conventions}
\label{section-conventions}

\noindent
On suppose une fois pour toutes que tous les schémas appartiennent à
un grand site fppf $\Sch_{fppf}$. De plus, tous les anneaux $A$ considérés
sont tels que $\Spec(A)$ soit (isomorphe à) un
objet de ce grand site.

\medskip\noindent
Soient $S$ un schéma et $X$ un espace algébrique sur $S$.
Dans ce chapitre et le suivant, nous noterons $X \times_S X$
le produit de $X$ par lui-même (dans la catégorie des espaces
algébriques sur $S$), au lieu de $X \times X$. Nous voulons ainsi
éviter toute confusion lorsque nous changeons de schéma de base, comme dans
Spaces, section \ref{spaces-section-change-base-scheme}.


\section{Axiomes de séparation}
\label{section-separation}

\noindent
Dans cette section, nous réunissons toutes les conditions de séparation « absolues »
des espaces algébriques. Puisque, selon nos conventions, tout espace algébrique est un
espace algébrique sur un schéma de base bien déterminé, toute propriété absolue
de $X$ sur $S$ correspond à une condition imposée à $X$ considéré
comme espace algébrique sur $\Spec(\mathbf{Z})$. En voici la
formulation précise.

\begin{definition}
\label{definition-separated}
(Comparer avec Spaces, définition \ref{spaces-definition-separated}.)
Considérons un grand site fppf
$\Sch_{fppf} = (\Sch/\Spec(\mathbf{Z}))_{fppf}$.
Soit $X$ un espace algébrique sur
$\Spec(\mathbf{Z})$. Soit $\Delta : X \to X \times X$
le morphisme diagonal.
\begin{enumerate}
\item On dit que $X$ est {\it séparé} si $\Delta$ est une immersion fermée.
\item On dit que $X$ est {\it localement séparé}\footnote{Dans la
littérature, cela désigne souvent des espaces algébriques quasi-séparés et
localement séparés.} si $\Delta$ est une
immersion.
\item On dit que $X$ est {\it quasi-séparé} si $\Delta$ est quasi-compact.
\item On dit que $X$ est {\it localement quasi-séparé pour la topologie de Zariski}\footnote{
Cette notion a été proposée par B.\ Conrad.} s'il
existe un recouvrement de Zariski $X = \bigcup_{i \in I} X_i$ (voir Spaces,
définition \ref{spaces-definition-Zariski-open-covering}) tel que
chaque $X_i$ soit quasi-séparé.
\end{enumerate}
Soit $S$ un schéma appartenant à $\Sch_{fppf}$, et soit
$X$ un espace algébrique sur $S$. On dit alors que $X$ est {\it séparé},
{\it localement séparé}, {\it quasi-séparé} ou
{\it localement quasi-séparé pour la topologie de Zariski}
si $X$, considéré comme espace algébrique sur $\Spec(\mathbf{Z})$ (voir
Spaces, définition \ref{spaces-definition-base-change}),
possède la propriété correspondante.
\end{definition}

\noindent
Un espace algébrique $X$ sur $S$ qui est séparé
(au sens absolu ci-dessus) est bien séparé sur $S$ (et il en va de même
pour les autres propriétés absolues de séparation ci-dessus). Ce point sera étudié
en détail dans
Morphisms of Spaces, section \ref{spaces-morphisms-section-separation-axioms}.
Nous verrons au
lemme \ref{lemma-quasi-separated-quasi-compact-pieces}
que la propriété d'être localement quasi-séparé pour la topologie de Zariski est indépendante du schéma de base
(et donc équivalente à la notion absolue).

\begin{lemma}
\label{lemma-trivial-implications}
Soit $S$ un schéma.
Soit $X$ un espace algébrique sur $S$.
On a les implications suivantes entre les axiomes de séparation
de la définition \ref{definition-separated} :
\begin{enumerate}
\item la propriété d'être séparé entraîne toutes les autres,
\item la propriété d'être quasi-séparé entraîne celle d'être localement quasi-séparé pour la topologie de Zariski.
\end{enumerate}
\end{lemma}

\begin{proof}
La démonstration est omise.
\end{proof}

\begin{lemma}
\label{lemma-characterize-quasi-separated}
Soit $S$ un schéma. Soit $X$ un espace algébrique sur $S$.
Les conditions suivantes sont équivalentes :
\begin{enumerate}
\item $X$ est un espace algébrique quasi-séparé,
\item pour $U \to X$, $V \to X$, où $U$ et $V$ sont des schémas quasi-compacts,
le produit fibré $U \times_X V$ est quasi-compact,
\item pour $U \to X$, $V \to X$, où $U$ et $V$ sont affines,
le produit fibré $U \times_X V$ est quasi-compact.
\end{enumerate}
\end{lemma}

\begin{proof}
Le lemme
\ref{spaces-lemma-category-of-spaces-over-smaller-base-scheme}
de Spaces permet de supposer que $S = \Spec(\mathbf{Z})$.
Puisque $U \times_X V = X \times_{X \times X} (U \times V)$
et que $U \times V$ est quasi-compact lorsque $U$ et $V$ le sont,
on voit que (1) entraîne (2). Il est clair que (2) entraîne (3).
Supposons (3). Choisissons un schéma $W$ et un morphisme étale surjectif
$W \to X$. Alors $W \times W \to X \times X$ est étale surjectif.
Il suffit donc de montrer que
$$
j : W \times_X W = X \times_{(X \times X)} (W \times W) \to W \times W
$$
est quasi-compact, voir Spaces, lemme
\ref{spaces-lemma-descent-representable-transformations-property}.
Si $U \subset W$ et $V \subset W$ sont des ouverts affines, alors
$j^{-1}(U \times V) = U \times_X V$ est quasi-compact par hypothèse.
Puisque les ouverts affines $U \times V$ forment un recouvrement ouvert affine de
$W \times W$ (Schemes, lemme \ref{schemes-lemma-affine-covering-fibre-product}),
on conclut par
Schemes, lemme \ref{schemes-lemma-quasi-compact-affine}.
\end{proof}

\begin{lemma}
\label{lemma-characterize-separated}
Soit $S$ un schéma. Soit $X$ un espace algébrique sur $S$.
Les conditions suivantes sont équivalentes :
\begin{enumerate}
\item $X$ est un espace algébrique séparé,
\item pour $U \to X$, $V \to X$, où $U$ et $V$ sont affines,
le produit fibré $U \times_X V$ est affine et
$$
\mathcal{O}(U) \otimes_\mathbf{Z} \mathcal{O}(V)
\longrightarrow
\mathcal{O}(U \times_X V)
$$
est surjectif.
\end{enumerate}
\end{lemma}

\begin{proof}
Le lemme
\ref{spaces-lemma-category-of-spaces-over-smaller-base-scheme}
de Spaces permet de supposer que $S = \Spec(\mathbf{Z})$.
Puisque $U \times_X V = X \times_{X \times X} (U \times V)$
et que $U \times V$ est affine lorsque $U$ et $V$ le sont,
on voit que (1) entraîne (2).
Supposons (2). Choisissons un schéma $W$ et un morphisme étale surjectif
$W \to X$. Alors $W \times W \to X \times X$ est étale surjectif.
Il suffit donc de montrer que
$$
j : W \times_X W = X \times_{(X \times X)} (W \times W) \to W \times W
$$
est une immersion fermée, voir Spaces, lemme
\ref{spaces-lemma-descent-representable-transformations-property}.
Si $U \subset W$ et $V \subset W$ sont des ouverts affines, alors
$j^{-1}(U \times V) = U \times_X V$ est affine par hypothèse
et le morphisme $U \times_X V \to U \times V$ est une immersion fermée,
car l'homomorphisme d'anneaux correspondant est surjectif.
Puisque les ouverts affines $U \times V$ forment un recouvrement ouvert affine
de $W \times W$
(Schemes, lemme \ref{schemes-lemma-affine-covering-fibre-product}),
on conclut par
Morphisms, lemme \ref{morphisms-lemma-closed-immersion}.
\end{proof}







\section{Points des espaces algébriques}
\label{section-points}

\noindent
Comme il ressort clairement de
Spaces, exemple \ref{spaces-example-affine-line-translation},
il ne convient pas de définir un point d'un espace algébrique comme un monomorphisme
ayant pour source le spectre d'un corps.
Nous définissons plutôt les points comme des classes d'équivalence de morphismes de spectres
de corps, exactement comme expliqué dans
Schemes, section \ref{schemes-section-points}.

\medskip\noindent
Soit $S$ un schéma.
Soit $F$ un préfaisceau sur $(\Sch/S)_{fppf}$.
Soit $K$ un corps. Considérons un morphisme
$$
\Spec(K) \longrightarrow F.
$$
D'après le lemme de Yoneda, celui-ci est donné par un
élément $p \in F(\Spec(K))$. On dit que deux tels
couples $(\Spec(K), p)$ et $(\Spec(L), q)$
sont {\it équivalents} s'il existe
un troisième corps $\Omega$ et un diagramme commutatif
$$
\xymatrix{
\Spec(\Omega) \ar[r] \ar[d] &
\Spec(L) \ar[d]^q \\
\Spec(K) \ar[r]^p &
F.
}
$$
Autrement dit, il existe des extensions de corps
$K \to \Omega$ et $L \to \Omega$ telles que
les images de $p$ et de $q$ coïncident dans
$F(\Spec(\Omega))$. Nous omettons la vérification du fait que ceci
définit une relation d'équivalence.

\begin{definition}
\label{definition-points}
Soit $S$ un schéma. Soit $X$ un espace algébrique sur $S$.
Un {\it point} de $X$ est une classe d'équivalence de morphismes
de spectres de corps vers $X$.
L'ensemble des points de $X$ est noté $|X|$.
\end{definition}

\noindent
Notons que, si $f : X \to Y$ est un morphisme d'espaces algébriques
sur $S$, alors on dispose d'une application induite $|f| : |X| \to |Y|$ qui
envoie un représentant $x : \Spec(K) \to X$ sur le représentant
$f \circ x : \Spec(K) \to Y$.

\begin{lemma}
\label{lemma-scheme-points}
Soit $S$ un schéma. Soit $X$ un schéma sur $S$.
Les points de $X$ en tant que schéma sont canoniquement en bijection
avec les points de $X$ en tant qu'espace algébrique.
\end{lemma}

\begin{proof}
C'est Schemes, lemme \ref{schemes-lemma-characterize-points}.
\end{proof}

\begin{lemma}
\label{lemma-points-cartesian}
Soit $S$ un schéma. Soit
$$
\xymatrix{
Z \times_Y X \ar[r] \ar[d] & X \ar[d] \\
Z \ar[r] & Y
}
$$
un diagramme cartésien d'espaces algébriques sur $S$. Alors l'application entre les ensembles
de points
$$
|Z \times_Y X|
\longrightarrow
|Z| \times_{|Y|} |X|
$$
est surjective.
\end{lemma}

\begin{proof}
En effet, soient des corps $K$, $L$ et des morphismes
$\Spec(K) \to X$, $\Spec(L) \to Z$. L'hypothèse
que leurs images dans $|Y|$ coïncident signifie qu'il
existe une extension commune $M/K$ et $M/L$
telle que les morphismes
$\Spec(M) \to \Spec(K) \to X \to Y$ et
$\Spec(M) \to \Spec(L) \to Z \to Y$ coïncident.
C'est précisément la condition qui assure l'existence d'un
morphisme $\Spec(M) \to Z \times_Y X$.
\end{proof}

\begin{lemma}
\label{lemma-characterize-surjective}
Soit $S$ un schéma.
Soit $X$ un espace algébrique sur $S$.
Soit $f : T \to X$ un morphisme d'un schéma vers $X$.
Les assertions suivantes sont équivalentes :
\begin{enumerate}
\item $f : T \to X$ est surjectif (au sens de
Spaces, définition \ref{spaces-definition-relative-representable-property}),
et
\item $|f| : |T| \to |X|$ est surjectif.
\end{enumerate}
\end{lemma}

\begin{proof}
Supposons (1). Soit $x : \Spec(K) \to X$ un morphisme
du spectre d'un corps vers $X$. Par hypothèse, le morphisme de
schémas $\Spec(K) \times_X T \to \Spec(K)$ est surjectif.
Il existe donc une extension de corps $K'/K$ et un morphisme
$\Spec(K') \to \Spec(K) \times_X T$ tels que le carré
de gauche du diagramme
$$
\xymatrix{
\Spec(K') \ar[r] \ar[d] &
\Spec(K) \times_X T \ar[d] \ar[r] &
T \ar[d] \\
\Spec(K) \ar@{=}[r] &
\Spec(K) \ar[r]^-x & X
}
$$
soit commutatif. Cela montre que $|f| : |T| \to |X|$ est surjectif.

\medskip\noindent
Supposons (2). Soit $Z \to X$ un morphisme, où $Z$ est
un schéma. Il faut montrer que le morphisme de schémas $Z \times_X T \to Z$
est surjectif, c'est-à-dire que $|Z \times_X T| \to |Z|$ est surjectif.
Cela résulte de (2) et du
lemme \ref{lemma-points-cartesian}.
\end{proof}

\begin{lemma}
\label{lemma-points-presentation}
Soit $S$ un schéma.
Soit $X$ un espace algébrique sur $S$.
Soit $X = U/R$ une présentation de $X$, voir
Spaces, définition \ref{spaces-definition-presentation}.
Alors l'image de $|R| \to |U| \times |U|$ est une relation d'équivalence
et $|X|$ est le quotient de $|U|$ par cette relation d'équivalence.
\end{lemma}

\begin{proof}
L'hypothèse signifie que $U$ est un schéma, que $p : U \to X$ est un morphisme
étale surjectif, que $R = U \times_X U$ est un schéma et définit une relation
d'équivalence étale sur $U$ telle que $X = U/R$ comme faisceaux. D'après le
lemme \ref{lemma-characterize-surjective},
on voit que $|U| \to |X|$ est surjective. D'après le
lemme \ref{lemma-points-cartesian},
l'application
$$
|R| \longrightarrow |U| \times_{|X|} |U|
$$
est surjective. L'image de $|R| \to |U| \times |U|$ est donc
exactement l'ensemble des couples $(u_1, u_2) \in |U| \times |U|$
tels que $u_1$ et $u_2$ aient la même image dans $|X|$.
En réunissant ces deux assertions, on obtient le résultat du lemme.
\end{proof}

\begin{lemma}
\label{lemma-topology-points}
Soit $S$ un schéma. Il existe une unique topologie sur les ensembles de points
des espaces algébriques sur $S$ qui possède les propriétés suivantes :
\begin{enumerate}
\item si $X$ est un schéma sur $S$, la topologie sur $|X|$ est la topologie usuelle
(via l'identification du lemme \ref{lemma-scheme-points}),
\item pour tout morphisme d'espaces algébriques $X \to Y$ sur $S$,
l'application $|X| \to |Y|$ est continue, et
\item pour tout morphisme étale $U \to X$, où $U$ est un schéma,
l'application d'espaces topologiques $|U| \to |X|$ est continue et ouverte.
\end{enumerate}
\end{lemma}

\begin{proof}
Soit $X$ un espace algébrique sur $S$. Soit $p : U \to X$ un
morphisme étale surjectif, où $U$ est un schéma sur $S$.
On déclare que $W \subset |X|$ est ouvert si et seulement si $|p|^{-1}(W)$
est une partie ouverte de $|U|$. Cela définit une topologie sur $|X|$
(c'est la topologie quotient sur $|X|$, voir
Topology, lemme \ref{topology-lemma-quotient}).

\medskip\noindent
Montrons que la topologie est indépendante du choix de
la présentation. Pour cela, il suffit de montrer que, si $U'$ est un schéma
et $U' \to X$ un morphisme étale, alors l'application $|U'| \to |X|$
(la topologie sur $|X|$ étant définie à l'aide de $U \to X$ comme ci-dessus)
est ouverte et continue, ce qui démontrera en outre (3).
Posons $U'' = U \times_X U'$ ; on a alors le diagramme commutatif
$$
\xymatrix{
U'' \ar[r] \ar[d] & U' \ar[d] \\
U \ar[r] & X
}
$$
Puisque $U \to X$ et $U' \to X$ sont étales, on voit que
$U'' \to U$ et $U'' \to U'$ sont tous deux des morphismes étales de schémas.
De plus, $U'' \to U'$ est surjectif. On obtient donc
un diagramme commutatif d'applications d'ensembles
$$
\xymatrix{
|U''| \ar[r] \ar[d] & |U'| \ar[d] \\
|U| \ar[r] & |X|
}
$$
La flèche horizontale inférieure est surjective (voir le
lemme \ref{lemma-characterize-surjective}
ou le
lemme \ref{lemma-points-presentation})
et continue par définition de la topologie sur $|X|$.
La flèche horizontale supérieure est surjective, continue et ouverte d'après
Morphisms, lemme \ref{morphisms-lemma-etale-open}.
La flèche verticale gauche est également continue et ouverte d'après
Morphisms, lemme \ref{morphisms-lemma-etale-open}.
Il en résulte formellement que la flèche verticale droite
est continue et ouverte.

\medskip\noindent
Pour achever la démonstration, prouvons (2).
Soit $a : X \to Y$ un morphisme d'espaces algébriques. D'après
Spaces, lemme \ref{spaces-lemma-lift-morphism-presentations},
on peut trouver un diagramme
$$
\xymatrix{
U \ar[d]_p \ar[r]_\alpha & V \ar[d]^q \\
X \ar[r]^a & Y
}
$$
où $U$ et $V$ sont des schémas, et où $p$ et $q$ sont étales et surjectifs.
On en déduit le diagramme
$$
\xymatrix{
|U| \ar[d]_p \ar[r]_\alpha & |V| \ar[d]^q \\
|X| \ar[r]^a & |Y|
}
$$
où toutes les flèches sauf l'horizontale inférieure sont continues et
où les deux flèches verticales sont surjectives et ouvertes. Il s'ensuit que la
flèche horizontale inférieure est continue, comme voulu.
\end{proof}

\begin{definition}
\label{definition-topological-space}
Soit $S$ un schéma. Soit $X$ un espace algébrique sur $S$.
L'{\it espace topologique sous-jacent} à $X$ est l'ensemble des points
$|X|$ muni de la topologie construite au
lemme \ref{lemma-topology-points}.
\end{definition}

\noindent
Il se trouve que cet espace topologique porte la même information
que le petit site de Zariski $X_{Zar}$ de
Spaces, définition \ref{spaces-definition-small-Zariski-site}.

\begin{lemma}
\label{lemma-open-subspaces}
Soit $S$ un schéma.
Soit $X$ un espace algébrique sur $S$.
\begin{enumerate}
\item La règle $X' \mapsto |X'|$ définit une bijection respectant les inclusions
entre les sous-espaces ouverts $X'$ (voir
Spaces, définition \ref{spaces-definition-immersion})
de $X$ et les ouverts de l'espace topologique $|X|$.
\item Une famille $\{X_i \subset X\}_{i \in I}$ de sous-espaces ouverts de $X$
est un recouvrement de Zariski (voir
Spaces, définition \ref{spaces-definition-Zariski-open-covering})
si et seulement si $|X| = \bigcup |X_i|$.
\end{enumerate}
Autrement dit, le petit site de Zariski $X_{Zar}$ de $X$ s'identifie canoniquement
à un site associé à l'espace topologique $|X|$ (voir
Sites, exemple \ref{sites-example-site-topological}).
\end{lemma}

\begin{proof}
Pour démontrer (1), construisons l'inverse de cette règle.
Supposons donc que $W \subset |X|$ soit ouvert. Choisissons une présentation
$X = U/R$ correspondant au morphisme étale surjectif
$p : U \to X$ et aux morphismes étales $s, t : R \to U$.
Par construction, on voit que $|p|^{-1}(W)$ est un
ouvert de $U$. Notons $W' \subset U$ le sous-schéma ouvert correspondant.
Il est clair que $R' = s^{-1}(W') = t^{-1}(W')$ est un ouvert de Zariski
de $R$ qui définit une relation d'équivalence étale sur $W'$.
D'après Spaces, lemme \ref{spaces-lemma-finding-opens}, le morphisme
$X' = W'/R' \to X$ est une immersion ouverte. Ainsi, $X'$ est un espace algébrique
d'après Spaces, lemme \ref{spaces-lemma-representable-over-space}.
Par construction, $|X'| = W$, c'est-à-dire que $X'$ est un sous-espace de $X$
correspondant à $W$. Cela démontre (1).

\medskip\noindent
Pour démontrer (2), notons que, si $\{X_i \subset X\}_{i \in I}$ est une famille
de sous-espaces ouverts, elle est un recouvrement de Zariski si et seulement si
$U = \bigcup U \times_X X_i$ est un recouvrement ouvert. Cela résulte de
la définition d'un recouvrement de Zariski et du fait que le morphisme
$U \to X$ est surjectif comme morphisme de préfaisceaux sur $(\Sch/S)_{fppf}$.
D'autre part, on voit que $|X| = \bigcup |X_i|$ si et seulement si
$U = \bigcup U \times_X X_i$ d'après le lemme \ref{lemma-points-presentation}
(et le fait que les projections $U \times_X X_i \to X_i$ sont surjectives
et étales). L'équivalence de (2) en résulte.
\end{proof}

\begin{lemma}
\label{lemma-factor-through-open-subspace}
Soit $S$ un schéma.
Soient $X$, $Y$ des espaces algébriques sur $S$.
Soit $X' \subset X$ un sous-espace ouvert.
Soit $f : Y \to X$ un morphisme d'espaces algébriques sur $S$.
Alors $f$ se factorise par $X'$ si et seulement si $|f| : |Y| \to |X|$
se factorise par $|X'| \subset |X|$.
\end{lemma}

\begin{proof}
D'après Spaces, lemme \ref{spaces-lemma-base-change-immersions},
on voit que $Y' = Y \times_X X' \to Y$ est une immersion ouverte.
Si $|f|(|Y|) \subset |X'|$, alors $|Y'| = |Y|$ manifestement. Ainsi, $Y' = Y$ d'après le
lemme \ref{lemma-open-subspaces}.
\end{proof}

\begin{lemma}
\label{lemma-etale-image-open}
Soit $S$ un schéma. Soit $X$ un espace algébrique sur $S$.
Soit $U$ un schéma et soit $f : U \to X$ un morphisme étale.
Soit $X' \subset X$ le sous-espace ouvert correspondant à
l'ouvert $|f|(|U|) \subset |X|$ par le
lemme \ref{lemma-open-subspaces}.
Alors $f$ se factorise par un morphisme étale surjectif $f' : U \to X'$.
En outre, si $R = U \times_X U$, alors $R = U \times_{X'} U$ et $X'$ admet
la présentation $X' = U/R$.
\end{lemma}

\begin{proof}
L'existence de la factorisation résulte du
lemme \ref{lemma-factor-through-open-subspace}.
Le morphisme $f'$ est surjectif d'après le
lemme \ref{lemma-characterize-surjective}.
Pour voir que $f'$ est étale, soit $T \to X'$ un morphisme,
où $T$ est un schéma. Alors $T \times_X U = T \times_{X'} U$,
car $X' \to X$ est un monomorphisme de faisceaux. Ainsi, la projection
$T \times_{X'} U \to T$ est étale puisque $f$ est supposé étale.
On a $U \times_X U = U \times_{X'} U$, car $X' \to X$ est un monomorphisme.
Alors $X' = U/R$ résulte de
Spaces, lemme \ref{spaces-lemma-space-presentation}.
\end{proof}

\begin{lemma}
\label{lemma-equivalence-class-point-monomorphism}
Soit $S$ un schéma. Soit $X$ un espace algébrique sur $S$.
Soient $p : \Spec(K) \to X$ et $q : \Spec(L) \to X$
des morphismes, où $K$ et $L$ sont des corps. Supposons que $p$ et $q$
déterminent le même point de $|X|$ et que $p$ soit un monomorphisme.
Alors $q$ se factorise de manière unique par $p$.
\end{lemma}

\begin{proof}
Puisque $p$ et $q$ définissent le même point de $|X|$, on voit que le schéma
$$
Y = \Spec(K) \times_{p, X, q} \Spec(L)
$$
est non vide. Comme tout changement de base d'un monomorphisme est un monomorphisme,
cela signifie que le morphisme de projection $Y \to \Spec(L)$
est un monomorphisme. Ainsi, $Y = \Spec(L)$, voir
Schemes, lemme \ref{schemes-lemma-mono-towards-spec-field}.
On en conclut que $q$ se factorise par $p$. L'unicité
résulte du fait que $p$ est un monomorphisme.
\end{proof}

\begin{lemma}
\label{lemma-points-monomorphism}
Soit $S$ un schéma. Soit $X$ un espace algébrique sur $S$.
Considérons l'application
$$
\{\Spec(k) \to X \text{ monomorphisme où }k\text{ est un corps}\}
\longrightarrow
|X|
$$
Cette application est injective.
\end{lemma}

\begin{proof}
Cela résulte du lemme \ref{lemma-equivalence-class-point-monomorphism}.
\end{proof}

\noindent
Nous verrons dans Decent Spaces,
lemme \ref{decent-spaces-lemma-decent-points-monomorphism},
que l'application du lemme \ref{lemma-points-monomorphism}
est bijective lorsque $X$ est décent.

















\section{Espaces quasi-compacts}
\label{section-quasi-compact}

\begin{definition}
\label{definition-quasi-compact}
Soit $S$ un schéma.
Soit $X$ un espace algébrique sur $S$.
On dit que $X$ est {\it quasi-compact} s'il existe un morphisme
étale surjectif $U \to X$, où $U$ est un schéma quasi-compact.
\end{definition}

\begin{lemma}
\label{lemma-quasi-compact-space}
Soit $S$ un schéma.
Soit $X$ un espace algébrique sur $S$.
Alors $X$ est quasi-compact si et seulement si $|X|$ est quasi-compact.
\end{lemma}

\begin{proof}
Choisissons un schéma $U$ et un morphisme étale surjectif $U \to X$.
Nous utiliserons le lemme \ref{lemma-characterize-surjective}.
Si $U$ est quasi-compact, alors, puisque $|U| \to |X|$ est surjective,
on en conclut que $|X|$ est quasi-compact.
Si $|X|$ est quasi-compact, alors, puisque $|U| \to |X|$ est ouverte,
on voit qu'il existe un ouvert quasi-compact $U' \subset U$
tel que $|U'| \to |X|$ soit surjective (et encore étale).
D'où le résultat.
\end{proof}

\begin{lemma}
\label{lemma-finite-disjoint-quasi-compact}
Une somme disjointe finie d'espaces algébriques quasi-compacts est
un espace algébrique quasi-compact.
\end{lemma}

\begin{proof}
Cela résulte immédiatement du
lemme \ref{lemma-quasi-compact-space}
et du fait topologique correspondant.
\end{proof}

\begin{example}
\label{example-quasi-compact-not-very-reasonable}
L'espace $\mathbf{A}^1_{\mathbf{Q}}/\mathbf{Z}$ est un espace algébrique
quasi-compact.
\end{example}

\begin{lemma}
\label{lemma-space-locally-quasi-compact}
Soit $S$ un schéma.
Soit $X$ un espace algébrique sur $S$.
Tout point de $|X|$ admet un système fondamental de voisinages
ouverts quasi-compacts.
En particulier, $|X|$ est localement quasi-compact au sens de
Topology, définition \ref{topology-definition-locally-quasi-compact}.
\end{lemma}

\begin{proof}
Cela découle formellement du fait qu'il existe un schéma
$U$ et une application surjective, ouverte et continue
$U \to |X|$ entre espaces topologiques. Plus précisément, si
$u \in U$ a pour image $x \in |X|$, alors les images des voisinages
affines de $u$ forment un système fondamental de voisinages ouverts
quasi-compacts de $x$.
\end{proof}







\section{Recouvrements spéciaux}
\label{section-special-coverings}

\noindent
Dans cette section, nous réunissons quelques lemmes immédiats sur l'existence
de recouvrements étales surjectifs d'espaces algébriques.

\begin{lemma}
\label{lemma-cover-by-union-affines}
Soit $S$ un schéma.
Soit $X$ un espace algébrique sur $S$.
Il existe un morphisme étale surjectif $U \to X$, où
$U$ est une somme disjointe de schémas affines.
On peut en outre supposer que chacun de ces schémas affines
est envoyé dans un ouvert affine de $S$.
\end{lemma}

\begin{proof}
Soit $V \to X$ un morphisme étale surjectif.
Soit $V = \bigcup_{i \in I} V_i$ un recouvrement ouvert de Zariski
tel que chaque $V_i$ soit envoyé dans un ouvert affine de $S$.
Posons alors $U = \coprod_{i \in I} V_i$, muni du morphisme induit
$U \to V \to X$. Celui-ci est étale et surjectif, car il est composé
de transformations représentables de foncteurs qui sont étales et
surjectives (d'après le principe général de
Spaces, lemme
\ref{spaces-lemma-composition-representable-transformations-property}
et de
Morphisms, lemmes \ref{morphisms-lemma-composition-surjective} et
\ref{morphisms-lemma-composition-etale}).
\end{proof}

\begin{lemma}
\label{lemma-union-of-quasi-compact}
Soit $S$ un schéma.
Soit $X$ un espace algébrique sur $S$.
Il existe un recouvrement de Zariski $X = \bigcup X_i$
tel que chaque espace algébrique $X_i$ admette un recouvrement
étale surjectif par un schéma affine. On peut en outre supposer
que chaque $X_i$ soit envoyé dans un ouvert affine de $S$.
\end{lemma}

\begin{proof}
D'après le lemme \ref{lemma-cover-by-union-affines}, on peut trouver un morphisme
étale surjectif $U = \coprod U_i \to X$, où chaque $U_i$ est affine et est envoyé
dans un ouvert affine de $S$. Soit $X_i \subset X$ le sous-espace ouvert
de $X$ tel que $U_i \to X$ se factorise par un morphisme étale surjectif
$U_i \to X_i$, voir le
lemme \ref{lemma-etale-image-open}.
Puisque $U = \bigcup U_i$, on voit que $X = \bigcup X_i$.
Comme $U_i \to X_i$ est surjectif, il s'ensuit que $X_i \to S$ se factorise par
un ouvert affine de $S$.
\end{proof}

\begin{lemma}
\label{lemma-quasi-compact-affine-cover}
Soit $S$ un schéma.
Soit $X$ un espace algébrique sur $S$.
Alors $X$ est quasi-compact si et seulement s'il
existe un morphisme étale surjectif $U \to X$
où $U$ est un schéma affine.
\end{lemma}

\begin{proof}
S'il existe un morphisme étale surjectif $U \to X$ où $U$
est affine, alors $X$ est quasi-compact par la définition \ref{definition-quasi-compact}.
Réciproquement, si $X$ est quasi-compact, alors $|X|$ est quasi-compact.
Soit $U = \coprod_{i \in I} U_i$ une somme disjointe de schémas affines
munie d'un morphisme étale surjectif $\varphi : U \to X$
(lemme \ref{lemma-cover-by-union-affines}).
Alors $|X| = \bigcup \varphi(|U_i|)$ et,
par quasi-compacité, il existe des indices $i_1, \ldots, i_n$
tels que $|X| = \bigcup \varphi(|U_{i_j}|)$. Par conséquent,
$U_{i_1} \cup \ldots \cup U_{i_n}$ est un schéma affine muni d'un
morphisme étale surjectif vers $X$.
\end{proof}

\noindent
Le lemme suivant deviendra caduc lorsque nous traiterons des
morphismes séparés dans le chapitre consacré aux morphismes d'espaces algébriques.

\begin{lemma}
\label{lemma-separated-cover}
Soit $S$ un schéma.
Soit $X$ un espace algébrique sur $S$.
Soit $U$ un schéma séparé et soit $U \to X$ étale.
Alors $U \to X$ est séparé, et $R = U \times_X U$ est un schéma séparé.
\end{lemma}

\begin{proof}
Soit $X' \subset X$ le sous-espace ouvert tel que $U \to X$ se factorise
par un morphisme étale surjectif $U \to X'$, voir le
lemme \ref{lemma-etale-image-open}.
Si $U \to X'$ est séparé, alors $U \to X$ l'est aussi, voir
Spaces, lemme
\ref{spaces-lemma-composition-representable-transformations-property}
(car l'immersion ouverte $X' \to X$ est séparée d'après
Spaces, lemme
\ref{spaces-lemma-representable-transformations-property-implication}
et
Schemes, lemme \ref{schemes-lemma-immersions-monomorphisms}).
De plus, puisque $U \times_{X'} U = U \times_X U$, il suffit
de démontrer le résultat après avoir remplacé $X$ par $X'$, c'est-à-dire qu'on peut
supposer $U \to X$ surjectif.
Considérons le diagramme commutatif
$$
\xymatrix{
R = U \times_X U \ar[r] \ar[d] & U \ar[d] \\
U \ar[r] & X
}
$$
Dans la démonstration de
Spaces, lemme \ref{spaces-lemma-properties-diagonal},
nous avons vu que $j : R \to U \times_S U$ est séparé.
Le morphisme de schémas $U \to S$ est séparé, puisque $U$ est un schéma
séparé, voir
Schemes, lemme \ref{schemes-lemma-compose-after-separated}.
Par conséquent, $U \times_S U \to U$ est séparé en tant que changement de base, voir
Schemes, lemme \ref{schemes-lemma-separated-permanence}.
Le schéma $U \times_S U$ est donc séparé (d'après le même lemme).
Puisque $j$ est séparé, on voit de même que $R$ est séparé.
Ainsi, $R \to U$ est un morphisme séparé (d'après
Schemes, lemme \ref{schemes-lemma-compose-after-separated},
là encore). Par conséquent, d'après
Spaces, lemme \ref{spaces-lemma-representable-morphisms-spaces-property}
et le diagramme ci-dessus, on conclut que $U \to X$ est séparé.
\end{proof}

\begin{lemma}
\label{lemma-quasi-separated}
Soit $S$ un schéma. Soit $X$ un espace algébrique sur $S$.
S'il existe un schéma quasi-séparé $U$ et un morphisme
étale surjectif $U \to X$ tel que l'une ou l'autre des projections
$U \times_X U \to U$ soit quasi-compacte, alors $X$ est quasi-séparé.
\end{lemma}

\begin{proof}
On peut considérer $X$ comme un espace algébrique sur $\mathbf{Z}$.
Considérons le diagramme cartésien
$$
\xymatrix{
U \times_X U \ar[r] \ar[d]_j & X \ar[d]^\Delta \\
U \times U \ar[r] & X \times X
}
$$
Puisque $U$ est quasi-séparé, la projection $U \times U \to U$ est
quasi-séparée (elle s'obtient par changement de base à partir d'un morphisme
quasi-séparé de schémas, voir Schemes, lemme \ref{schemes-lemma-separated-permanence}).
L'hypothèse du lemme implique donc que $j$ est quasi-compact d'après
Schemes, lemme \ref{schemes-lemma-quasi-compact-permanence}.
D'après Spaces, lemme \ref{spaces-lemma-representable-morphisms-spaces-property},
on voit que le morphisme $\Delta$ est quasi-compact, comme voulu.
\end{proof}

\begin{lemma}
\label{lemma-quasi-separated-quasi-compact-pieces}
Soit $S$ un schéma.
Soit $X$ un espace algébrique sur $S$.
Les assertions suivantes sont équivalentes :
\begin{enumerate}
\item $X$ est localement quasi-séparé pour la topologie de Zariski sur $S$,
\item $X$ est localement quasi-séparé pour la topologie de Zariski,
\item il existe un recouvrement ouvert de Zariski $X = \bigcup X_i$
tel que, pour tout $i$, il existe un schéma affine
$U_i$ et un morphisme étale, surjectif et quasi-compact
$U_i \to X_i$, et
\item il existe un recouvrement ouvert de Zariski $X = \bigcup X_i$
tel que, pour tout $i$, il existe un schéma affine
$U_i$ qui est envoyé dans un ouvert affine de $S$ et un morphisme
étale, surjectif et quasi-compact $U_i \to X_i$.
\end{enumerate}
\end{lemma}

\begin{proof}
Supposons que les morphismes $U_i \to X_i \subset X$ soient comme en (3). Pour démontrer (4),
choisissons, pour tout $i$, un recouvrement ouvert affine fini $U_i =
U_{i1} \cup \ldots \cup U_{in_i}$ tel que chaque $U_{ij}$ soit envoyé
dans un ouvert affine de $S$. Les composés
$U_{ij} \to U_i \to X_i$ sont étales et quasi-compacts (voir
Spaces, lemme
\ref{spaces-lemma-composition-representable-transformations-property}).
Soit $X_{ij} \subset X_i$ le sous-espace ouvert correspondant à
l'image de $|U_{ij}| \to |X_i|$, voir le
lemme \ref{lemma-etale-image-open}.
Notons que $U_{ij} \to X_{ij}$ est quasi-compact, car $X_{ij} \subset X_i$
est un monomorphisme et que $U_{ij} \to X$ est quasi-compact.
Alors $X = \bigcup X_{ij}$ est un recouvrement comme en (4).
L'implication (4) $\Rightarrow$ (3) est immédiate.

\medskip\noindent
Supposons (4). Pour montrer que $X$ est localement quasi-séparé pour la topologie de Zariski sur $S$,
il suffit de montrer que $X_i$ est quasi-séparé sur $S$.
On peut donc supposer qu'il existe un schéma affine $U$ envoyé dans
un ouvert affine de $S$ et un morphisme étale, surjectif et quasi-compact
$U \to X$. Considérons le carré cartésien
$$
\xymatrix{
U \times_X U \ar[r] \ar[d] & U \times_S U \ar[d] \\
X \ar[r]^-{\Delta_{X/S}} & X \times_S X
}
$$
La flèche verticale droite est étale et surjective (voir
Spaces, lemme
\ref{spaces-lemma-product-representable-transformations-property})
et $U \times_S U$ est affine (car $U$ est envoyé dans un ouvert affine de $S$, voir
Schemes, section \ref{schemes-section-fibre-products}),
et $U \times_X U$ est quasi-compact,
car la projection $U \times_X U \to U$ est quasi-compacte en tant que
changement de base de $U \to X$. Il résulte de
Spaces, lemme \ref{spaces-lemma-representable-morphisms-spaces-property}
que $\Delta_{X/S}$ est quasi-compact, comme voulu.

\medskip\noindent
Supposons (1). Pour démontrer (3), on se ramène immédiatement au cas
où $X$ est quasi-séparé sur $S$. D'après le
lemme \ref{lemma-union-of-quasi-compact},
on peut trouver un recouvrement ouvert de Zariski $X = \bigcup X_i$ tel que chaque
$X_i$ soit envoyé dans un ouvert affine de $S$ et qu'il existe des schémas affines
$U_i$ et des morphismes étales surjectifs $U_i \to X_i$.
Puisque $U_i \to S$ se factorise par un ouvert affine de $S$, on voit que
$U_i \times_S U_i$ est affine, voir
Schemes, section \ref{schemes-section-fibre-products}.
Comme $X$ est quasi-séparé sur $S$, les morphismes
$$
R_i = U_i \times_{X_i} U_i = U_i \times_X U_i
\longrightarrow
U_i \times_S U_i
$$
sont quasi-compacts en tant que changements de base de $\Delta_{X/S}$. On en conclut
que $R_i$ est un schéma quasi-compact. Cela implique à son tour que chaque
projection $R_i \to U_i$ est quasi-compacte. Par conséquent, en appliquant
Spaces, lemme \ref{spaces-lemma-representable-morphisms-spaces-property}
au recouvrement $U_i \to X_i$ et au morphisme $U_i \to X_i$,
on conclut que les morphismes $U_i \to X_i$ sont quasi-compacts, comme voulu.

\medskip\noindent
On voit maintenant que (1), (3) et (4) sont équivalentes. Puisque (3) ne
fait pas intervenir le schéma de base, on en conclut qu'elles sont également équivalentes
à (2).
\end{proof}

\noindent
Le lemme suivant se révélera très utile.

\begin{lemma}
\label{lemma-finite-fibres-presentation}
Soit $S$ un schéma.
Soit $X$ un espace algébrique sur $S$.
Soit $U$ un schéma. Soit $\varphi : U \to X$ un morphisme étale tel que
les projections $R = U \times_X U \to U$ soient quasi-compactes ; c'est par exemple le cas si
$\varphi$ est quasi-compact. Alors les fibres de
$$
|U| \to |X|
\quad\text{et}\quad
|R| \to |X|
$$
sont finies.
\end{lemma}

\begin{proof}
Notons $R = U \times_X U$, et $s, t : R \to U$ les projections.
Soit $u \in U$ un point, et soit $x \in |X|$ son image.
La fibre de $|U| \to |X|$ au-dessus de $x$ est égale à
$s(t^{-1}(\{u\}))$ d'après le
lemme \ref{lemma-points-cartesian},
et la fibre de $|R| \to |X|$ au-dessus de $x$ est $t^{-1}(s(t^{-1}(\{u\})))$.
Puisque $t : R \to U$ est étale et quasi-compact, ses fibres sont finies
(ses fibres sont des sommes disjointes de spectres de corps d'après
Morphisms, lemme \ref{morphisms-lemma-etale-over-field},
et elles sont quasi-compactes). D'où le résultat.
\end{proof}








\section{Propriétés des espaces définies par des propriétés des schémas}
\label{section-types-properties}

\noindent
Toute propriété de schémas qui est locale pour la topologie \'etale donne lieu à une
propriété correspondante des espaces algébriques grâce au lemme suivant.

\begin{lemma}
\label{lemma-type-property}
Soit $S$ un schéma.
Soit $X$ un espace algébrique sur $S$.
Soit $\mathcal{P}$ une propriété de schémas locale pour la topologie \'etale,
voir
Descent, définition \ref{descent-definition-property-local}.
Les assertions suivantes sont équivalentes :
\begin{enumerate}
\item il existe un schéma $U$ et un morphisme \'etale surjectif $U \to X$
tels que le schéma $U$ possède la propriété $\mathcal{P}$, et
\item pour tout schéma $U$ et tout morphisme \'etale $U \to X$,
le schéma $U$ possède la propriété $\mathcal{P}$.
\end{enumerate}
Si $X$ est représentable, ces conditions équivalent à $\mathcal{P}(X)$.
\end{lemma}

\begin{proof}
L'implication (2) $\Rightarrow$ (1) est immédiate.
Réciproquement, choisissons un morphisme \'etale surjectif $U \to X$
où $U$ est un schéma qui possède $\mathcal{P}$, et soit $V$ un schéma \'etale sur
$X$. Alors $U \times_X V \rightarrow V$ est un morphisme \'etale surjectif
de schémas ; ainsi, $V$ hérite de $\mathcal{P}$ de $U \times_X V$, lequel
hérite à son tour de $\mathcal{P}$ de $U$ (voir la discussion qui suit
Descent, définition \ref{descent-definition-property-local}).
La dernière assertion résulte clairement de (1) et de
Descent, définition \ref{descent-definition-property-local}.
\end{proof}

\begin{definition}
\label{definition-type-property}
Soit $\mathcal{P}$ une propriété de schémas qui est
locale pour la topologie \'etale.
Soit $S$ un schéma.
Soit $X$ un espace algébrique sur $S$.
On dit que $X$ {\it possède la propriété $\mathcal{P}$}
si l'une des conditions équivalentes du
lemme \ref{lemma-type-property}
est satisfaite.
\end{definition}

\begin{remark}
\label{remark-list-properties-local-etale-topology}
Voici une liste de propriétés qui sont locales pour la topologie \'etale
(rappelons que les topologies fpqc, fppf, syntomique et lisse sont
plus fines que la topologie \'etale) :
\begin{enumerate}
\item localement noethérien, voir
Descent, lemme \ref{descent-lemma-Noetherian-local-fppf},
\item de Jacobson, voir
Descent, lemme \ref{descent-lemma-Jacobson-local-fppf},
\item localement noethérien et vérifiant $(S_k)$, voir
Descent, lemme \ref{descent-lemma-Sk-local-syntomic},
\item Cohen-Macaulay, voir
Descent, lemme \ref{descent-lemma-CM-local-syntomic},
\item de Gorenstein, voir
Duality for Schemes, lemme \ref{duality-lemma-gorenstein-local-syntomic},
\item réduit, voir
Descent, lemme \ref{descent-lemma-reduced-local-smooth},
\item normal, voir
Descent, lemme \ref{descent-lemma-normal-local-smooth},
\item localement noethérien et vérifiant $(R_k)$, voir
Descent, lemme \ref{descent-lemma-Rk-local-smooth},
\item régulier, voir
Descent, lemme \ref{descent-lemma-regular-local-smooth},
\item de Nagata, voir
Descent, lemme \ref{descent-lemma-Nagata-local-smooth}.
\end{enumerate}
\end{remark}

\noindent
Toute propriété de germes de schémas locale pour la topologie \'etale donne lieu à une
propriété correspondante des espaces algébriques. Voici le lemme de rigueur.

\begin{lemma}
\label{lemma-local-source-target-at-point}
Soit $\mathcal{P}$ une propriété de germes de schémas locale pour la topologie \'etale, voir
Descent, définition \ref{descent-definition-local-at-point}.
Soit $S$ un schéma.
Soit $X$ un espace algébrique sur $S$.
Soit $x \in |X|$ un point de $X$.
Considérons les morphismes \'etales $a : U \to X$, où $U$ est un schéma.
Les assertions suivantes sont équivalentes :
\begin{enumerate}
\item pour tout $U \to X$ comme ci-dessus et tout $u \in U$ tel que $a(u) = x$, on a
$\mathcal{P}(U, u)$, et
\item il existe $U \to X$ comme ci-dessus et $u \in U$ tel que $a(u) = x$ et que l'on ait
$\mathcal{P}(U, u)$.
\end{enumerate}
Si $X$ est représentable, ces conditions équivalent à $\mathcal{P}(X, x)$.
\end{lemma}

\begin{proof}
Démonstration omise.
\end{proof}

\begin{definition}
\label{definition-property-at-point}
Soit $S$ un schéma. Soit $X$ un espace algébrique sur $S$.
Soit $x \in |X|$. Soit $\mathcal{P}$ une propriété de germes de schémas qui est
locale pour la topologie \'etale.
On dit que $X$ {\it possède la propriété $\mathcal{P}$ en $x$} si l'une des
conditions équivalentes du
lemme \ref{lemma-local-source-target-at-point}
est satisfaite.
\end{definition}

\begin{remark}
\label{remark-list-properties-local-ring-local-etale-topology}
Soit $P$ une propriété d'anneaux locaux. Supposons que, pour tout
homomorphisme \'etale d'anneaux $A \to B$ et tout idéal premier $\mathfrak q$ de $B$ au-dessus de
l'idéal premier $\mathfrak p$ de $A$, on ait
$P(A_\mathfrak p) \Leftrightarrow P(B_\mathfrak q)$.
On obtient alors une propriété locale pour la topologie \'etale des germes $(U, u)$ de schémas
en posant $\mathcal{P}(U, u) = P(\mathcal{O}_{U, u})$.
Dans cette situation, nous emploierons l'expression
``l'anneau local de $X$ en $x$ possède $P$'' pour signifier que
$X$ possède la propriété $\mathcal{P}$ en $x$.
Voici une liste de telles propriétés $P$ :
\begin{enumerate}
\item noethérien, voir
More on Algebra, lemme \ref{more-algebra-lemma-Noetherian-etale-extension},
\item de dimension $d$, voir
More on Algebra, lemme \ref{more-algebra-lemma-dimension-etale-extension},
\item régulier, voir
More on Algebra, lemme \ref{more-algebra-lemma-regular-etale-extension},
\item anneau de valuation discrète, cela résulte de (2), (3) et du
lemme \ref{algebra-lemma-characterize-dvr} d'Algebra,
\item réduit, voir
More on Algebra, lemme \ref{more-algebra-lemma-henselization-reduced},
\item normal, voir
More on Algebra, lemme \ref{more-algebra-lemma-henselization-normal},
\item noethérien et de profondeur $k$, voir
More on Algebra, lemme \ref{more-algebra-lemma-henselization-depth},
\item noethérien et Cohen-Macaulay, voir
More on Algebra, lemme \ref{more-algebra-lemma-henselization-CM},
\item noethérien et de Gorenstein, voir
Dualizing Complexes, lemme \ref{dualizing-lemma-flat-under-gorenstein}.
\end{enumerate}
Il existe d'autres propriétés pour lesquelles ceci vaut, par exemple celles d'être un anneau G ou un
anneau de Nagata. Si nous en avons un jour besoin, nous les ajouterons ici,
ainsi que des références à des démonstrations détaillées des faits
d'algèbre correspondants.
\end{remark}







\section{Ensembles constructibles}
\label{section-constructible}

\begin{lemma}
\label{lemma-locally-constructible}
Soit $S$ un schéma. Soit $X$ un espace algébrique sur $S$.
Soit $E \subset |X|$ une partie. Les assertions suivantes sont équivalentes :
\begin{enumerate}
\item pour tout morphisme \'etale $U \to X$, où $U$ est un
schéma, l'image réciproque de $E$ dans $U$ est une partie localement constructible
de $U$,
\item pour tout morphisme \'etale $U \to X$, où $U$ est un
schéma affine, l'image réciproque de $E$ dans $U$ est une partie constructible
de $U$,
\item il existe un morphisme \'etale surjectif $U \to X$, où $U$ est un
schéma, tel que l'image réciproque de $E$ dans $U$ soit une partie localement constructible
de $U$.
\end{enumerate}
\end{lemma}

\begin{proof}
D'après Properties, lemme \ref{properties-lemma-locally-constructible},
on voit que (1) et (2) sont équivalentes. Il est immédiat que (1)
implique (3). Supposons donc donné un morphisme \'etale surjectif
$\varphi : U \to X$, où $U$ est un schéma tel que $\varphi^{-1}(E)$
soit localement constructible. Soit $\varphi' : U' \to X$ un autre
morphisme \'etale, où $U'$ est un schéma. On a alors
$$
E'' = \text{pr}_1^{-1}(\varphi^{-1}(E)) = \text{pr}_2^{-1}((\varphi')^{-1}(E))
$$
où $\text{pr}_1 : U \times_X U' \to U$ et
$\text{pr}_2 : U \times_X U' \to U'$ sont les projections.
D'après Morphisms, lemme \ref{morphisms-lemma-inverse-image-constructible},
on voit que $E''$ est localement constructible dans $U \times_X U'$.
Soit $W' \subset U'$ un ouvert affine. Puisque $\text{pr}_2$ est
\'etale et donc ouverte, on peut choisir un ouvert quasi-compact
$W'' \subset U \times_X U'$ tel que $\text{pr}_2(W'') = W'$.
Alors $\text{pr}_2|_{W''} : W'' \to W'$ est quasi-compact.
On a $W' \cap (\varphi')^{-1}(E) = \text{pr}_2(E'' \cap W'')$
puisque $\varphi$ est surjectif, voir le lemme \ref{lemma-points-cartesian}.
Ainsi, $W' \cap (\varphi')^{-1}(E) = \text{pr}_2(E'' \cap W'')$
est localement constructible d'après
Morphisms, théorème \ref{morphisms-theorem-chevalley}, comme voulu.
\end{proof}

\begin{definition}
\label{definition-locally-constructible}
Soit $S$ un schéma. Soit $X$ un espace algébrique sur $S$.
Soit $E \subset |X|$ une partie. On dit que $E$ est
{\it localement constructible pour la topologie \'etale} si les conditions
équivalentes du lemme \ref{lemma-locally-constructible} sont satisfaites.
\end{definition}

\noindent
Bien entendu, si $X$ est représentable, c'est-à-dire si $X$ est un schéma,
cela signifie simplement que $E$ est une partie localement constructible
de l'espace topologique sous-jacent.







\section{Dimension en un point}
\label{section-dimension}

\noindent
On peut utiliser
Descent, lemme \ref{descent-lemma-dimension-at-point-local},
pour définir la dimension d'un espace
algébrique $X$ en un point $x$. On obtient ainsi une notion différente de la
notion topologique (c'est-à-dire la dimension de $|X|$ en $x$).

\begin{definition}
\label{definition-dimension-at-point}
Soit $S$ un schéma.
Soit $X$ un espace algébrique sur $S$.
Soit $x \in |X|$ un point de $X$.
On définit la {\it dimension de $X$ en $x$} comme
l'élément $\dim_x(X) \in \{0, 1, 2, \ldots, \infty\}$
tel que $\dim_x(X) = \dim_u(U)$ pour tout (ou, de manière équivalente, pour un)
couple $(a : U \to X, u)$ constitué d'un morphisme \'etale $a : U \to X$
d'un schéma vers $X$ et d'un point $u \in U$ tel que $a(u) = x$.
Voir la
définition \ref{definition-property-at-point}, le
lemme \ref{lemma-local-source-target-at-point} et
Descent, lemme \ref{descent-lemma-dimension-at-point-local}.
\end{definition}

\noindent
Attention : en général, on n'a {\bf pas} $\dim_x(X) = \dim_x(|X|)$.
Un contre-exemple est l'espace algébrique $X$ de
Spaces, exemple \ref{spaces-example-infinite-product}.
Plus précisément, soit $x \in |X|$ un point distinct du point générique $x_0$
de $|X|$. Alors $\dim_x(X) = 0$, mais $\dim_x(|X|) = 1$.
En particulier, la dimension de $X$ (telle qu'elle est définie
ci-dessous) est différente de la dimension de $|X|$.

\begin{definition}
\label{definition-dimension}
Soit $S$ un schéma. Soit $X$ un espace algébrique sur $S$.
La {\it dimension} $\dim(X)$ de $X$ est définie par la formule
$$
\dim(X) = \sup\nolimits_{x \in |X|} \dim_x(X)
$$
\end{definition}

\noindent
D'après
Properties, lemme \ref{properties-lemma-dimension},
on voit que c'est la notion usuelle lorsque $X$ est un schéma.
Il existe un autre invariant qui mesure la dimension d'un schéma
en un point, à savoir la dimension de l'anneau local. Cet invariant
est lui aussi compatible avec les morphismes \'etales, voir la
section \ref{section-dimension-local-ring}.






\section{Dimension des anneaux locaux}
\label{section-dimension-local-ring}

\noindent
La dimension de l'anneau local d'un espace algébrique est une notion bien
définie.

\begin{lemma}
\label{lemma-pre-dimension-local-ring}
Soit $S$ un schéma.
Soit $X$ un espace algébrique sur $S$.
Soit $x \in |X|$ un point.
Soit $d \in \{0, 1, 2, \ldots, \infty\}$.
Les assertions suivantes sont équivalentes :
\begin{enumerate}
\item il existe un schéma $U$, un morphisme \'etale $a : U \to X$ et un point
$u \in U$ tel que $a(u) = x$ et $\dim(\mathcal{O}_{U, u}) = d$,
\item pour tout schéma $U$, tout morphisme \'etale $a : U \to X$ et tout point
$u \in U$ tel que $a(u) = x$, on a $\dim(\mathcal{O}_{U, u}) = d$.
\end{enumerate}
Si $X$ est un schéma, ces conditions équivalent à $\dim(\mathcal{O}_{X, x}) = d$.
\end{lemma}

\begin{proof}
Combiner le
lemme \ref{lemma-local-source-target-at-point} et
Descent, lemme \ref{descent-lemma-dimension-local-ring-local}.
\end{proof}

\begin{definition}
\label{definition-dimension-local-ring}
Soit $S$ un schéma. Soit $X$ un espace algébrique sur $S$. Soit $x \in |X|$
un point. La {\it dimension de l'anneau local de $X$ en $x$} est
l'élément $d \in \{0, 1, 2, \ldots, \infty\}$ qui satisfait les conditions
équivalentes du lemme \ref{lemma-pre-dimension-local-ring}. Dans ce cas, on
dira aussi que {\it $x$ est un point de codimension $d$ dans $X$}.
\end{definition}

\noindent
Outre le lemme ci-dessous, signalons aussi au lecteur les
lemmes \ref{lemma-dimension-local-ring} et
\ref{lemma-dimension-decent-invariant-under-etale}.

\begin{lemma}
\label{lemma-dimension}
Soit $S$ un schéma. Soit $X$ un espace algébrique sur $S$.
Les quantités suivantes sont égales :
\begin{enumerate}
\item La dimension de $X$.
\item La borne supérieure des dimensions des anneaux locaux de $X$.
\item La borne supérieure des $\dim_x(X)$ pour $x \in |X|$.
\end{enumerate}
\end{lemma}

\begin{proof}
Les nombres de (1) et (3) sont égaux par la définition \ref{definition-dimension}.
Soit $U \to X$ un morphisme \'etale surjectif issu d'un schéma $U$.
La borne supérieure des $\dim_x(X)$ pour $x \in |X|$ est égale à la
borne supérieure des $\dim_u(U)$ pour les points $u$ de $U$ par définition.
Celle-ci est égale à la borne supérieure des $\dim(\mathcal{O}_{U, u})$ d'après
Properties, lemme \ref{properties-lemma-dimension}. Ce dernier
est à son tour égal à (2) par définition.
\end{proof}




\section{Points génériques}
\label{section-generic-points}

\noindent
Soit $T$ un espace topologique. D'après la deuxième édition de
l'EGA I, un {\it point maximal de $T$} est un point générique d'une composante
irréductible de $T$. Si $T = |X|$ est l'espace topologique associé à
un espace algébrique $X$, il existe au moins deux notions de points maximaux :
on peut considérer les points maximaux de $T$ vu comme espace topologique, ou
les images des points maximaux de $U$, où $U \to X$ est un
morphisme \'etale et $U$ un schéma. La seconde notion correspond à
l'ensemble des points de codimension $0$ (lemme \ref{lemma-codimension-0-points}).
Les points de codimension $0$ sont plus faciles
à manier pour les espaces algébriques généraux ; les deux notions
coïncident pour les espaces algébriques quasi-séparés et, plus généralement, décents
(Decent Spaces, lemme \ref{decent-spaces-lemma-decent-generic-points}).

\begin{lemma}
\label{lemma-codimension-0-points}
Soit $S$ un schéma et soit $X$ un espace algébrique sur $S$.
Soit $x \in |X|$. Considérons les morphismes \'etales $a : U \to X$, où
$U$ est un schéma. Les assertions suivantes sont équivalentes :
\begin{enumerate}
\item $x$ est un point de codimension $0$ dans $X$,
\item il existe $U \to X$ comme ci-dessus et $u \in U$ tel que $a(u) = x$,
et le point $u$ est le point générique d'une composante irréductible de $U$,
\item pour tout $U \to X$ comme ci-dessus et tout $u \in U$ ayant pour image $x$,
le point $u$ est le point générique d'une composante irréductible de $U$.
\end{enumerate}
Si $X$ est représentable, ces conditions équivalent à ce que $x$ soit un point générique
d'une composante irréductible de $|X|$.
\end{lemma}

\begin{proof}
Remarquons qu'un point $u$ d'un schéma $U$ est le point générique d'une
composante irréductible de $U$ si et seulement si $\dim(\mathcal{O}_{U, u}) = 0$
(Properties, lemme \ref{properties-lemma-generic-point}).
Le résultat découle donc de la définition de la codimension d'un
point dans $X$ (définition \ref{definition-dimension-local-ring}).
\end{proof}

\begin{lemma}
\label{lemma-codimension-0-points-dense}
Soit $S$ un schéma et soit $X$ un espace algébrique sur $S$.
L'ensemble des points de codimension $0$ de $X$ est dense dans $|X|$.
\end{lemma}

\begin{proof}
Si $U$ est un schéma, l'ensemble des points génériques de ses composantes
irréductibles est dense dans $U$ (cela vaut pour tout espace topologique quasi-sobre).
Ainsi, si $U \to X$ est un morphisme \'etale surjectif, l'ensemble
des points de codimension $0$ de $X$ est l'image d'une partie dense de
$|U|$ (lemme \ref{lemma-codimension-0-points}). Puisque $|X|$ possède la
topologie quotient induite par $|U| \to |X|$, on conclut.
\end{proof}





\section{Espaces réduits}
\label{section-reduced}

\noindent
Nous avons déjà défini les espaces algébriques réduits à la
section \ref{section-types-properties}.
Nous démontrons ici seulement quelques lemmes simples sur les espaces
algébriques réduits.

\begin{lemma}
\label{lemma-subspace-induced-topology}
Soit $S$ un schéma. Soit $Z \to X$ une immersion d'espaces algébriques.
Alors $|Z| \to |X|$ est un homéomorphisme de $|Z|$ sur une partie localement fermée
de $|X|$.
\end{lemma}

\begin{proof}
Soit $U$ un schéma et soit $U \to X$ un morphisme \'etale surjectif.
Alors $Z \times_X U \to U$ est une immersion de schémas, donc induit un
homéomorphisme de $|Z \times_X U|$ sur une partie localement fermée $T'$
de $|U|$. D'après le lemme \ref{lemma-points-cartesian}, la partie
$T'$ est l'image réciproque de l'image $T$ de $|Z| \to |X|$.
L'application $|Z| \to |X|$ est injective, car la transformation de
foncteurs $Z \to X$ est injective, voir
Spaces, section \ref{spaces-section-Zariski}. D'après
Topology, lemme \ref{topology-lemma-open-morphism-quotient-topology}
on voit que $T$ est localement fermée dans $|X|$. De plus, l'application continue
$|Z| \to T$ est un homéomorphisme, puisque l'application $|Z \times_X U| \to T'$
est un homéomorphisme et que $|Z \times_X U| \to |Z|$ est submersive.
\end{proof}

\noindent
Le lemme suivant nous aidera à construire des sous-espaces
(localement) fermés.

\begin{lemma}
\label{lemma-subspaces-presentation}
Soit $S$ un schéma. Soit $j : R \to U \times_S U$ une relation d'équivalence \'etale.
Soit $X = U/R$ l'espace algébrique associé
(Spaces, théorème \ref{spaces-theorem-presentation}). Il existe une
bijection canonique
$$
\text{Sous-schémas localement fermés stables par }R : Z'\text{ de }U
\leftrightarrow
\text{sous-espaces localement fermés }Z\text{ de }X
$$
De plus, $Z \to X$ est fermé (resp.\ ouvert) si et seulement si
$Z' \to U$ est fermé (resp.\ ouvert).
\end{lemma}

\begin{proof}
Notons $\varphi : U \to X$ l'application canonique. La bijection envoie
$Z \to X$ sur $Z' = Z \times_X U \to U$. Il résulte immédiatement de la définition
que $Z' \to U$ est une immersion, resp.\ une immersion fermée, resp.\ une
immersion ouverte, si tel est le cas de $Z \to X$. Il est aussi clair que $Z'$ est stable par $R$
(voir Groupoids, définition \ref{groupoids-definition-invariant-open}).

\medskip\noindent
Réciproquement, supposons que $Z' \to U$ soit une immersion stable par $R$.
Soit $R'$ la restriction de $R$ à $Z'$, voir
Groupoids, définition \ref{groupoids-definition-restrict-groupoid}.
Puisque $R' = R \times_{s, U} Z' = Z' \times_{U, t} R$ dans ce cas,
on voit que $R'$ est une relation d'équivalence \'etale sur $Z'$. D'après
Spaces, théorème \ref{spaces-theorem-presentation}, on voit que
$Z = Z'/R'$ est un espace algébrique. Par construction, on a
$U \times_X Z = Z'$ ; ainsi, $U \times_X Z \to U$ est une immersion.
Notons que la propriété ``immersion'' est préservée par changement de base
et locale pour la topologie fppf sur la base (voir Spaces, section \ref{spaces-section-lists}).
De plus, les immersions sont séparées et localement quasi-finies (voir
Schemes, lemme \ref{schemes-lemma-immersions-monomorphisms}
et
Morphisms, lemme \ref{morphisms-lemma-immersion-locally-quasi-finite}).
Par conséquent, d'après More on Morphisms, lemme
\ref{more-morphisms-lemma-separated-locally-quasi-finite-morphisms-fppf-descend}
les immersions satisfont à la descente pour les recouvrements fppf. Toutes les hypothèses de
Spaces,
lemme \ref{spaces-lemma-morphism-sheaves-with-P-effective-descent-etale}
sont donc satisfaites pour $Z \to X$, $\mathcal{P}=$``immersion'',
et le morphisme \'etale surjectif $U \to X$. On conclut que $Z \to X$
est représentable et est une immersion, ce qui est la
définition d'un sous-espace (voir
Spaces, définition \ref{spaces-definition-immersion}).

\medskip\noindent
Il est clair que ces constructions sont inverses l'une de l'autre, d'où le résultat.
\end{proof}

\begin{lemma}
\label{lemma-reduced-closed-subspace}
Soit $S$ un schéma.
Soit $X$ un espace algébrique sur $S$.
Soit $T \subset |X|$ une partie fermée.
Il existe un unique sous-espace fermé $Z \subset X$ possédant
les propriétés suivantes : (a) on a $|Z| = T$, et (b) $Z$ est réduit.
\end{lemma}

\begin{proof}
Soit $U \to X$ un morphisme \'etale surjectif, où $U$ est un schéma.
Posons $R = U \times_X U$, de sorte que $X = U/R$, voir
Spaces, lemme \ref{spaces-lemma-space-presentation}.
Comme d'habitude, on note $s, t : R \to U$ les deux morphismes de projection.
D'après le lemme \ref{lemma-points-presentation},
on voit que $T$ correspond à une partie fermée $T' \subset |U|$ telle
que $s^{-1}(T') = t^{-1}(T')$.
Soit $Z' \subset U$ le sous-schéma porté par $T'$ et muni de la structure réduite induite.
Dans ce cas, les produits fibrés
$Z' \times_{U, t} R$ et $Z' \times_{U, s} R$ sont des sous-schémas fermés
de $R$
(Schemes, lemme \ref{schemes-lemma-base-change-immersion})
qui sont \'etales sur $Z'$
(Morphisms, lemme \ref{morphisms-lemma-base-change-etale}),
et sont donc réduits
(car la propriété d'être réduit est locale pour la topologie \'etale, voir la
remarque \ref{remark-list-properties-local-etale-topology}).
Comme ils ont le même espace topologique sous-jacent (voir ci-dessus),
on conclut que $Z' \times_{U, t} R = Z' \times_{U, s} R$.
On peut donc appliquer le lemme \ref{lemma-subspaces-presentation}
pour obtenir un sous-espace fermé $Z \subset X$ dont l'image réciproque sur $U$ est $Z'$.
Par construction, $|Z| = T$ et $Z$ est réduit. Cela prouve l'existence.
Nous omettons la démonstration de l'unicité.
\end{proof}

\begin{lemma}
\label{lemma-map-into-reduction}
Soit $S$ un schéma.
Soient $X$, $Y$ des espaces algébriques sur $S$.
Soit $Z \subset X$ un sous-espace fermé.
Supposons $Y$ réduit.
Un morphisme $f : Y \to X$ se factorise par $Z$ si et seulement si
$f(|Y|) \subset |Z|$.
\end{lemma}

\begin{proof}
Supposons $f(|Y|) \subset |Z|$. Choisissons un diagramme
$$
\xymatrix{
V \ar[d]_b \ar[r]_h & U \ar[d]^a \\
Y \ar[r]^f & X
}
$$
où $U$, $V$ sont des schémas et les flèches verticales sont surjectives et
\'etales. Le schéma $V$ est réduit, voir le
lemme \ref{lemma-type-property}.
Par conséquent, $h$ se factorise par $a^{-1}(Z)$ d'après
Schemes, lemme \ref{schemes-lemma-map-into-reduction}.
Ainsi, $a \circ h$ se factorise par $Z$.
Puisque $Z \subset X$ est un sous-faisceau et que $V \to Y$ est une surjection de faisceaux
sur $(\Sch/S)_{fppf}$, on conclut que $Y \to X$ se factorise
par $Z$.
\end{proof}

\begin{definition}
\label{definition-reduced-induced-space}
Soit $S$ un schéma et soit $X$ un espace algébrique sur $S$.
Soit $Z \subset |X|$ une partie fermée.
Une {\it structure d'espace algébrique sur $Z$} est donnée par un sous-espace fermé
$Z'$ de $X$ tel que $|Z'|$ soit égal à $Z$.
La {\it structure d'espace algébrique réduite induite}
sur $Z$ est celle qui a été construite dans le
lemme \ref{lemma-reduced-closed-subspace}.
La {\it réduction $X_{red}$ de $X$} est la structure d'espace
algébrique réduite induite sur $|X|$.
\end{definition}











\section{Le lieu schématique}
\label{section-schematic}

\noindent
Tout espace algébrique possède un plus grand sous-espace ouvert qui est un
schéma ; cela est plus ou moins clair, mais nous en donnons aussi la preuve ci-dessous.
Bien entendu, ce sous-espace peut être vide, par exemple si
$X = \mathbf{A}^1_{\mathbf{Q}}/\mathbf{Z}$ (le
contre-exemple universel). En revanche, si $X$ est par exemple quasi-séparé,
alors ce plus grand sous-espace ouvert est en fait dense dans $X$ !

\begin{lemma}
\label{lemma-subscheme}
Soit $S$ un schéma.
Soit $X$ un espace algébrique sur $S$.
Il existe un plus grand sous-espace ouvert $X' \subset X$ qui est un schéma.
\end{lemma}

\begin{proof}
Soit $U \to X$ un morphisme \'etale surjectif, où $U$ est un schéma.
Soit $R = U \times_X U$. Les sous-espaces ouverts de $X$ correspondent biunivoquement
avec les sous-schémas ouverts de $U$ qui sont stables par $R$. Il existe donc un
ensemble de tels sous-espaces. Soient $X_i$, $i \in I$, les sous-espaces ouverts
de $X$ qui sont des schémas, c'est-à-dire qui sont représentables. Considérons le
sous-espace ouvert $X' \subset X$ dont l'ensemble de points sous-jacent est
l'ouvert $\bigcup |X_i|$ de $|X|$. D'après le
lemme \ref{lemma-characterize-surjective},
on voit que
$$
\coprod X_i \longrightarrow X'
$$
est une application surjective de faisceaux sur $(\Sch/S)_{fppf}$.
Mais puisque chaque $X_i \to X'$ est représentable par des immersions ouvertes,
on voit qu'en fait l'application est surjective pour la topologie de Zariski.
En effet, si $T \to X'$ est un morphisme d'un schéma
vers $X'$, alors $X_i \times_{X'} T$ est un sous-schéma ouvert de $T$.
On peut donc appliquer
Schemes, lemme \ref{schemes-lemma-glue-functors}
pour voir que $X'$ est un schéma.
\end{proof}

\noindent
Dans la suite de cette section, on dit qu'un sous-espace ouvert
$X'$ d'un espace algébrique $X$ est {\it dense} si la partie ouverte
correspondante $|X'| \subset |X|$ est dense.

\begin{lemma}
\label{lemma-quasi-separated-finite-etale-cover-dense-open-scheme}
Soit $S$ un schéma. Soit $X$ un espace algébrique sur $S$.
S'il existe un morphisme fini, \'etale et surjectif
$U \to X$ où $U$ est un schéma quasi-séparé, alors
il existe un sous-espace ouvert dense $X'$ de $X$ qui est un schéma.
Plus précisément, tout point $x \in |X|$ de codimension $0$ dans $X$
appartient à $X'$.
\end{lemma}

\begin{proof}
Soit $X' \subset X$ le sous-espace ouvert maximal qui est un schéma
(lemme \ref{lemma-subscheme}).
Soit $x \in |X|$ un point de codimension $0$ dans $X$.
D'après le lemme \ref{lemma-codimension-0-points-dense},
il suffit de montrer que $x \in X'$.
Soit $U \to X$ comme dans l'énoncé du lemme.
Écrivons $R = U \times_X U$ et notons $s, t : R \to U$ les projections usuelles.
Notons que $s, t$ sont surjectifs, finis et \'etales.
D'après le lemme \ref{lemma-finite-fibres-presentation},
la fibre de $|U| \to |X|$ au-dessus de $x$ est finie, disons
$\{\eta_1, \ldots, \eta_n\}$. D'après le lemme \ref{lemma-codimension-0-points},
chaque $\eta_i$ est le point générique d'une composante irréductible de $U$.
D'après Properties, lemme \ref{properties-lemma-maximal-points-affine},
on peut trouver un ouvert affine $W \subset U$ contenant
$\{\eta_1, \ldots, \eta_n\}$
(c'est ici que l'on utilise le fait que $U$ est quasi-séparé). D'après
Groupoids, lemme \ref{groupoids-lemma-find-invariant-affine}
on peut supposer que $W$ est stable par $R$.
Puisque $W \subset U$ est un ouvert affine stable par $R$, la restriction
$R_W$ de $R$ à $W$ est égale à $R_W = s^{-1}(W) = t^{-1}(W)$ (voir
Groupoids, définition \ref{groupoids-definition-invariant-open}
et la discussion qui la suit). En particulier, les morphismes $R_W \to W$ sont
eux aussi finis et \'etales. Il en résulte que $R_W$ est affine.
On voit donc que $W/R_W$ est un schéma, d'après
Groupoids, proposition \ref{groupoids-proposition-finite-flat-equivalence}.
D'autre part, $W/R_W$ est un sous-espace ouvert de $X$ d'après
Spaces, lemme \ref{spaces-lemma-finding-opens}, et il contient
$x$ par construction.
\end{proof}

\noindent
Nous étendrons la proposition suivante au cas des
espaces algébriques décents dans
Decent Spaces, théorème
\ref{decent-spaces-theorem-decent-open-dense-scheme}.

\begin{proposition}
\label{proposition-locally-quasi-separated-open-dense-scheme}
Soit $S$ un schéma. Soit $X$ un espace algébrique sur $S$. Si $X$ est
localement quasi-séparé pour la topologie de Zariski (par exemple si $X$ est quasi-séparé), alors
il existe un sous-espace ouvert dense $X'$ de $X$ qui est un schéma.
Plus précisément, tout point $x \in |X|$ de codimension $0$ dans $X$
appartient à $X'$.
\end{proposition}

\begin{proof}
La question est locale sur $X$ d'après le lemme \ref{lemma-subscheme}.
Ainsi, d'après le lemme \ref{lemma-quasi-separated-quasi-compact-pieces},
on peut supposer qu'il existe un schéma affine $U$ et un morphisme
surjectif, quasi-compact et \'etale $U \to X$.
De plus, $U \to X$ est séparé (lemme \ref{lemma-separated-cover}).
Posons $R = U \times_X U$ et notons $s, t : R \to U$ les projections
usuelles. Alors $s, t$ sont surjectifs, quasi-compacts, séparés et
\'etales. Par conséquent, $s, t$ sont aussi quasi-finis et ont des fibres finies
(Morphisms,
lemmes \ref{morphisms-lemma-etale-locally-quasi-finite},
\ref{morphisms-lemma-quasi-finite-locally-quasi-compact}, et
\ref{morphisms-lemma-quasi-finite}).
D'après Morphisms, lemme \ref{morphisms-lemma-generically-finite},
pour tout $\eta \in U$ qui est le point générique d'une
composante irréductible de $U$, il existe un voisinage ouvert
$V \subset U$ de $\eta$ tel que $s^{-1}(V) \to V$ soit fini. D'après
Descent, lemme \ref{descent-lemma-descending-property-finite},
la propriété d'être fini est locale pour la topologie fpqc (et en particulier \'etale) sur le but.
On peut donc appliquer
More on Groupoids, lemme \ref{more-groupoids-lemma-property-invariant},
qui affirme que le plus grand ouvert $W \subset U$ au-dessus duquel $s$ est
fini est stable par $R$. D'après ce qui précède, $W$ contient tout
point générique d'une composante irréductible de $U$.
La restriction $R_W$ de $R$ à $W$ est égale à $R_W = s^{-1}(W) = t^{-1}(W)$
(voir Groupoids, définition \ref{groupoids-definition-invariant-open}
et la discussion qui la suit).
Par construction, $s_W, t_W : R_W \to W$ sont finis \'etales.
Considérons le sous-espace ouvert $X' = W/R_W \subset X$ (voir
Spaces, lemme \ref{spaces-lemma-finding-opens}).
Par construction, l'application d'inclusion $X' \to X$
induit une bijection sur les points de codimension $0$.
On est ainsi ramené au
lemme \ref{lemma-quasi-separated-finite-etale-cover-dense-open-scheme}.
\end{proof}




\section{Obtention d'un schéma}
\label{section-getting-a-scheme}

\noindent
Nous avons utilisé dans la section précédente le fait que le quotient $U/R$ d'un
schéma affine $U$ par une relation d'équivalence $R$ est un schéma si les
morphismes $s, t : R \to U$ sont finis \'etales. C'est un cas particulier
du résultat suivant.

\begin{proposition}
\label{proposition-finite-flat-equivalence-global}
Soit $S$ un schéma.
Soit $(U, R, s, t, c)$ un groupoïde en schémas sur $S$.
Supposons que
\begin{enumerate}
\item $s, t : R \to U$ soient finis localement libres,
\item $j = (t, s)$ soit une relation d'équivalence, et
\item toute partie fermée non vide $Z$ de $U$ contienne un point
$u$ dont la classe d'équivalence pour $R$, $t(s^{-1}(\{u\}))$, soit contenue
dans un ouvert affine de $U$\footnote{Soit $E \subset U$ la
partie des points $u$ tels que $t(s^{-1}(\{u\}))$
soit contenu dans un ouvert affine de $U$. La condition (3) est satisfaite
si $E = U$, ou si tout point de type fini de $U$ appartient à $E$, ou
si tout $u \in U$ se spécialise en un point de $E$.}.
\end{enumerate}
Alors il existe un morphisme fini localement libre $U \to M$
de schémas sur $S$ tel que $R = U \times_M U$ et tel que $M$
représente le faisceau quotient $U/R$ pour la topologie fppf.
\end{proposition}

\begin{proof}
Par l'hypothèse (3) et
Groupoids, lemme \ref{groupoids-lemma-find-invariant-affine}
nous pouvons trouver un recouvrement ouvert $U = \bigcup U_i$ tel que chaque $U_i$
soit stable par $R$ et ouvert affine dans $U$. Posons $R_i = R|_{U_i}$.
Considérons les faisceaux fppf $F = U/R$ et $F_i = U_i/R_i$.
D'après Spaces, lemme \ref{spaces-lemma-finding-opens}, les morphismes
$F_i \to F$ sont représentables et sont des immersions ouvertes.
D'après Groupoids, proposition \ref{groupoids-proposition-finite-flat-equivalence},
les faisceaux $F_i$ sont représentables par des schémas affines.
Si $T$ est un schéma et $T \to F$ un morphisme, alors $V_i = F_i \times_F T$
est ouvert dans $T$, et nous affirmons que $T = \bigcup V_i$. En effet,
localement pour la topologie fppf sur $T$, nous pouvons relever $T \to F$ en un morphisme
$f : T \to U$ et, dans ce cas, $f^{-1}(U_i) \subset V_i$.
Nous en concluons que $F$ est représentable par un schéma, voir
Schemes, lemme \ref{schemes-lemma-glue-functors}.
\end{proof}

\noindent
Par exemple, si $U$ est isomorphe à un sous-schéma localement fermé d'un
schéma affine ou à un sous-schéma localement fermé de
$\text{Proj}(A)$ pour un anneau gradué $A$, alors la troisième hypothèse est satisfaite d'après
Properties, lemme \ref{properties-lemma-ample-finite-set-in-affine}.
En particulier, nous pouvons appliquer ce résultat aux actions libres de groupes finis et de
schémas en groupes finis sur des schémas quasi-affines ou quasi-projectifs.
Par exemple, le quotient $X/G$ d'une variété quasi-projective
$X$ par l'action libre d'un groupe fini $G$ est un schéma. Voici un
énoncé détaillé.

\begin{lemma}
\label{lemma-quotient-scheme}
Soit $S$ un schéma. Soit $G \to S$ un schéma en groupes. Soit $X \to S$ un
morphisme de schémas. Soit $a : G \times_S X \to X$ une action. Supposons que
\begin{enumerate}
\item $G \to S$ soit fini localement libre,
\item l'action $a$ soit libre,
\item $X \to S$ soit affine, ou quasi-affine, ou projectif, ou
quasi-projectif, ou que $X$ soit isomorphe à un sous-schéma ouvert d'un
schéma affine, ou que $X$ soit isomorphe à un sous-schéma ouvert de $\text{Proj}(A)$
pour un anneau gradué $A$, ou que $G \to S$ soit radiciel.
\end{enumerate}
Alors le faisceau quotient $X/G$ pour la topologie fppf est un schéma et $X \to X/G$
est un torseur sous $G$ pour la topologie fppf.
\end{lemma}

\begin{proof}
Montrons d'abord que $X/G$ est un schéma. Puisque l'action est libre, le morphisme
$j = (a, \text{pr}) : G \times_S X \to X \times_S X$
est un monomorphisme et donc une relation d'équivalence, voir
Groupoids, lemme \ref{groupoids-lemma-free-action}. Les morphismes
$s, t : G \times_S X \to X$
sont finis localement libres, puisque nous avons supposé que $G \to S$ est fini localement
libre. Pour conclure, il suffit maintenant de vérifier la dernière hypothèse de la
proposition \ref{proposition-finite-flat-equivalence-global}.
Puisque l'action de $G$ est définie sur $S$, il suffit de montrer que
toute partie finie d'une fibre de $X \to S$ est contenue dans un
ouvert affine de $X$. Si $X$ est isomorphe à un sous-schéma ouvert d'un
schéma affine ou à un sous-schéma ouvert de $\text{Proj}(A)$
pour un anneau gradué $A$, cela résulte de
Properties, lemme \ref{properties-lemma-ample-finite-set-in-affine}.
Si $X \to S$ est affine, ou quasi-affine, ou projectif, ou
quasi-projectif, nous pouvons remplacer $S$ par un ouvert affine et nous
nous ramenons au cas que nous venons de traiter. Si $G \to S$ est radiciel,
alors les orbites des points de $X$ sous l'action de $G$ sont réduites à un point
et la condition est trivialement satisfaite. Nous omettons quelques détails.

\medskip\noindent
Pour voir que $X \to X/G$ est un torseur sous $G$ pour la topologie fppf
(Groupoids, définition \ref{groupoids-definition-principal-homogeneous-space})
il faut montrer que $G \times_S X \to X \times_{X/G} X$
est un isomorphisme et que $X \to X/G$ admet localement des sections pour la topologie fppf.
La seconde assertion résulte du fait que $X \to X/G$ est surjectif
comme morphisme de faisceaux fppf (par construction). La première résulte de
l'isomorphisme $R = U \times_M U$ dans la conclusion de la
proposition \ref{proposition-finite-flat-equivalence-global}
(notons que $R = G \times_S X$ dans notre cas).
\end{proof}

\begin{lemma}
\label{lemma-quotient-separated}
Notations et hypothèses comme dans la
proposition \ref{proposition-finite-flat-equivalence-global}. Alors
\begin{enumerate}
\item si $U$ est quasi-séparé sur $S$, alors $U/R$ est quasi-séparé
sur $S$,
\item si $U$ est quasi-séparé, alors $U/R$ est quasi-séparé,
\item si $U$ est séparé sur $S$, alors $U/R$ est séparé sur $S$,
\item si $U$ est séparé, alors $U/R$ est séparé, et
\item d'autres énoncés pourront être ajoutés ici.
\end{enumerate}
Des résultats analogues valent dans la situation du Lemme \ref{lemma-quotient-scheme}.
\end{lemma}

\begin{proof}
Puisque $M$ représente le faisceau quotient, nous avons un diagramme cartésien
$$
\xymatrix{
R \ar[r]_-j \ar[d] & U \times_S U \ar[d] \\
M \ar[r] & M \times_S M
}
$$
de schémas. Puisque $U \times_S U \to M \times_S M$ est surjectif, fini et localement
libre, pour montrer que $M \to M \times_S M$ est quasi-compact, resp.\ une
immersion fermée, il suffit de montrer que $j : R \to U \times_S U$ est
quasi-compact, resp.\ une immersion fermée, voir
Descent, lemmes \ref{descent-lemma-descending-property-quasi-compact} et
\ref{descent-lemma-descending-property-closed-immersion}.
Puisque $j : R \to U \times_S U$ est un morphisme sur $U$ et que
$R$ est fini sur $U$, nous voyons que $j$ est quasi-compact dès que
le morphisme de projection $U \times_S U \to U$ est quasi-séparé
(Schemes, lemme \ref{schemes-lemma-quasi-compact-permanence}).
Puisque $j$ est un monomorphisme localement de type fini, nous voyons que
$j$ est une immersion fermée dès qu'il est propre
(\'Etale Morphisms, lemme \ref{etale-lemma-characterize-closed-immersion})
ce qui est le cas dès que le morphisme de projection
$U \times_S U \to U$ est séparé
(Morphisms, lemme \ref{morphisms-lemma-image-proper-scheme-closed}).
Ceci prouve (1) et (3). Pour prouver (2) et (4), nous remplaçons $S$ par
$\Spec(\mathbf{Z})$, voir définition \ref{definition-separated}.
Puisque le Lemme \ref{lemma-quotient-scheme} se démontre en appliquant la
proposition \ref{proposition-finite-flat-equivalence-global},
le dernier énoncé est lui aussi clair.
\end{proof}




\section{Points des espaces quasi-séparés}
\label{section-points-quasi-separated}

\noindent
Les points peuvent se comporter très mal pour les espaces algébriques, au degré de généralité adopté
dans le Stacks Project. Cependant, pour les espaces quasi-séparés, leur comportement
est essentiellement semblable à celui des points des schémas. Nous démontrons quelques résultats
à ce sujet dans cette section ; le chapitre sur les espaces décents en contient bien davantage,
voir par exemple
Decent Spaces, section \ref{decent-spaces-section-points}.

\begin{lemma}
\label{lemma-quasi-separated-sober}
Soit $S$ un schéma. Soit $X$ un espace algébrique qui soit localement
quasi-séparé pour la topologie de Zariski et défini sur $S$. Alors l'espace topologique $|X|$ est sobre (voir
Topology, définition \ref{topology-definition-generic-point}).
\end{lemma}

\begin{proof}
En combinant
Topology, lemme \ref{topology-lemma-sober-local}
et le
lemme \ref{lemma-quasi-separated-quasi-compact-pieces}
on voit que l'on peut supposer qu'il existe un schéma affine $U$
et un morphisme \'etale, surjectif et quasi-compact $U \to X$.
Posons $R = U \times_X U$, avec pour projections $s, t : R \to U$. En appliquant le
lemme \ref{lemma-finite-fibres-presentation}
on voit que les fibres de $s, t$ sont finies. Il s'ensuit que toutes les hypothèses de
Topology, lemme \ref{topology-lemma-quotient-kolmogorov}
sont satisfaites, et l'on conclut que $|X|$ est un espace de Kolmogoroff\footnote{
En fait, nous utilisons aussi ici le
Schemes, lemme \ref{schemes-lemma-scheme-sober} (sobriété des schémas),
Morphisms, les lemmes \ref{morphisms-lemma-etale-flat}
et \ref{morphisms-lemma-generalizations-lift-flat} (les générisations
se relèvent le long des morphismes \'etales), le
lemme \ref{lemma-points-presentation} (les points d'un espace algébrique en
termes d'une présentation), et le
lemme \ref{lemma-topology-points} (l'application quotient est ouverte).}.

\medskip\noindent
Il reste à montrer que toute partie fermée irréductible
$T \subset |X|$ possède un point générique. D'après le
lemme \ref{lemma-reduced-closed-subspace}
il existe un sous-espace fermé $Z \subset X$ tel que $|Z| = T$.
Notons que $U \times_X Z \to Z$ est un morphisme \'etale, surjectif et
quasi-compact d'un schéma affine vers $Z$ ; ainsi, $Z$ est localement
quasi-séparé pour la topologie de Zariski d'après le
lemme \ref{lemma-quasi-separated-quasi-compact-pieces}.
Par la
proposition \ref{proposition-locally-quasi-separated-open-dense-scheme}
nous voyons qu'il existe un sous-espace ouvert dense $Z' \subset Z$
qui est un schéma. Cela signifie que $|Z'| \subset T$ est ouvert et dense.
L'espace topologique $|Z'|$ est donc irréductible, ce qui signifie que
$Z'$ est un schéma irréductible. D'après
Schemes, lemme \ref{schemes-lemma-scheme-sober}
nous concluons que $|Z'|$ est l'adhérence d'un unique point
$\eta \in |Z'| \subset T$ et donc aussi que $T = \overline{\{\eta\}}$, ce qui conclut.
\end{proof}

\begin{lemma}
\label{lemma-quasi-compact-quasi-separated-spectral}
Soit $S$ un schéma. Soit $X$ un espace algébrique quasi-compact et
quasi-séparé sur $S$. L'espace topologique $|X|$ est un espace spectral.
\end{lemma}

\begin{proof}
D'après Topology, définition \ref{topology-definition-spectral-space},
nous devons vérifier que $|X|$ est sobre et quasi-compact, qu'il possède une base
d'ouverts quasi-compacts et que l'intersection de deux
ouverts quasi-compacts quelconques est quasi-compacte. D'après le
Lemme \ref{lemma-quasi-separated-sober}, nous voyons que $|X|$ est sobre.
D'après le Lemme \ref{lemma-quasi-compact-space}, nous voyons que $|X|$ est quasi-compact.
D'après le Lemme \ref{lemma-quasi-compact-affine-cover}, il existe un schéma affine
$U$ et un morphisme \'etale surjectif $f : U \to X$.
Puisque $|f| : |U| \to |X|$ est ouverte et continue, et que $|U|$ possède
une base d'ouverts quasi-compacts, nous concluons que $|X|$ possède une base
d'ouverts quasi-compacts. Enfin, supposons que
$A, B \subset |X|$ soient des ouverts quasi-compacts. Alors $A = |X'|$ et $B = |X''|$
pour des sous-espaces ouverts $X', X'' \subset X$ (Lemme \ref{lemma-open-subspaces}),
et nous pouvons choisir des schémas affines $V$ et $W$ et des morphismes
\'etales surjectifs $V \to X'$ et $W \to X''$
(Lemme \ref{lemma-quasi-compact-affine-cover}).
Alors $A \cap B$ est l'image de
$|V \times_X W| \to |X|$ (Lemme \ref{lemma-points-cartesian}).
Puisque $V \times_X W$ est quasi-compact car $X$ est quasi-séparé
(Lemme \ref{lemma-characterize-quasi-separated})
nous concluons que $A \cap B$ est quasi-compacte, ce qui achève la démonstration.
\end{proof}

\noindent
Le lemme suivant permet de prouver qu'un
espace algébrique est isomorphe au spectre d'un corps.

\begin{lemma}
\label{lemma-point-like-spaces}
Soit $S$ un schéma. Soit $k$ un corps.
Soit $X$ un espace algébrique sur $S$ et supposons qu'il existe
un morphisme \'etale surjectif $\Spec(k) \to X$.
Si $X$ est quasi-séparé, alors $X \cong \Spec(k')$,
où $k/k'$ est une extension finie séparable.
\end{lemma}

\begin{proof}
Posons $R = \Spec(k) \times_X \Spec(k)$ ; nous avons donc un
diagramme de produit fibré
$$
\xymatrix{
R \ar[r]_-s \ar[d]_-t & \Spec(k) \ar[d] \\
\Spec(k) \ar[r] & X
}
$$
D'après
Spaces, lemme \ref{spaces-lemma-space-presentation}
nous savons que $X = \Spec(k)/R$ est le faisceau quotient.
Comme $\Spec(k) \to X$ est \'etale, les morphismes $s$ et $t$ sont
\'etales. Ainsi, $R = \coprod_{i \in I} \Spec(k_i)$ est une somme
disjointe de spectres de corps, et $s$ et $t$
induisent tous deux des extensions finies séparables $s, t : k \subset k_i$, voir
Morphisms, lemme \ref{morphisms-lemma-etale-over-field}.
Comme
$$
R = \Spec(k) \times_X \Spec(k)
= (\Spec(k) \times_S \Spec(k)) \times_{X \times_S X, \Delta} X
$$
et que $\Delta$ est quasi-compact par hypothèse, nous concluons que
$R \to \Spec(k) \times_S \Spec(k)$ est quasi-compact.
Ainsi, $R$ est quasi-compact puisque $\Spec(k) \times_S \Spec(k)$ est
affine. Nous concluons que $I$ est fini. Il s'ensuit
que $s$ et $t$ sont des morphismes finis localement libres. D'après
Groupoids, proposition \ref{groupoids-proposition-finite-flat-equivalence}
nous concluons que $\Spec(k)/R$ est
représenté par $\Spec(k')$, l'inclusion $k' \subset k$ étant finie localement libre,
où
$$
k' = \{x \in k \mid s_i(x) = t_i(x)\text{ pour tout }i \in I\}
$$
Il est facile de voir que $k'$ est un corps.
\end{proof}

\begin{remark}
\label{remark-cannot-decide-yet}
Le Lemme \ref{lemma-point-like-spaces} vaut pour les espaces algébriques décents, voir
Decent Spaces, lemme \ref{decent-spaces-lemma-decent-point-like-spaces}.
En fait, un espace algébrique décent à un seul point est un schéma, voir
Decent Spaces, lemme \ref{decent-spaces-lemma-when-field}.
Cela vaut aussi lorsque $X$ est localement séparé, car un
espace algébrique localement séparé est décent, voir
Decent Spaces, lemme \ref{decent-spaces-lemma-locally-separated-decent}.
\end{remark}







\section{Morphismes \'etales d'espaces algébriques}
\label{section-etale-morphisms}

\noindent
Cette section devrait en réalité figurer dans le chapitre consacré aux morphismes d'espaces
algébriques, mais nous avons besoin de la notion d'espace algébrique \'etale sur un autre
afin de définir le petit site \'etale d'un espace algébrique.
Nous devons donc effectuer quelques travaux préliminaires sur les morphismes \'etales de schémas vers des
espaces algébriques et sur les morphismes \'etales entre espaces algébriques.
Pour davantage de résultats sur les morphismes \'etales d'espaces algébriques, voir
Morphisms of Spaces, section \ref{spaces-morphisms-section-etale}.

\begin{lemma}
\label{lemma-etale-over-space}
Soit $S$ un schéma.
Soit $X$ un espace algébrique sur $S$.
Soient $U$, $U'$ des schémas sur $S$.
\begin{enumerate}
\item Si $U \to U'$ est un morphisme \'etale de schémas, et
si $U' \to X$ est un morphisme \'etale de $U'$ vers $X$, alors la
composée $U \to X$ est un morphisme \'etale de $U$ vers $X$.
\item Si $\varphi : U \to X$ et $\varphi' : U' \to X$ sont des
morphismes \'etales vers $X$, et si $\chi : U \to U'$ est un
morphisme de schémas tel que $\varphi = \varphi' \circ \chi$,
alors $\chi$ est un morphisme \'etale de schémas.
\item Si $\chi : U \to U'$ est un morphisme \'etale surjectif
de schémas et si $\varphi' : U' \to X$ est un morphisme tel que
$\varphi = \varphi' \circ \chi$ soit \'etale, alors $\varphi'$
est \'etale.
\end{enumerate}
\end{lemma}

\begin{proof}
Rappelons que la définition d'un morphisme \'etale d'un schéma vers un
espace algébrique est celle de
Spaces, définition
\ref{spaces-definition-relative-representable-property}
puisque tout morphisme d'un schéma vers un espace algébrique
est représentable.

\medskip\noindent
L'assertion (1) du lemme résulte de ceci, du fait que
les morphismes \'etales sont stables par composition
(Morphisms, lemme
\ref{morphisms-lemma-composition-etale})
et des lemmes
de Spaces
\ref{spaces-lemma-composition-representable-transformations-property} et
\ref{spaces-lemma-morphism-schemes-gives-representable-transformation-property}
(qui sont formels).

\medskip\noindent
Pour prouver l'assertion (2), choisissons un schéma $W$ sur $S$ et un
morphisme \'etale surjectif $W \to X$. Considérons le changement de base
$\chi_W : W \times_X U \to W \times_X U'$ de $\chi$.
Comme $W \times_X U$ et $W \times_X U'$ sont \'etales sur $W$, nous concluons que
$\chi_W$ est \'etale d'après
Morphisms, lemme \ref{morphisms-lemma-etale-permanence}.
D'autre part, dans le diagramme commutatif
$$
\xymatrix{
W \times_X U \ar[r] \ar[d] & W \times_X U' \ar[d] \\
U \ar[r] & U'
}
$$
les deux flèches verticales sont \'etales et surjectives.
Ainsi, d'après
Descent, lemme \ref{descent-lemma-syntomic-smooth-etale-permanence}
nous concluons que $U \to U'$ est \'etale.

\medskip\noindent
Pour prouver l'assertion (3), choisissons un schéma $W$ sur $S$ et un morphisme $W \to X$.
Comme ci-dessus, considérons le diagramme
$$
\xymatrix{
W \times_X U \ar[r] \ar[d] & W \times_X U' \ar[d] \ar[r] & W \ar[d] \\
U \ar[r] & U' \ar[r] & X
}
$$
Nous savons que $W \times_X U \to W \times_X U'$ est \'etale surjectif
(comme changement de base de $U \to U'$)
et que $W \times_X U \to W$ est \'etale. Ainsi, $W \times_X U' \to W$
est \'etale d'après Descent, lemme
\ref{descent-lemma-syntomic-smooth-etale-permanence}. Par définition,
cela signifie que $\varphi'$ est \'etale.
\end{proof}

\begin{definition}
\label{definition-etale}
Soit $S$ un schéma.
Un morphisme $f : X \to Y$ entre espaces algébriques sur $S$ est
dit {\it \'etale} si et seulement si, pour tout morphisme \'etale
$\varphi : U \to X$, où $U$ est un schéma, la composée
$f \circ \varphi$ est elle aussi \'etale.
\end{definition}

\noindent
Si $X$ et $Y$ sont des schémas, cela coïncide avec la notion usuelle de
morphisme \'etale de schémas. En fait, dès que $X \to Y$ est un morphisme
représentable d'espaces algébriques, cela coïncide avec la notion définie dans
Spaces, définition \ref{spaces-definition-relative-representable-property}.
Cela résulte en combinant le lemme \ref{lemma-etale-local} ci-dessous avec
Spaces, lemme \ref{spaces-lemma-representable-morphisms-spaces-property}.

\begin{lemma}
\label{lemma-etale-local}
Soit $S$ un schéma.
Soit $f : X \to Y$ un morphisme d'espaces algébriques sur $S$.
Les assertions suivantes sont équivalentes :
\begin{enumerate}
\item $f$ est \'etale,
\item il existe un morphisme \'etale surjectif $\varphi : U \to X$,
où $U$ est un schéma, tel que la composée $f \circ \varphi$ soit
\'etale (en tant que morphisme d'espaces algébriques),
\item il existe un morphisme \'etale surjectif $\psi : V \to Y$,
où $V$ est un schéma, tel que le changement de base $V \times_Y X \to V$
soit \'etale (en tant que morphisme d'espaces algébriques),
\item il existe un diagramme commutatif
$$
\xymatrix{
U \ar[d] \ar[r] & V \ar[d] \\
X \ar[r] & Y
}
$$
où $U$, $V$ sont des schémas, les flèches verticales sont \'etales, la
flèche verticale gauche est surjective et la flèche horizontale est \'etale.
\end{enumerate}
\end{lemma}

\begin{proof}
Montrons que (4) implique (1). Supposons donné un diagramme comme en (4).
Soit $W \to X$ un morphisme \'etale, où $W$ est un schéma. Alors
$W \times_X U \to U$ est \'etale. Ainsi, $W \times_X U \to V$ est \'etale
comme composée des morphismes \'etales de schémas $W \times_X U \to U$
et $U \to V$. Par conséquent, $W \times_X U \to Y$ est \'etale d'après le
Lemme \ref{lemma-etale-over-space} (1). Comme, de plus,
la projection $W \times_X U \to W$ est surjective et \'etale, nous concluons
du Lemme \ref{lemma-etale-over-space} (3) que $W \to Y$ est \'etale.

\medskip\noindent
Montrons que (1) implique (4). Supposons (1). Choisissons un diagramme commutatif
$$
\xymatrix{
U \ar[d] \ar[r] & V \ar[d] \\
X \ar[r] & Y
}
$$
où $U \to X$ et $V \to Y$ sont surjectifs et \'etales, voir
Spaces, lemme \ref{spaces-lemma-lift-morphism-presentations}.
Par hypothèse, le morphisme $U \to Y$ est \'etale ;
donc $U \to V$ est \'etale d'après le Lemme \ref{lemma-etale-over-space} (2).

\medskip\noindent
Nous omettons la preuve que (2) et (3) sont aussi équivalents à (1).
\end{proof}

\begin{lemma}
\label{lemma-composition-etale}
La composée de deux morphismes \'etales d'espaces algébriques
est \'etale.
\end{lemma}

\begin{proof}
Cela résulte immédiatement de la définition.
\end{proof}

\begin{lemma}
\label{lemma-base-change-etale}
Le changement de base d'un morphisme \'etale d'espaces algébriques
par un morphisme quelconque d'espaces algébriques est \'etale.
\end{lemma}

\begin{proof}
Soit $X \to Y$ un morphisme \'etale d'espaces algébriques sur $S$.
Soit $Z \to Y$ un morphisme d'espaces algébriques.
Choisissons un schéma $U$ et un morphisme \'etale surjectif $U \to X$.
Choisissons un schéma $W$ et un morphisme \'etale surjectif $W \to Z$.
Alors $U \to Y$ est \'etale ; ainsi, dans le diagramme
$$
\xymatrix{
W \times_Y U \ar[d] \ar[r] & W \ar[d] \\
Z \times_Y X \ar[r] & Z
}
$$
la flèche horizontale supérieure est \'etale.
De plus, la flèche verticale gauche est surjective
et \'etale (nous omettons la vérification). Nous concluons donc que la flèche
horizontale inférieure est \'etale d'après le Lemme \ref{lemma-etale-local}.
\end{proof}

\begin{lemma}
\label{lemma-etale-permanence}
Soit $S$ un schéma. Soient $X, Y, Z$ des espaces algébriques.
Soient $g : X \to Z$, $h : Y \to Z$ des morphismes \'etales, et soit
$f : X \to Y$ un morphisme tel que $h \circ f = g$.
Alors $f$ est \'etale.
\end{lemma}

\begin{proof}
Choisissons un diagramme commutatif
$$
\xymatrix{
U \ar[d] \ar[r]_\chi & V \ar[d] \\
X \ar[r] & Y
}
$$
où $U \to X$ et $V \to Y$ sont surjectifs et \'etales, voir
Spaces, lemme \ref{spaces-lemma-lift-morphism-presentations}.
Par hypothèse, les morphismes $\varphi : U \to X \to Z$ et
$\psi : V \to Y \to Z$ sont \'etales. De plus, $\psi \circ \chi = \varphi$
par notre hypothèse sur $f, g, h$.
Ainsi, $U \to V$ est \'etale d'après le Lemme \ref{lemma-etale-over-space},
assertion (2).
\end{proof}

\begin{lemma}
\label{lemma-etale-open}
Soit $S$ un schéma.
Si $X \to Y$ est un morphisme \'etale d'espaces algébriques sur $S$,
alors l'application associée $|X| \to |Y|$ entre espaces topologiques
est ouverte.
\end{lemma}

\begin{proof}
Cela résulte clairement du diagramme du
lemme \ref{lemma-etale-local}
et du
lemme \ref{lemma-topology-points}.
\end{proof}

\noindent
Enfin, voici un joli lemme. Il est faux qu'un espace algébrique
muni d'un morphisme \'etale vers un schéma soit nécessairement un schéma, voir
Spaces, exemple \ref{spaces-example-non-representable-descent}.
Mais cela est vrai si le but est le spectre d'un corps.

\begin{lemma}
\label{lemma-etale-over-field-scheme}
Soit $S$ un schéma. Soit $X \to \Spec(k)$
un morphisme \'etale sur $S$, où $k$ est un corps.
Alors $X$ est un schéma.
\end{lemma}

\begin{proof}
Soit $U$ un schéma affine et soit $U \to X$ un morphisme \'etale. D'après la
définition \ref{definition-etale}
nous voyons que $U \to \Spec(k)$ est un morphisme \'etale.
Ainsi, $U = \coprod_{i = 1, \ldots, n} \Spec(k_i)$
est une somme disjointe finie de spectres d'extensions finies séparables
$k_i$ de $k$, voir
Morphisms, lemme \ref{morphisms-lemma-etale-over-field}.
Le morphisme $R = U \times_X U \to U \times_{\Spec(k)} U$ est un monomorphisme,
et $U \times_{\Spec(k)} U$ est aussi une somme disjointe finie de
spectres d'extensions finies séparables de $k$. Ainsi, d'après
Schemes, lemme \ref{schemes-lemma-mono-towards-spec-field}
nous voyons que $R$ est de même une somme disjointe finie de
spectres d'extensions finies séparables de $k$.
Ainsi, $U$ et $R$ sont affines, et
les deux projections $R \to U$ sont finies localement libres.
Par conséquent, $U/R$ est un schéma d'après
Groupoids, proposition \ref{groupoids-proposition-finite-flat-equivalence}.
D'après
Spaces, lemme \ref{spaces-lemma-finding-opens}
il est aussi un sous-espace ouvert de $X$. D'après le
lemme \ref{lemma-subscheme}
nous concluons que $X$ est un schéma.
\end{proof}













\section{Espaces et recouvrements fpqc}
\label{section-fpqc}

\noindent
Soit $S$ un schéma.
Un espace algébrique sur $S$ est défini comme un faisceau pour la topologie fppf muni de
propriétés supplémentaires. Il n'est donc pas immédiatement clair qu'il vérifie
la condition de faisceau pour la topologie fpqc (voir
Topologies, définition \ref{topologies-definition-sheaf-property-fpqc}).
Dans cette section, nous donnons l'argument de Gabber qui le démontre.
Toutefois, lorsque nous disons que l'espace algébrique $X$ vérifie la
condition de faisceau pour la topologie fpqc, nous ne considérons en réalité que les
recouvrements fpqc $\{f_i : T_i \to T\}_{i \in I}$ tels que $T, T_i$ soient des
objets du grand site $(\Sch/S)_{fppf}$ (conformément à
nos conventions, voir la section \ref{section-conventions}).

\begin{proposition}[Gabber]
\label{proposition-sheaf-fpqc}
Soit $S$ un schéma. Soit $X$ un espace algébrique sur $S$. Alors
$X$ vérifie la condition de faisceau pour la topologie fpqc.
\end{proposition}

\begin{proof}
Puisque $X$ est un faisceau pour la topologie de Zariski, il suffit de montrer
ce qui suit. Étant donné un morphisme plat surjectif de schémas affines
$f : T' \to T$, on a :
$X(T)$ est le noyau du couple d'applications $X(T') \to X(T' \times_T T')$.
Voir Topologies, lemme \ref{topologies-lemma-sheaf-property-fpqc}
(nous omettons ici un petit argument, car le lemme cité
est formulé pour des foncteurs définis sur la catégorie de tous les schémas).

\medskip\noindent
Soient $a, b : T \to X$ deux morphismes tels que $a \circ f = b \circ f$.
Il faut montrer que $a = b$. Considérons le produit fibré
$$
E = X \times_{\Delta_{X/S}, X \times_S X, (a, b)} T.
$$
D'après Spaces, lemme \ref{spaces-lemma-properties-diagonal},
le morphisme $\Delta_{X/S}$ est un monomorphisme représentable. Ainsi,
$E \to T$ est un monomorphisme de schémas. Notre hypothèse
$a \circ f = b \circ f$ implique que $T' \to T$ se factorise (uniquement) par $E$.
Considérons le diagramme commutatif
$$
\xymatrix{
T' \times_T E \ar[r] \ar[d] & E \ar[d] \\
T' \ar[r] \ar@/^5ex/[u] \ar[ru] & T
}
$$
Puisque la projection $T' \times_T E \to T'$ est un monomorphisme
admettant une section, nous concluons que c'est un isomorphisme. Nous en déduisons que
$E \to T$ est un isomorphisme d'après
Descent, lemme \ref{descent-lemma-descending-property-isomorphism}.
Cela signifie que $a = b$, comme voulu.

\medskip\noindent
Soit ensuite $c : T' \to X$ un morphisme tel que les deux composées
$T' \times_T T' \to T' \to X$ soient égales. Il faut trouver un morphisme
$a : T \to X$ dont la composée avec $T' \to T$ soit $c$. Choisissons un
schéma affine $U$ et un morphisme \'etale $U \to X$ tels que l'image de
$|U| \to |X|$ contienne l'image de $|c| : |T'| \to |X|$.
C'est possible d'après les lemmes \ref{lemma-topology-points} et
\ref{lemma-cover-by-union-affines}, le fait qu'une somme disjointe finie de
schémas affines est affine, et le fait que $|T'|$ est quasi-compact
(nous omettons un petit argument). Puisque $U \to X$ est séparé
(lemme \ref{lemma-separated-cover}), on voit que
$$
V = U \times_{X, c} T' \longrightarrow T'
$$
est un morphisme de schémas \'etale, surjectif et séparé
(pour voir qu'il est surjectif, utiliser le lemme \ref{lemma-points-cartesian}
et notre choix de $U \to X$). L'égalité
$c \circ \text{pr}_0 = c \circ \text{pr}_1$ signifie que nous obtenons une
donnée de descente sur $V/T'/T$
(Descent, définition \ref{descent-definition-descent-datum}),
car
\begin{align*}
V \times_{T'} (T' \times_T T')
& =
U \times_{X, c \circ \text{pr}_0} (T' \times_T T') \\
& =
(T' \times_T T') \times_{c \circ \text{pr}_1, X} U \\
& =
(T' \times_T T') \times_{T'} V
\end{align*}
Le morphisme $V \to T'$ est ind-quasi-affine d'après
More on Morphisms, lemme
\ref{more-morphisms-lemma-etale-separated-ind-quasi-affine}
(car les morphismes \'etales sont localement quasi-finis, voir
Morphisms, lemme \ref{morphisms-lemma-etale-locally-quasi-finite}).
D'après More on Groupoids, lemme \ref{more-groupoids-lemma-ind-quasi-affine},
la donnée de descente est effective. Soit $W \to T$ un morphisme
tel qu'il existe un isomorphisme $\alpha : T' \times_T W \to V$
compatible avec la donnée de descente donnée sur $V$ et la donnée de descente
canonique sur $T' \times_T W$. Alors $W \to T$ est surjectif et \'etale
(Descent, lemmes \ref{descent-lemma-descending-property-surjective} et
\ref{descent-lemma-descending-property-etale}).
Considérons la composée
$$
b' : T' \times_T W \longrightarrow V = U \times_{X, c} T' \longrightarrow U
$$
Les deux composées
$b' \circ (\text{pr}_0, 1), 
b' \circ (\text{pr}_1, 1) :
(T' \times_T T') \times_T W \to T' \times_T W \to U$
coïncident par notre choix de $\alpha$ et la propriété correspondante de $c$
(calcul omis). Ainsi, $b'$ descend en un morphisme $b : W \to U$ d'après
Descent, lemme \ref{descent-lemma-fpqc-universal-effective-epimorphisms}.
Le diagramme
$$
\xymatrix{
T' \times_T W \ar[r] \ar[d] & W \ar[r]_b & U \ar[d] \\
T' \ar[rr]^c &  & X
}
$$
est commutatif. Cela signifie que nous avons démontré l'existence
d'une solution $a$ localement pour la topologie \'etale sur $T$, c'est-à-dire d'un $a' : W \to X$.
Cependant, puisque nous avons démontré l'unicité
dans le premier paragraphe, nous constatons que cette solution locale pour la topologie \'etale
satisfait à la condition de recollement, c'est-à-dire que
$\text{pr}_0^*a' = \text{pr}_1^*a'$ comme éléments de $X(W \times_T W)$.
Puisque $X$ est un faisceau pour la topologie \'etale, il existe un unique $a \in X(T)$ dont la restriction
est $a'$ sur $W$.
\end{proof}









\section{Le site \'etale d'un espace algébrique}
\label{section-etale-site}

\noindent
Dans cette section, nous définissons le petit site \'etale d'un espace algébrique.
C'est l'analogue du petit site \'etale $S_\etale$ d'un schéma.
Le lemme \ref{lemma-etale-over-space} implique que, dans la définition ci-dessous,
tout morphisme entre objets du site \'etale de $X$ est \'etale, et que
tout schéma \'etale sur un objet de $X_\etale$ est encore un objet de
$X_\etale$.

\begin{definition}
\label{definition-etale-site}
Soit $S$ un schéma.
Soit $\Sch_{fppf}$ un grand site fppf contenant $S$,
et soit $\Sch_\etale$ le grand site \'etale correspondant
(c'est-à-dire ayant la même catégorie sous-jacente).
Soit $X$ un espace algébrique sur $S$.
Le {\it petit site \'etale $X_\etale$} de $X$ est défini comme suit :
\begin{enumerate}
\item un objet de $X_\etale$ est un morphisme $\varphi : U \to X$
où $U \in \Ob((\Sch/S)_\etale)$ est un schéma et
$\varphi$ est un morphisme \'etale,
\item un morphisme $(\varphi : U \to X) \to (\varphi' : U' \to X)$
est donné par un morphisme de schémas $\chi : U \to U'$ tel que
$\varphi = \varphi' \circ \chi$, et
\item une famille de morphismes $\{(U_i \to X) \to (U \to X)\}_{i \in I}$
de $X_\etale$ est un recouvrement si et seulement si $\{U_i \to U\}_{i \in I}$
est un recouvrement de $(\Sch/S)_\etale$.
\end{enumerate}
\end{definition}

\noindent
Une conséquence de notre choix est que le site \'etale d'un espace algébrique
n'a en général pas d'objet final ! En revanche, si $X$ est
un schéma, la définition ci-dessus coïncide avec
Topologies, définition \ref{topologies-definition-big-small-etale}.

\medskip\noindent
Le site précédent est celui que nous employons par défaut, mais nous utiliserons aussi
deux variantes. On peut considérer tous les {\it espaces algébriques}
$U$ \'etales sur $X$, ce qui donne le site
$X_{spaces, \etale}$ défini ci-dessous, ou tous les {\it schémas affines}
$U$ \'etales sur $X$, ce qui donne le site
$X_{affine, \etale}$ défini ci-dessous. La première de ces deux notions
sert à étudier la fonctorialité du petit site \'etale, voir le
lemme \ref{lemma-functoriality-etale-site}.

\begin{definition}
\label{definition-spaces-etale-site}
Soit $S$ un schéma.
Soit $\Sch_{fppf}$ un grand site fppf contenant $S$,
et soit $\Sch_\etale$ le grand site \'etale correspondant
(c'est-à-dire ayant la même catégorie sous-jacente).
Soit $X$ un espace algébrique sur $S$.
Le site {\it $X_{spaces, \etale}$} de $X$ est défini comme suit :
\begin{enumerate}
\item un objet de $X_{spaces, \etale}$ est un morphisme
$\varphi : U \to X$, où $U$ est un espace algébrique sur $S$ et
$\varphi$ est un morphisme \'etale d'espaces algébriques sur $S$,
\item un morphisme $(\varphi : U \to X) \to (\varphi' : U' \to X)$ de
$X_{spaces, \etale}$ est donné par un morphisme d'espaces algébriques
$\chi : U \to U'$ tel que $\varphi = \varphi' \circ \chi$, et
\item une famille de morphismes
$\{\varphi_i : (U_i \to X) \to (U \to X)\}_{i \in I}$
de $X_{spaces, \etale}$ est un recouvrement si et seulement si
$|U| = \bigcup \varphi_i(|U_i|)$.
\end{enumerate}
Comme d'habitude, nous choisissons un ensemble de recouvrements de ce type, contenant au moins
les recouvrements de $X_\etale$, comme dans
Sets, lemme \ref{sets-lemma-coverings-site},
afin de faire de $X_{spaces, \etale}$ un site.
\end{definition}

\noindent
Puisque le morphisme identité de $X$ est \'etale, il est clair que
$X_{spaces, \etale}$ possède un objet final.
Montrons immédiatement que le topos correspondant est équivalent au
petit topos \'etale de $X$.

\begin{lemma}
\label{lemma-compare-etale-sites}
Le foncteur
$$
X_\etale \longrightarrow X_{spaces, \etale}, \quad
U/X \longmapsto U/X
$$
est un foncteur cocontinu spécial
(Sites, définition \ref{sites-definition-special-cocontinuous-functor})
et induit donc une équivalence de topos
$\Sh(X_\etale) \to \Sh(X_{spaces, \etale})$.
\end{lemma}

\begin{proof}
Il faut montrer que le foncteur satisfait aux hypothèses (1) -- (5) du
Sites, lemme \ref{sites-lemma-equivalence}.
Il est clair que le foncteur est continu et cocontinu, ce qui
prouve les hypothèses (1) et (2).
Les hypothèses (3) et (4) sont satisfaites simplement parce que le foncteur est pleinement fidèle.
L'hypothèse (5) est satisfaite parce que, par définition, un espace algébrique admet
un recouvrement par un schéma.
\end{proof}

\begin{remark}
\label{remark-explain-equivalence}
Expliquons le sens du lemme \ref{lemma-compare-etale-sites}.
Soit $S$ un schéma et soit $X$ un espace algébrique sur $S$.
Soit $\mathcal{F}$ un faisceau sur le petit site \'etale $X_\etale$ de
$X$. Le lemme affirme qu'il existe un unique faisceau $\mathcal{F}'$ sur
$X_{spaces, \etale}$ dont la restriction est $\mathcal{F}$ sur la
sous-catégorie $X_\etale$. Si $U \to X$ est un morphisme \'etale
d'espaces algébriques, comment calculer $\mathcal{F}'(U)$ ? Par définition
d'un espace algébrique, il existe un schéma $U'$ et un morphisme
\'etale surjectif $U' \to U$. Alors $\{U' \to U\}$ est un recouvrement de
$X_{spaces, \etale}$, et l'on obtient donc un diagramme exact
$$
\xymatrix{
\mathcal{F}'(U) \ar[r] &
\mathcal{F}(U') \ar@<1ex>[r] \ar@<-1ex>[r] &
\mathcal{F}(U' \times_U U').
}
$$
Notons que $U' \times_U U'$ est un schéma ; nous pouvons donc
écrire $\mathcal{F}$ et non $\mathcal{F}'$.
Nous voyons ainsi comment calculer $\mathcal{F}'$
à partir du faisceau $\mathcal{F}$.
\end{remark}

\begin{definition}
\label{definition-affine-etale-site}
Soit $S$ un schéma.
Soit $\Sch_{fppf}$ un grand site fppf contenant $S$,
et soit $\Sch_\etale$ le grand site \'etale correspondant
(c'est-à-dire ayant la même catégorie sous-jacente).
Soit $X$ un espace algébrique sur $S$.
Le site {\it $X_{affine, \etale}$} de $X$ est défini comme suit :
\begin{enumerate}
\item un objet de $X_{affine, \etale}$ est un morphisme
$\varphi : U \to X$, où $U \in \Ob((\Sch/S)_\etale)$ est un schéma affine et
$\varphi$ est un morphisme \'etale,
\item un morphisme $(\varphi : U \to X) \to (\varphi' : U' \to X)$ de
$X_{affine, \etale}$ est donné par un morphisme de schémas
$\chi : U \to U'$ tel que $\varphi = \varphi' \circ \chi$, et
\item une famille de morphismes
$\{\varphi_i : (U_i \to X) \to (U \to X)\}_{i \in I}$
de $X_{affine, \etale}$ est un recouvrement si et seulement si
$\{U_i \to U\}$ est un recouvrement \'etale standard, voir
Topologies, définition \ref{topologies-definition-standard-etale}.
\end{enumerate}
Comme d'habitude, nous choisissons un ensemble de recouvrements de ce type, comme dans
Sets, lemme \ref{sets-lemma-coverings-site},
afin de faire de $X_{affine, \etale}$ un site.
\end{definition}

\begin{lemma}
\label{lemma-alternative}
Soit $S$ un schéma. Soit $X$ un espace algébrique sur $S$.
Le foncteur $X_{affine, \etale} \to X_\etale$
est cocontinu spécial et induit une équivalence de topos de
$\Sh(X_{affine, \etale})$ vers $\Sh(X_\etale)$.
\end{lemma}

\begin{proof}
Omis. Indication : comparer avec la preuve de
Topologies, lemme \ref{topologies-lemma-affine-big-site-etale}.
\end{proof}

\begin{definition}
\label{definition-etale-topos}
Soit $S$ un schéma. Soit $X$ un espace algébrique sur $S$.
Le {\it topos \'etale} de $X$, ou plus précisément le
{\it petit topos \'etale} de $X$, est la catégorie
$\Sh(X_\etale)$
des faisceaux d'ensembles sur $X_\etale$.
\end{definition}

\noindent
D'après le
lemme \ref{lemma-compare-etale-sites},
on a
$\Sh(X_\etale) = \Sh(X_{spaces, \etale})$,
de sorte que l'on peut aussi voir ce topos comme la catégorie des faisceaux d'ensembles sur
$X_{spaces, \etale}$. De même, d'après le
lemme \ref{lemma-alternative},
on voit que
$\Sh(X_\etale) = \Sh(X_{affine, \etale})$.
Ce topos est fonctoriel par rapport aux morphismes
d'espaces algébriques. Voici un énoncé précis.

\begin{lemma}
\label{lemma-functoriality-etale-site}
Soit $S$ un schéma.
Soit $f : X \to Y$ un morphisme d'espaces algébriques sur $S$.
\begin{enumerate}
\item Le foncteur continu
$$
Y_{spaces, \etale} \longrightarrow X_{spaces, \etale}, \quad
V \longmapsto X \times_Y V
$$
induit un morphisme de sites
$$
f_{spaces, \etale} :
X_{spaces, \etale}
\to
Y_{spaces, \etale}.
$$
\item La règle $f \mapsto f_{spaces, \etale}$ est compatible aux
compositions, autrement dit $(f \circ g)_{spaces, \etale}
= f_{spaces, \etale} \circ g_{spaces, \etale}$ (voir
Sites, définition \ref{sites-definition-composition-morphisms-sites}).
\item Le morphisme de topos associé à $f_{spaces, \etale}$
induit, via le lemme \ref{lemma-compare-etale-sites}, un morphisme de topos
$f_{small} : \Sh(X_\etale) \to \Sh(Y_\etale)$
dont la construction est compatible aux compositions.
\item Si $f$ est un morphisme représentable d'espaces algébriques,
alors $f_{small}$ provient d'un morphisme de sites
$X_\etale \to Y_\etale$,
correspondant au foncteur continu $V \mapsto X \times_Y V$.
\end{enumerate}
\end{lemma}

\begin{proof}
Montrons que le foncteur décrit en (1) satisfait aux hypothèses
de Sites, proposition \ref{sites-proposition-get-morphism}.
Il faut donc montrer que
$Y_{spaces, \etale}$ possède un objet final (à savoir $Y$) et que
le foncteur transforme celui-ci en un objet final de $X_{spaces, \etale}$
(à savoir $X$). C'est clair, car $X \times_Y Y = X$ dans toute catégorie.
Il faut ensuite montrer que $Y_{spaces, \etale}$ possède des produits fibrés.
C'est vrai puisque la catégorie des espaces algébriques possède des produits fibrés,
et puisque $V \times_Y V'$ est \'etale sur $Y$ si $V$ et $V'$ sont \'etales
sur $Y$ (voir les lemmes \ref{lemma-composition-etale} et
\ref{lemma-base-change-etale} ci-dessus).
La proposition s'applique donc et fournit un morphisme
de sites comme en (1).

\medskip\noindent
L'assertion (2) s'obtient en explicitant les définitions.
L'assertion (3) est claire en utilisant les équivalences pour $X$ et $Y$
du lemme \ref{lemma-compare-etale-sites} ci-dessus.
L'assertion (4) en résulte car, si $f$ est représentable, les
foncteurs ci-dessus s'insèrent dans un diagramme commutatif
$$
\xymatrix{
X_\etale \ar[r] &
X_{spaces, \etale} \\
Y_\etale \ar[r] \ar[u] &
Y_{spaces, \etale} \ar[u]
}
$$
de catégories.
\end{proof}

\noindent
On peut préciser quelque peu le lemme précédent lorsqu'on décrit
la relation entre les faisceaux sur $X$ et ceux sur $Y$.
Plus précisément, on peut la formuler en termes de $f$-applications ; comparer
Sheaves, définition \ref{sheaves-definition-f-map}, comme suit.

\begin{definition}
\label{definition-f-map}
Soit $S$ un schéma.
Soit $f : X \to Y$ un morphisme d'espaces algébriques sur $S$.
Soit $\mathcal{F}$ un faisceau d'ensembles sur $X_\etale$ et
soit $\mathcal{G}$ un faisceau d'ensembles sur $Y_\etale$.
Une {\it $f$-application $\varphi : \mathcal{G} \to \mathcal{F}$}
est une collection d'applications
$\varphi_{(U, V, g)} : \mathcal{G}(V) \to \mathcal{F}(U)$
indexée par les diagrammes commutatifs
$$
\xymatrix{
U \ar[d]_g \ar[r] & X \ar[d]^f \\
V \ar[r] & Y
}
$$
où $U \in X_\etale$, $V \in Y_\etale$, telle que, pour tout
diagramme prolongé
$$
\xymatrix{
U' \ar[r] \ar[d]_{g'} & U \ar[d]_g \ar[r] & X \ar[d]^f \\
V' \ar[r] & V \ar[r] & Y
}
$$
où $V' \to V$ et $U' \to U$ sont des morphismes \'etales de schémas, le diagramme
$$
\xymatrix{
\mathcal{G}(V)
\ar[rr]_{\varphi_{(U, V, g)}}
\ar[d]_{\text{restriction de }\mathcal{G}} & &
\mathcal{F}(U)
\ar[d]^{\text{restriction de }\mathcal{F}} \\
\mathcal{G}(V')
\ar[rr]^{\varphi_{(U', V', g')}} & &
\mathcal{F}(U')
}
$$
est commutatif.
\end{definition}

\begin{lemma}
\label{lemma-f-map}
Soit $S$ un schéma.
Soit $f : X \to Y$ un morphisme d'espaces algébriques sur $S$.
Soit $\mathcal{F}$ un faisceau d'ensembles sur $X_\etale$ et
soit $\mathcal{G}$ un faisceau d'ensembles sur $Y_\etale$.
Il existe des bijections canoniques entre les trois ensembles suivants :
\begin{enumerate}
\item l'ensemble des morphismes $\mathcal{G} \to f_{small, *}\mathcal{F}$ ;
\item l'ensemble des morphismes $f_{small}^{-1}\mathcal{G} \to \mathcal{F}$ ;
\item l'ensemble des $f$-applications $\varphi : \mathcal{G} \to \mathcal{F}$.
\end{enumerate}
\end{lemma}

\begin{proof}
Notons que (1) et (2) coïncident parce que les foncteurs $f_{small, *}$
et $f_{small}^{-1}$ forment un couple de foncteurs adjoints.
Supposons que $\alpha : f_{small}^{-1}\mathcal{G} \to \mathcal{F}$
soit un morphisme de faisceaux sur $X_\etale$. Considérons un diagramme
$$
\xymatrix{
U \ar[d]_g \ar[r]_{j_U} & X \ar[d]^f \\
V \ar[r]^{j_V} & Y
}
$$
comme dans la définition \ref{definition-f-map}.
La commutativité du diagramme fournit également un morphisme
$g_{small}^{-1}(j_V)^{-1}\mathcal{G} \to (j_U)^{-1}\mathcal{F}$
(comparer Sites, section \ref{sites-section-localize}, pour la
description des foncteurs de localisation). Nous obtenons donc en particulier
un morphisme
$\varphi_{(U, V, g)} :
\mathcal{G}(V) = (j_V)^{-1}\mathcal{G}(V)
\to
(j_U)^{-1}\mathcal{F}(U) = \mathcal{F}(U)$.
Nous omettons de vérifier que cette règle est compatible aux
restrictions ultérieures et définit une $f$-application de $\mathcal{G}$ vers
$\mathcal{F}$.

\medskip\noindent
Réciproquement, supposons donnée une $f$-application
$\varphi = (\varphi_{(U, V, g)})$.
Désignons par $\mathcal{G}'$ (resp.\ $\mathcal{F}'$) le prolongement de
$\mathcal{G}$ (resp.\ $\mathcal{F}$) à $Y_{spaces, \etale}$
(resp.\ $X_{spaces, \etale}$), voir le
lemme \ref{lemma-compare-etale-sites}.
Il faut alors construire un morphisme de faisceaux
$$
\mathcal{G}' \longrightarrow (f_{spaces, \etale})_*\mathcal{F}'
$$
Pour cela, soit $V \to Y$ un morphisme \'etale d'espaces algébriques.
Il faut construire une application d'ensembles
$$
\mathcal{G}'(V) \to \mathcal{F}'(X \times_Y V)
$$
Choisissons un morphisme \'etale surjectif $V' \to V$ avec $V'$ schéma,
puis un morphisme \'etale surjectif
$U' \to X \times_Y V'$ avec $U'$ schéma. Posons $U'' = U' \times_{X \times_Y V} U'$
et $V'' = V' \times_V V'$. Nous obtenons un morphisme de schémas $g' : U' \to V'$ ainsi qu'un morphisme de schémas
$$
g'' : U' \times_{X \times_Y V} U' \longrightarrow V' \times_V V'
$$
Considérons le diagramme suivant
$$
\xymatrix{
\mathcal{F}'(X \times_Y V) \ar[r] &
\mathcal{F}(U') \ar@<1ex>[r] \ar@<-1ex>[r] &
\mathcal{F}(U' \times_{X \times_Y V} U') \\
\mathcal{G}'(V) \ar[r] \ar@{..>}[u] &
\mathcal{G}(V') \ar@<1ex>[r] \ar@<-1ex>[r] \ar[u]_{\varphi_{(U', V', g')}} &
\mathcal{G}(V' \times_V V') \ar[u]_{\varphi_{(U'', V'', g'')}}
}
$$
La compatibilité des applications $\varphi_{...}$
aux restrictions montre que les deux carrés de droite sont commutatifs.
La définition des recouvrements de $X_{spaces, \etale}$ montre que
les lignes horizontales sont des diagrammes exacts. On obtient donc
la flèche pointillée. Nous laissons au lecteur le soin de montrer que ces
flèches sont compatibles aux applications de restriction.
\end{proof}

\noindent
Si le morphisme d'espaces algébriques $X \to Y$ est \'etale, alors le morphisme
de topos $\Sh(X_\etale) \to \Sh(Y_\etale)$
est un morphisme de localisation. Voici un énoncé précis.

\begin{lemma}
\label{lemma-etale-morphism-topoi}
Soit $S$ un schéma, et soit $f : X \to Y$ un morphisme d'espaces algébriques
sur $S$. Supposons $f$ \'etale. Il existe alors un foncteur
$$
j : X_\etale \to Y_\etale, \quad
(\varphi : U \to X) \mapsto (f \circ \varphi : U \to Y)
$$
qui est cocontinu. Le morphisme de topos $f_{small}$ est le
morphisme de topos associé à $j$, voir
Sites, lemme \ref{sites-lemma-cocontinuous-morphism-topoi}.
De plus, $j$ est aussi continu ; par conséquent,
Sites, lemme \ref{sites-lemma-when-shriek},
s'applique. En particulier, $f_{small}^{-1}\mathcal{G}(U) = \mathcal{G}(jU)$
pour tout faisceau $\mathcal{G}$ sur $Y_\etale$.
\end{lemma}

\begin{proof}
Notons que notre définition même d'un morphisme \'etale d'espaces algébriques
(définition \ref{definition-etale}) assure
que la règle donnée définit bien le foncteur $j$ indiqué.
Il est clair que $j$ est cocontinu et continu, simplement parce qu'un recouvrement
$\{U_i \to U\}$ de $j(\varphi : U \to X)$ dans $Y_\etale$ est
la même chose qu'un recouvrement de $(\varphi : U \to X)$ dans $X_\etale$. Il
reste à montrer que $j$ induit le même morphisme de topos que $f_{small}$.
Considérons pour cela le diagramme
$$
\xymatrix{
X_\etale \ar[r] \ar[d]^j &
X_{spaces, \etale} \ar@/_/[d]_{j_{spaces}} \\
Y_\etale \ar[r] &
Y_{spaces, \etale} \ar@/_/[u]_{v : V \mapsto X \times_Y V}
}
$$
de catégories. Ici, le foncteur $j_{spaces}$ est le prolongement évident de $j$
à la catégorie $X_{spaces, \etale}$. Le carré intérieur est donc
commutatif. En fait, $j_{spaces}$ s'identifie au
foncteur de localisation
$j_X : Y_{spaces, \etale}/X \to Y_{spaces, \etale}$
étudié dans
Sites, section \ref{sites-section-localize}.
Par conséquent, d'après
Sites, lemme \ref{sites-lemma-localize-given-products},
le foncteur cocontinu $j_{spaces}$ et le foncteur $v$ du diagramme
induisent le même morphisme de topos. D'après
Sites, lemme \ref{sites-lemma-composition-cocontinuous},
la commutativité du carré intérieur (formé de foncteurs cocontinus
entre sites) fournit un diagramme commutatif des morphismes de topos associés.
La construction de $f_{small}$ dans le
lemme \ref{lemma-functoriality-etale-site} achève donc la preuve.
\end{proof}

\noindent
Le lemme précédent dit que l'image réciproque de $\mathcal{G}$ par un morphisme \'etale
$f : X \to Y$ d'espaces algébriques est simplement la restriction de $\mathcal{G}$
à la catégorie $X_\etale$. Nous emploierons souvent l'abréviation
\begin{equation}
\label{equation-restrict}
\mathcal{G}|_{X_\etale} = f_{small}^{-1}\mathcal{G}
\end{equation}
pour le signaler. Notons que le foncteur
$j : X_\etale \to Y_\etale$
du lemme est alors fidèle, mais n'est pas pleinement fidèle en
général. Nous l'étudierons plus techniquement dans la
section \ref{section-localize}.

\begin{lemma}
\label{lemma-pushforward-etale-base-change}
Soit $S$ un schéma. Soit
$$
\xymatrix{
X' \ar[r] \ar[d]_{f'} & X \ar[d]^f \\
Y' \ar[r]^g & Y
}
$$
un carré cartésien d'espaces algébriques sur $S$. Soit
$\mathcal{F}$ un faisceau sur $X_\etale$. Si $g$ est \'etale, alors
\begin{enumerate}
\item $f'_{small, *}(\mathcal{F}|_{X'}) = (f_{small, *}\mathcal{F})|_{Y'}$
dans $\Sh(Y'_\etale)$\footnote{On a aussi
$(f')_{small}^{-1}(\mathcal{G}|_{Y'}) = (f_{small}^{-1}\mathcal{G})|_{X'}$
en raison de la commutativité du diagramme et de (\ref{equation-restrict})}, et
\item si $\mathcal{F}$ est un faisceau abélien, alors
$R^if'_{small, *}(\mathcal{F}|_{X'}) = (R^if_{small, *}\mathcal{F})|_{Y'}$.
\end{enumerate}
\end{lemma}

\begin{proof}
Considérons le diagramme de foncteurs suivant
$$
\xymatrix{
X'_{spaces, \etale} \ar[r]_j &
X_{spaces, \etale} \\
Y'_{spaces, \etale} \ar[r]^j \ar[u]^{V' \mapsto V' \times_{Y'} X'} &
Y_{spaces, \etale} \ar[u]_{V \mapsto V \times_Y X}
}
$$
Les flèches horizontales sont des foncteurs de localisation et les flèches verticales induisent
des morphismes de sites. La dernière assertion de
Sites, lemme \ref{sites-lemma-localize-morphism},
donne donc (1). Pour obtenir (2), appliquons (1) à une résolution injective de $\mathcal{F}$
et utilisons que la restriction est exacte et préserve les injectifs (voir
Cohomology on Sites, lemme \ref{sites-cohomology-lemma-cohomology-of-open}).
\end{proof}

\noindent
Le lemme suivant dit que l'on peut considérer un faisceau sur le petit
site \'etale d'un espace algébrique comme une famille compatible de faisceaux
sur les petits sites \'etales des schémas \'etales sur cet espace. Notons bien
que tous les morphismes de comparaison $c_f$ du lemme sont des isomorphismes,
ce qui est compatible avec
Topologies, lemme \ref{topologies-lemma-characterize-sheaf-big-etale},
et avec le fait que tous les morphismes entre objets de $X_\etale$
sont \'etales.

\begin{lemma}
\label{lemma-characterize-sheaf-small-etale}
Soit $S$ un schéma. Soit $X$ un espace algébrique sur $S$.
Un faisceau $\mathcal{F}$ sur $X_\etale$ équivaut aux données suivantes :
\begin{enumerate}
\item pour tout $U \in \Ob(X_\etale)$, un faisceau
$\mathcal{F}_U$ sur $U_\etale$ ;
\item pour tout $f : U' \to U$ dans $X_\etale$, un isomorphisme
$c_f : f_{small}^{-1}\mathcal{F}_U \to \mathcal{F}_{U'}$.
\end{enumerate}
Ces données sont soumises à la condition que, pour tous $f : U' \to U$
et $g : U'' \to U'$ dans $X_\etale$, la composée
$c_g \circ g_{small}^{-1} c_f$ soit égale à $c_{f \circ g}$.
\end{lemma}

\begin{proof}
On peut interpréter $g_{small}^{-1}$ comme dans le lemme \ref{lemma-etale-morphism-topoi}.
Le lemme résulte alors d'un fait général sur les sites, voir
Sites, lemme \ref{sites-lemma-glue-sheaves-absolute}.
\end{proof}

\noindent
Soit $S$ un schéma. Soit $X$ un espace algébrique sur $S$.
Soit $X = U/R$ une présentation de $X$ provenant d'un morphisme
\'etale surjectif quelconque $\varphi : U \to X$, voir
Spaces, définition \ref{spaces-definition-presentation}.
On obtient en particulier un groupoïde $(U, R, s, t, c, e, i)$ tel que
$j = (t, s) : R \to U \times_S U$, voir
Groupoids, lemme \ref{groupoids-lemma-equivalence-groupoid}.

\begin{lemma}
\label{lemma-descent-sheaf}
Avec $S$, $\varphi : U \to X$ et $(U, R, s, t, c, e, i)$ comme ci-dessus,
pour tout faisceau $\mathcal{F}$ sur $X_\etale$, le
faisceau\footnote{Dans ce lemme
et sa preuve, nous écrivons simplement $\varphi^{-1}$ au lieu de $\varphi_{small}^{-1}$,
et de même pour toutes les autres images réciproques.}
$\mathcal{G} = \varphi^{-1}\mathcal{F}$ est muni d'un
isomorphisme canonique
$$
\alpha :
t^{-1}\mathcal{G}
\longrightarrow
s^{-1}\mathcal{G}
$$
tel que le diagramme
$$
\xymatrix{
& \text{pr}_1^{-1}t^{-1}\mathcal{G} \ar[r]_-{\text{pr}_1^{-1}\alpha} &
\text{pr}_1^{-1}s^{-1}\mathcal{G} \ar@{=}[rd] & \\
\text{pr}_0^{-1}s^{-1}\mathcal{G} \ar@{=}[ru] & & &
c^{-1}s^{-1}\mathcal{G} \\
&
\text{pr}_0^{-1}t^{-1}\mathcal{G} \ar[lu]^{\text{pr}_0^{-1}\alpha} \ar@{=}[r] &
c^{-1}t^{-1}\mathcal{G} \ar[ru]_{c^{-1}\alpha}
}
$$
est commutatif. Le foncteur $\mathcal{F} \mapsto (\mathcal{G}, \alpha)$
définit une équivalence de catégories entre les faisceaux sur
$X_\etale$ et les couples $(\mathcal{G}, \alpha)$ ci-dessus.
\end{lemma}

\begin{proof}[Première preuve du lemme \ref{lemma-descent-sheaf}]
Soit $\mathcal{C} = X_{spaces, \etale}$. D'après le
lemme \ref{lemma-etale-morphism-topoi}
et sa preuve, on a $U_{spaces, \etale} = \mathcal{C}/U$
et le foncteur image réciproque $\varphi^{-1}$ est simplement le foncteur de restriction.
De plus, $\{U \to X\}$ est un recouvrement du site $\mathcal{C}$ et
$R = U \times_X U$. L'isomorphisme $\alpha$ est simplement l'identification
canonique
$$
\left(\mathcal{F}|_{\mathcal{C}/U}\right)|_{\mathcal{C}/U \times_X U}
=
\left(\mathcal{F}|_{\mathcal{C}/U}\right)|_{\mathcal{C}/U \times_X U}
$$
et la commutativité du diagramme est la condition de cocycle pour les données
de recollement. Ce lemme est donc un cas particulier du recollement des faisceaux, voir
Sites, section \ref{sites-section-glueing-sheaves}.
\end{proof}

\begin{proof}[Seconde preuve du lemme \ref{lemma-descent-sheaf}]
L'existence de $\alpha$ vient de l'égalité
$\varphi \circ t = \varphi \circ s$ et de la fonctorialité de l'image réciproque
par rapport au morphisme, voir le
lemme \ref{lemma-functoriality-etale-site}.
De la même manière, c'est-à-dire par fonctorialité de l'image réciproque, on voit
que l'isomorphisme $\alpha$ s'insère dans le diagramme commutatif.
La construction $\mathcal{F} \mapsto (\varphi^{-1}\mathcal{F}, \alpha)$
est manifestement fonctorielle en le faisceau $\mathcal{F}$.
On obtient ainsi le foncteur.

\medskip\noindent
Réciproquement, supposons que $(\mathcal{G}, \alpha)$ soit un couple.
Soit $V \to X$ un objet de $X_\etale$.
Alors le morphisme $V' = U \times_X V \to V$ est un morphisme \'etale surjectif
de schémas, et $\{V' \to V\}$ est donc un recouvrement \'etale
de $V$. Posons $\mathcal{G}' = (V' \to V)^{-1}\mathcal{G}$.
Puisque $R = U \times_X U$, avec $t = \text{pr}_0$
et $s = \text{pr}_1$, on a $V' \times_V V' = R \times_X V$,
les applications de projection $s', t' : V' \times_V V' \to V'$ étant obtenues par changement de base
à partir de $t$ et $s$. Ainsi, $\alpha$ s'induit en un isomorphisme
$\alpha' : (t')^{-1}\mathcal{G}' \to (s')^{-1}\mathcal{G}'$. Cela étant posé,
nous définissons simplement
$$
\xymatrix{
\mathcal{F}(V) \ar@{=}[r] &
\text{Noyau du couple}(\mathcal{G}(V') \ar@<1ex>[r] \ar@<-1ex>[r] &
\mathcal{G}(V' \times_V V')).
}
$$
Nous omettons de vérifier que cela définit un faisceau. Pour voir que
$\mathcal{G}(V) = \mathcal{F}(V)$ s'il existe un morphisme $V \to U$,
notons que le noyau du couple est alors
$H^0(\{V' \to V\}, \mathcal{G}) = \mathcal{G}(V)$.
\end{proof}







\section{Points of the small \'etale site}
\label{section-points-small-etale-site}

\noindent
This section is the analogue of
\'Etale Cohomology, Section
\ref{etale-cohomology-section-stalks}.

\begin{definition}
\label{definition-geometric-point}
Let $S$ be a scheme. Let $X$ be an algebraic space over $S$.
\begin{enumerate}
\item A {\it geometric point} of $X$ is a morphism
$\overline{x} : \Spec(k) \to X$, where $k$ is an algebraically
closed field. We often abuse notation and
write $\overline{x} = \Spec(k)$.
\item For every geometric point $\overline{x}$ we have the corresponding
``image'' point $x \in |X|$. We say that $\overline{x}$ is a
{\it geometric point lying over $x$}.
\end{enumerate}
\end{definition}

\noindent
It turns out that we can take stalks of sheaves on $X_\etale$
at geometric points exactly in the same way as was done in the
case of the small \'etale site of a scheme. In order to do this we
define the notion of an \'etale neighbourhood as follows.

\begin{definition}
\label{definition-etale-neighbourhood}
Let $S$ be a scheme. Let $X$ be an algebraic space over $S$.
Let $\overline{x}$ be a geometric point of $X$.
\begin{enumerate}
\item An {\it \'etale neighborhood} of $\overline{x}$
of $X$ is a commutative diagram
$$
\xymatrix{
& U \ar[d]^\varphi \\
{\bar x} \ar[r]^{\bar x} \ar[ur]^{\bar u} & X
}
$$
where $\varphi$ is an \'etale morphism of algebraic spaces over $S$.
We will use the notation $\varphi : (U, \overline{u}) \to (X, \overline{x})$
to indicate this situation.
\item A {\it morphism of \'etale neighborhoods}
$(U, \overline{u}) \to (U', \overline{u}')$
is an $X$-morphism $h : U \to U'$
such that $\overline{u}' = h \circ \overline{u}$.
\end{enumerate}
\end{definition}

\noindent
Note that we allow $U$ to be an algebraic space. When we take stalks
of a sheaf on $X_\etale$ we have to restrict to those $U$ which
are in $X_\etale$, and so in this case we will only consider the case
where $U$ is a scheme. Alternately we can work with the site
$X_{space, \etale}$ and consider all \'etale neighbourhoods. And there
won't be any difference because of the last assertion in the following lemma.

\begin{lemma}
\label{lemma-cofinal-etale}
Let $S$ be a scheme. Let $X$ be an algebraic space over $S$.
Let $\overline{x}$ be a geometric point of $X$.
The category of \'etale neighborhoods is cofiltered. More precisely:
\begin{enumerate}
\item Let $(U_i, \overline{u}_i)_{i = 1, 2}$ be two \'etale neighborhoods of
$\overline{x}$ in $X$. Then there exists a third \'etale neighborhood
$(U, \overline{u})$ and morphisms
$(U, \overline{u}) \to (U_i, \overline{u}_i)$, $i = 1, 2$.
\item Let $h_1, h_2: (U, \overline{u}) \to (U', \overline{u}')$ be two
morphisms between \'etale neighborhoods of $\overline{s}$. Then there exist an
\'etale neighborhood $(U'', \overline{u}'')$ and a morphism
$h : (U'', \overline{u}'') \to (U, \overline{u})$
which equalizes $h_1$ and $h_2$, i.e., such that
$h_1 \circ h = h_2 \circ h$.
\end{enumerate}
Moreover, given any \'etale neighbourhood
$(U, \overline{u}) \to (X, \overline{x})$
there exists a morphism of \'etale neighbourhoods
$(U', \overline{u}') \to (U, \overline{u})$
where $U'$ is a scheme.
\end{lemma}

\begin{proof}
For part (1), consider the fibre product $U = U_1 \times_X U_2$.
It is \'etale over both $U_1$ and $U_2$ because \'etale morphisms are
preserved under base change and composition, see
Lemmas \ref{lemma-base-change-etale} and \ref{lemma-composition-etale}.
The map $\overline{u} \to U$ defined by $(\overline{u}_1, \overline{u}_2)$
gives it the structure of an \'etale neighborhood mapping to both
$U_1$ and $U_2$.

\medskip\noindent
For part (2), define $U''$ as the fibre product
$$
\xymatrix{
U'' \ar[r] \ar[d] & U \ar[d]^{(h_1, h_2)} \\
U' \ar[r]^-\Delta & U' \times_X U'.
}
$$
Since $\overline{u}$ and $\overline{u}'$ agree over $X$ with $\overline{x}$,
we see that $\overline{u}'' = (\overline{u}, \overline{u}')$ is a geometric
point of $U''$. In particular $U'' \not = \emptyset$.
Moreover, since $U'$ is \'etale over $X$, so is the fibre product
$U'\times_X U'$ (as seen above in the case of $U_1 \times_X U_2$).
Hence the vertical arrow $(h_1, h_2)$ is \'etale by
Lemma \ref{lemma-etale-permanence}.
Therefore $U''$ is \'etale over $U'$ by base change, and hence also
\'etale over $X$ (because compositions of \'etale morphisms are \'etale).
Thus $(U'', \overline{u}'')$ is a solution to the problem posed by (2).

\medskip\noindent
To see the final assertion, choose any surjective \'etale morphism
$U' \to U$ where $U'$ is a scheme. Then
$U' \times_U \overline{u}$ is a scheme surjective and \'etale over
$\overline{u} = \Spec(k)$ with $k$ algebraically closed.
It follows (see
Morphisms, Lemma \ref{morphisms-lemma-etale-over-field})
that $U' \times_U \overline{u} \to \overline{u}$ has a section
which gives us the desired $\overline{u}'$.
\end{proof}

\begin{lemma}
\label{lemma-geometric-lift-to-usual}
Let $S$ be a scheme. Let $X$ be an algebraic space over $S$.
Let $\overline{x} : \Spec(k) \to X$ be a geometric point of $X$
lying over $x \in |X|$. Let $\varphi : U \to X$ be an \'etale morphism
of algebraic spaces and let $u \in |U|$ with $\varphi(u) = x$.
Then there exists a geometric point
$\overline{u} : \Spec(k) \to U$ lying over $u$ with
$\overline{x} = \varphi \circ \overline{u}$.
\end{lemma}

\begin{proof}
Choose an affine scheme $U'$ with $u' \in U'$ and an \'etale morphism
$U' \to U$ which maps $u'$ to $u$. If we can prove the lemma for
$(U', u') \to (X, x)$ then the lemma follows. Hence we may assume that
$U$ is a scheme, in particular that $U \to X$ is representable.
Then look at the cartesian diagram
$$
\xymatrix{
\Spec(k) \times_{\overline{x}, X, \varphi} U
\ar[d]_{\text{pr}_1} \ar[r]_-{\text{pr}_2} & U
\ar[d]^\varphi \\
\Spec(k) \ar[r]^-{\overline{x}} & X
}
$$
The projection $\text{pr}_1$ is the base change of an \'etale morphisms
so it is \'etale, see
Lemma \ref{lemma-base-change-etale}.
Therefore, the scheme $\Spec(k) \times_{\overline{x}, X, \varphi} U$
is a disjoint union of finite separable extensions of $k$, see
Morphisms, Lemma \ref{morphisms-lemma-etale-over-field}.
But $k$ is algebraically closed, so all these extensions are trivial,
so $\Spec(k) \times_{\overline{x}, X, \varphi} U$
is a disjoint union of copies of $\Spec(k)$ and each of
these corresponds to a geometric point $\overline{u}$ with
$\varphi \circ \overline{u} = \overline{x}$. By
Lemma \ref{lemma-points-cartesian}
the map
$$
|\Spec(k) \times_{\overline{x}, X, \varphi} U|
\longrightarrow
|\Spec(k)| \times_{|X|} |U|
$$
is surjective, hence we can pick $\overline{u}$ to lie over $u$.
\end{proof}

\begin{lemma}
\label{lemma-geometric-lift-to-cover}
Let $S$ be a scheme. Let $X$ be an algebraic space over $S$.
Let $\overline{x}$ be a geometric point of $X$.
Let $(U, \overline{u})$ an \'etale neighborhood of $\overline{x}$.
Let $\{\varphi_i : U_i \to U\}_{i \in I}$ be an \'etale covering in
$X_{spaces, \etale}$.
Then there exist $i \in I$ and $\overline{u}_i : \overline{x} \to U_i$
such that $\varphi_i : (U_i, \overline{u}_i) \to (U, \overline{u})$
is a morphism of \'etale neighborhoods.
\end{lemma}

\begin{proof}
Let $u \in |U|$ be the image of $\overline{u}$.
As $|U| = \bigcup_{i \in I} \varphi_i(|U_i|)$ there exists an
$i$ and a point $u_i \in U_i$ mapping to $x$. Apply
Lemma \ref{lemma-geometric-lift-to-usual}
to $(U_i, u_i) \to (U, u)$ and $\overline{u}$ to
get the desired geometric point.
\end{proof}

\begin{definition}
\label{definition-stalk}
Let $S$ be a scheme. Let $X$ be an algebraic space over $S$.
Let $\mathcal{F}$ be a presheaf on $X_\etale$.
Let $\overline{x}$ be a geometric point of $X$.
The {\it stalk} of $\mathcal{F}$ at $\overline{x}$ is
$$
\mathcal{F}_{\bar x}
=
\colim_{(U, \overline{u})} \mathcal{F}(U)
$$
where $(U, \overline{u})$ runs over all \'etale neighborhoods
of $\overline{x}$ in $X$ with $U \in \Ob(X_\etale)$.
\end{definition}

\noindent
By
Lemma \ref{lemma-cofinal-etale},
this colimit is over a filtered
index category, namely the opposite of the category of \'etale neighborhoods
in $X_\etale$. More precisely
Lemma \ref{lemma-cofinal-etale}
says the opposite of the category of all \'etale neighbourhoods is filtered,
and the full subcategory of those which are in $X_\etale$ is a cofinal
subcategory hence also filtered.

\medskip\noindent
This means an element of $\mathcal{F}_{\overline{x}}$ can be
thought of as a triple $(U, \overline{u}, \sigma)$ where
$U \in \Ob(X_\etale)$ and $\sigma \in \mathcal{F}(U)$.
Two triples $(U, \overline{u}, \sigma)$, $(U', \overline{u}', \sigma')$
define the same element of the stalk if there exists a third
\'etale neighbourhood
$(U'', \overline{u}'')$, $U'' \in \Ob(X_\etale)$
and morphisms of \'etale neighbourhoods
$h : (U'', \overline{u}'') \to (U, \overline{u})$,
$h' : (U'', \overline{u}'') \to (U', \overline{u}')$ such that
$h^*\sigma = (h')^*\sigma'$ in $\mathcal{F}(U'')$. See
Categories, Section \ref{categories-section-directed-colimits}.

\medskip\noindent
This also implies that if $\mathcal{F}'$ is the sheaf on
$X_{spaces, \etale}$ corresponding to $\mathcal{F}$ on
$X_\etale$, then
\begin{equation}
\label{equation-stalk-spaces-etale}
\mathcal{F}_{\overline{x}} = \colim_{(U, \overline{u})} \mathcal{F}'(U)
\end{equation}
where now the colimit is over all the \'etale neighbourhoods of $\overline{x}$.
We will often jump between the point of view of using $X_\etale$
and $X_{spaces, \etale}$ without further mention.

\medskip\noindent
In particular this means that if $\mathcal{F}$ is a presheaf of
abelian groups, rings, etc then $\mathcal{F}_{\overline{x}}$ is
an abelian group, ring, etc simply by the usual way of defining the
group structure on a directed colimit of abelian groups, rings, etc.

\begin{lemma}
\label{lemma-stalk-gives-point}
\begin{slogan}
A geometric point of an algebraic space gives a point of its \'etale topos.
\end{slogan}
Let $S$ be a scheme.
Let $X$ be an algebraic space over $S$.
Let $\overline{x}$ be a geometric point of $X$.
Consider the functor
$$
u : X_\etale \longrightarrow \textit{Sets}, \quad
U \longmapsto |U_{\overline{x}}|
$$
Then $u$ defines a point $p$ of the site $X_\etale$
(Sites, Definition \ref{sites-definition-point})
and its associated stalk functor $\mathcal{F} \mapsto \mathcal{F}_p$
(Sites, Equation \ref{sites-equation-stalk})
is the functor $\mathcal{F} \mapsto \mathcal{F}_{\overline{x}}$
defined above.
\end{lemma}

\begin{proof}
In the proof of
Lemma \ref{lemma-geometric-lift-to-usual}
we have seen that the scheme $U_{\overline{x}} = \overline{x} \times_X U$
is a disjoint union of
schemes isomorphic to $\overline{x}$. Thus we can also think of
$|U_{\overline{x}}|$ as the set of geometric points of $U$ lying over
$\overline{x}$, i.e., as the collection of morphisms
$\overline{u} : \overline{x} \to U$ fitting into the diagram of
Definition \ref{definition-etale-neighbourhood}.
From this it follows that $u(X)$ is a singleton, and that
$u(U \times_V W) = u(U) \times_{u(V)} u(W)$
whenever $U \to V$ and $W \to V$ are morphisms in $X_\etale$.
And, given a covering $\{U_i \to U\}_{i \in I}$ in $X_\etale$ we see
that $\coprod u(U_i) \to u(U)$ is surjective by
Lemma \ref{lemma-geometric-lift-to-cover}.
Hence
Sites, Proposition \ref{sites-proposition-point-limits}
applies, so $p$ is a point of the site $X_\etale$.
Finally, the our functor $\mathcal{F} \mapsto \mathcal{F}_{\overline{s}}$
is given by exactly the same colimit as the functor
$\mathcal{F} \mapsto \mathcal{F}_p$ associated to $p$ in
Sites, Equation \ref{sites-equation-stalk}
which proves the final assertion.
\end{proof}

\begin{lemma}
\label{lemma-stalk-exact}
Let $S$ be a scheme.
Let $X$ be an algebraic space over $S$.
Let $\overline{x}$ be a geometric point of $X$.
\begin{enumerate}
\item The stalk functor
$\textit{PAb}(X_\etale) \to \textit{Ab}$,
$\mathcal{F}  \mapsto  \mathcal{F}_{\overline{x}}$
is exact.
\item We have $(\mathcal{F}^\#)_{\overline{x}} = \mathcal{F}_{\overline{x}}$
for any presheaf of sets $\mathcal{F}$ on $X_\etale$.
\item The functor
$\textit{Ab}(X_\etale) \to \textit{Ab}$,
$\mathcal{F} \mapsto \mathcal{F}_{\overline{x}}$ is exact.
\item Similarly the functors
$\textit{PSh}(X_\etale) \to \textit{Sets}$ and
$\Sh(X_\etale) \to \textit{Sets}$ given by the stalk functor
$\mathcal{F} \mapsto \mathcal{F}_{\overline{x}}$ are exact (see
Categories, Definition \ref{categories-definition-exact})
and commute with arbitrary colimits.
\end{enumerate}
\end{lemma}

\begin{proof}
This result follows from the general material in
Modules on Sites, Section \ref{sites-modules-section-stalks}.
This is true because $\mathcal{F} \mapsto \mathcal{F}_{\overline{x}}$
comes from a point of the small \'etale site of $X$, see
Lemma \ref{lemma-stalk-gives-point}. See the proof of
\'Etale Cohomology, Lemma \ref{etale-cohomology-lemma-stalk-exact}
for a direct proof of some of these statements in the setting of
the small \'etale site of a scheme.
\end{proof}

\noindent
We will see below that the stalk functor
$\mathcal{F} \mapsto \mathcal{F}_{\overline{x}}$
is really the pullback along the morphism $\overline{x}$. In that sense
the following lemma is a generalization of the lemma above.

\begin{lemma}
\label{lemma-stalk-pullback}
Let $S$ be a scheme.
Let $f : X \to Y$ be a morphism of algebraic spaces over $S$.
\begin{enumerate}
\item The functor
$f_{small}^{-1} :
\textit{Ab}(Y_\etale)
\to
\textit{Ab}(X_\etale)$
is exact.
\item The functor
$f_{small}^{-1} :
\Sh(Y_\etale)
\to
\Sh(X_\etale)$
is exact, i.e., it commutes with finite limits and colimits, see
Categories, Definition \ref{categories-definition-exact}.
\item For any \'etale morphism $V \to Y$ of algebraic spaces
we have $f_{small}^{-1}h_V = h_{X \times_Y V}$.
\item Let $\overline{x} \to X$ be a geometric point.
Let $\mathcal{G}$ be a sheaf on $Y_\etale$.
Then there is a canonical identification
$$
(f_{small}^{-1}\mathcal{G})_{\overline{x}} = \mathcal{G}_{\overline{y}}.
$$
where $\overline{y} = f \circ \overline{x}$.
\end{enumerate}
\end{lemma}

\begin{proof}
Recall that $f_{small}$ is defined via $f_{spaces, small}$ in
Lemma \ref{lemma-functoriality-etale-site}.
Parts (1), (2) and (3) are general consequences of the fact that
$f_{spaces, \etale} :
X_{spaces, \etale}
\to
Y_{spaces, \etale}$
is a morphism of sites, see
Sites, Definition \ref{sites-definition-morphism-sites}
for (2),
Modules on Sites, Lemma \ref{sites-modules-lemma-flat-pullback-exact}
for (1), and
Sites, Lemma \ref{sites-lemma-pullback-representable-sheaf}
for (3).

\medskip\noindent
Proof of (4). This statement is a special case of
Sites, Lemma \ref{sites-lemma-point-morphism-sites}
via
Lemma \ref{lemma-stalk-gives-point}.
We also provide a direct proof. Note that by
Lemma \ref{lemma-stalk-exact}.
taking stalks commutes with sheafification.
Let $\mathcal{G}'$ be the sheaf on $Y_{spaces, \etale}$ whose restriction
to $Y_\etale$ is $\mathcal{G}$.
Recall that $f_{spaces, \etale}^{-1}\mathcal{G}'$ is the sheaf
associated to the presheaf
$$
U \longrightarrow \colim_{U \to X \times_Y V} \mathcal{G}'(V),
$$
see
Sites, Sections \ref{sites-section-continuous-functors} and
\ref{sites-section-functoriality-PSh}.
Thus we have
\begin{align*}
(f_{spaces, \etale}^{-1}\mathcal{G}')_{\overline{x}}
& =
\colim_{(U, \overline{u})} f_{spaces, \etale}^{-1}\mathcal{G}'(U)
\\
& = \colim_{(U, \overline{u})}
\colim_{a : U \to X \times_Y V} \mathcal{G}'(V) \\
& = \colim_{(V, \overline{v})} \mathcal{G}'(V) \\
& = \mathcal{G}'_{\overline{y}}
\end{align*}
in the third equality the pair $(U, \overline{u})$ and the map
$a : U \to X \times_Y V$ corresponds to the pair $(V, a \circ \overline{u})$.
Since the stalk of $\mathcal{G}'$
(resp.\ $f_{spaces, \etale}^{-1}\mathcal{G}'$)
agrees with the stalk of $\mathcal{G}$ (resp.\ $f_{small}^{-1}\mathcal{G}$),
see
Equation (\ref{equation-stalk-spaces-etale})
the result follows.
\end{proof}

\begin{remark}
\label{remark-stalk-pullback}
This remark is the analogue of
\'Etale Cohomology, Remark \ref{etale-cohomology-remark-stalk-pullback}.
Let $S$ be a scheme.
Let $X$ be an algebraic space over $S$.
Let $\overline{x} : \Spec(k) \to X$ be a geometric point of $X$.
By
\'Etale Cohomology,
Theorem \ref{etale-cohomology-theorem-equivalence-sheaves-point}
the category of sheaves on $\Spec(k)_\etale$ is
equivalent to the category of sets (by taking a sheaf to its global sections).
Hence it follows from
Lemma \ref{lemma-stalk-pullback} part (4)
applied to the morphism $\overline{x}$ that the functor
$$
\Sh(X_\etale) \longrightarrow \textit{Sets}, \quad
\mathcal{F} \longmapsto \mathcal{F}_{\overline{x}}
$$
is isomorphic to the functor
$$
\Sh(X_\etale)
\longrightarrow
\Sh(\Spec(k)_\etale) = \textit{Sets},
\quad
\mathcal{F} \longmapsto \overline{x}^*\mathcal{F}
$$
Hence we may view the stalk functors as pullback functors along
geometric morphisms (and not just some abstract morphisms of topoi
as in the result of
Lemma \ref{lemma-stalk-gives-point}).
\end{remark}

\begin{remark}
\label{remark-map-stalks}
Let $S$ be a scheme.
Let $X$ be an algebraic space over $S$.
Let $x \in |X|$.
We claim that for any pair of geometric points $\overline{x}$ and
$\overline{x}'$ lying over $x$ the stalk functors are isomorphic.
By definition of $|X|$ we can find a third geometric point
$\overline{x}''$ so that there exists a commutative diagram
$$
\xymatrix{
\overline{x}'' \ar[r] \ar[d] \ar[rd]^{\overline{x}''} &
\overline{x}' \ar[d]^{\overline{x}'} \\
\overline{x} \ar[r]^{\overline{x}} & X.
}
$$
Since the stalk functor $\mathcal{F} \mapsto \mathcal{F}_{\overline{x}}$
is given by pullback along the morphism $\overline{x}$ (and similarly for
the others) we conclude by functoriality of pullbacks.
\end{remark}




\noindent
The following theorem says that the small \'etale site of an algebraic
space has enough points.

\begin{theorem}
\label{theorem-exactness-stalks}
Let $S$ be a scheme. Let $X$ be an algebraic space over $S$.
A map $a : \mathcal{F} \to \mathcal{G}$ of sheaves of sets is injective
(resp.\ surjective) if and only if the map on stalks
$a_{\overline{x}} : \mathcal{F}_{\overline{x}} \to \mathcal{G}_{\overline{x}}$
is injective (resp.\ surjective) for all geometric points of $X$.
A sequence of abelian sheaves on $X_\etale$ is exact
if and only if it is exact on all stalks at geometric points of $S$.
\end{theorem}

\begin{proof}
We know the theorem is true if $X$ is a scheme, see
\'Etale Cohomology, Theorem \ref{etale-cohomology-theorem-exactness-stalks}.
Choose a surjective \'etale morphism $f : U \to X$ where $U$ is a scheme.
Since $\{U \to X\}$ is a covering (in $X_{spaces, \etale}$) we can check
whether a map of sheaves is injective, or surjective by restricting
to $U$. Now if $\overline{u} : \Spec(k) \to U$ is a geometric
point of $U$, then
$(\mathcal{F}|_U)_{\overline{u}} = \mathcal{F}_{\overline{x}}$
where $\overline{x} = f \circ \overline{u}$. (This is clear from the
colimits defining the stalks at $\overline{u}$ and $\overline{x}$, but
it also follows from
Lemma \ref{lemma-stalk-pullback}.)
Hence the result for $U$ implies the result for $X$ and we win.
\end{proof}

\noindent
The following lemma should be skipped on a first reading.

\begin{lemma}
\label{lemma-points-small-etale-site}
Let $S$ be a scheme.
Let $X$ be an algebraic space over $S$.
Let $p : \Sh(pt) \to \Sh(X_\etale)$
be a point of the small \'etale topos of $X$.
Then there exists a geometric point $\overline{x}$ of $X$
such that the stalk functor $\mathcal{F} \mapsto \mathcal{F}_p$
is isomorphic to the stalk functor
$\mathcal{F} \mapsto \mathcal{F}_{\overline{x}}$.
\end{lemma}

\begin{proof}
By
Sites, Lemma \ref{sites-lemma-point-site-topos}
there is a one to one correspondence between points of the site and points
of the associated topos. Hence we may assume that $p$ is given by
a functor $u : X_\etale \to \textit{Sets}$ which defines a point
of the site $X_\etale$.
Let $U \in \Ob(X_\etale)$ be an object whose structure morphism
$j : U \to X$ is surjective. Note that $h_U$ is a sheaf
which surjects onto the final sheaf. Since taking stalks is exact
we see that $(h_U)_p = u(U)$ is not empty (use
Sites, Lemma \ref{sites-lemma-points-recover}).
Pick $x \in u(U)$. By
Sites, Lemma \ref{sites-lemma-point-localize}
we obtain a point $q : \Sh(pt) \to \Sh(U_\etale)$
such that $p = j_{small} \circ q$, so that
$\mathcal{F}_p = (\mathcal{F}|_U)_q$ functorially.
By
\'Etale Cohomology, Lemma \ref{etale-cohomology-lemma-points-small-etale-site}
there is a geometric point $\overline{u}$ of $U$ and a functorial
isomorphism $\mathcal{G}_q = \mathcal{G}_{\overline{u}}$
for $\mathcal{G} \in \Sh(U_\etale)$. Set
$\overline{x} = j \circ \overline{u}$. Then we see that
$\mathcal{F}_{\overline{x}} \cong (\mathcal{F}|_U)_{\overline{u}}$
functorially in $\mathcal{F}$ on $X_\etale$ by
Lemma \ref{lemma-stalk-pullback}
and we win.
\end{proof}





\section{Supports of abelian sheaves}
\label{section-support}

\noindent
First we talk about supports of local sections.

\begin{lemma}
\label{lemma-support-subsheaf-final}
Let $S$ be a scheme. Let $X$ be an algebraic space over $S$.
Let $\mathcal{F}$ be a subsheaf of the final object of the \'etale
topos of $X$ (see
Sites, Example \ref{sites-example-singleton-sheaf}).
Then there exists a unique open
$W \subset X$ such that $\mathcal{F} = h_W$.
\end{lemma}

\begin{proof}
The condition means that $\mathcal{F}(U)$ is a singleton or
empty for all $\varphi : U \to X$ in $\Ob(X_{spaces, \etale})$.
In particular local sections always glue. If
$\mathcal{F}(U) \not = \emptyset$, then
$\mathcal{F}(\varphi(U)) \not = \emptyset$ because
$\varphi(U) \subset X$ is an open subspace
(Lemma \ref{lemma-etale-open})
and
$\{\varphi : U \to \varphi(U)\}$ is a covering in $X_{spaces, \etale}$.
Take
$W = \bigcup_{\varphi : U \to S, \mathcal{F}(U) \not = \emptyset} \varphi(U)$
to conclude.
\end{proof}

\begin{lemma}
\label{lemma-zero-over-image}
Let $S$ be a scheme.
Let $X$ be an algebraic space over $S$.
Let $\mathcal{F}$ be an abelian sheaf on $X_{spaces, \etale}$.
Let $\sigma \in \mathcal{F}(U)$ be a local section.
There exists an open subspace $W \subset U$ such that
\begin{enumerate}
\item $W \subset U$ is the largest open subspace of $U$ such
that $\sigma|_W = 0$,
\item for every $\varphi : V \to U$ in $X_{spaces, \etale}$ we have
$$
\sigma|_V = 0 \Leftrightarrow \varphi(V) \subset W,
$$
\item for every geometric point $\overline{u}$ of $U$ we have
$$
(U, \overline{u}, \sigma) = 0\text{ in }\mathcal{F}_{\overline{x}}
\Leftrightarrow
\overline{u} \in W
$$
where $\overline{x} = (U \to X) \circ \overline{u}$.
\end{enumerate}
\end{lemma}

\begin{proof}
Since $\mathcal{F}$ is a sheaf in the \'etale topology the restriction of
$\mathcal{F}$ to $U_{Zar}$ is a sheaf on $U$ in the Zariski topology.
Hence there exists a Zariski open $W$ having property (1), see
Modules, Lemma \ref{modules-lemma-support-section-closed}. Let
$\varphi : V \to U$ be an arrow of $X_{spaces, \etale}$. Note that
$\varphi(V) \subset U$ is an open subspace
(Lemma \ref{lemma-etale-open})
and that $\{V \to \varphi(V)\}$ is an \'etale covering. Hence if
$\sigma|_V = 0$, then by the sheaf condition for $\mathcal{F}$ we
see that $\sigma|_{\varphi(V)} = 0$. This proves (2).
To prove (3) we have to show that if $(U, \overline{u}, \sigma)$
defines the zero element of $\mathcal{F}_{\overline{x}}$, then
$\overline{u} \in W$. This is true because the assumption means
there exists a morphism of \'etale neighbourhoods
$(V, \overline{v}) \to (U, \overline{u})$ such that
$\sigma|_V = 0$. Hence by (2) we see that $V \to U$ maps into $W$, and
hence $\overline{u} \in W$.
\end{proof}

\noindent
Let $S$ be a scheme.
Let $X$ be an algebraic space over $S$.
Let $x \in |X|$.
Let $\mathcal{F}$ be a sheaf on $X_\etale$. By
Remark \ref{remark-map-stalks}
the isomorphism class of the stalk of the sheaf $\mathcal{F}$
at a geometric points lying over $x$ is well defined.

\begin{definition}
\label{definition-support}
Let $S$ be a scheme.
Let $X$ be an algebraic space over $S$.
Let $\mathcal{F}$ be an abelian sheaf on $X_\etale$.
\begin{enumerate}
\item The {\it support of $\mathcal{F}$} is the set of
points $x \in |X|$ such that $\mathcal{F}_{\overline{x}} \not = 0$
for any (some) geometric point $\overline{x}$ lying over $x$.
\item Let $\sigma \in \mathcal{F}(U)$ be a section.
The {\it support of $\sigma$} is the closed subset $U \setminus W$, where
$W \subset U$ is the largest open subset of $U$ on which $\sigma$
restricts to zero (see
Lemma \ref{lemma-zero-over-image}).
\end{enumerate}
\end{definition}

\begin{lemma}
\label{lemma-support-section-closed}
Let $S$ be a scheme.
Let $X$ be an algebraic space over $S$.
Let $\mathcal{F}$ be an abelian sheaf on $X_\etale$.
Let $U \in \Ob(X_\etale)$ and $\sigma \in \mathcal{F}(U)$.
\begin{enumerate}
\item The support of $\sigma$ is closed in $|X|$.
\item The support of $\sigma + \sigma'$ is contained in the union of
the supports of $\sigma, \sigma' \in \mathcal{F}(X)$.
\item If $\varphi : \mathcal{F} \to \mathcal{G}$ is a map of
abelian sheaves on $X_\etale$, then the support of $\varphi(\sigma)$ is
contained in the support of $\sigma \in \mathcal{F}(U)$.
\item The support of $\mathcal{F}$ is the union of the images of the
supports of all local sections of $\mathcal{F}$.
\item If $\mathcal{F} \to \mathcal{G}$ is surjective then the support
of $\mathcal{G}$ is a subset of the support of $\mathcal{F}$.
\item If $\mathcal{F} \to \mathcal{G}$ is injective then the support
of $\mathcal{F}$ is a subset of the support of $\mathcal{G}$.
\end{enumerate}
\end{lemma}

\begin{proof}
Part (1) holds by definition.
Parts (2) and (3) hold because they holds for the restriction of
$\mathcal{F}$ and $\mathcal{G}$ to $U_{Zar}$, see
Modules, Lemma \ref{modules-lemma-support-section-closed}.
Part (4) is a direct consequence of
Lemma \ref{lemma-zero-over-image} part (3).
Parts (5) and (6) follow from the other parts.
\end{proof}

\begin{lemma}
\label{lemma-support-sheaf-rings-closed}
The support of a sheaf of rings on the small \'etale site of an
algebraic space is closed.
\end{lemma}

\begin{proof}
This is true because (according to our conventions)
a ring is $0$ if and only if
$1 = 0$, and hence the support of a sheaf of rings
is the support of the unit section.
\end{proof}






\section{The structure sheaf of an algebraic space}
\label{section-structure-sheaf}

\noindent
The structure sheaf of an algebraic space is the sheaf of rings of the
following lemma.

\begin{lemma}
\label{lemma-sheaf-condition-holds}
Let $S$ be a scheme. Let $X$ be an algebraic space over $S$.
The rule $U \mapsto \Gamma(U, \mathcal{O}_U)$ defines
a sheaf of rings on $X_\etale$.
\end{lemma}

\begin{proof}
Immediate from the definition of a covering and
Descent, Lemma \ref{descent-lemma-sheaf-condition-holds}.
\end{proof}

\begin{definition}
\label{definition-structure-sheaf}
Let $S$ be a scheme.
Let $X$ be an algebraic space over $S$.
The {\it structure sheaf} of $X$
is the sheaf of rings $\mathcal{O}_X$
on the small \'etale site $X_\etale$ described in
Lemma \ref{lemma-sheaf-condition-holds}.
\end{definition}

\noindent
According to Lemma \ref{lemma-characterize-sheaf-small-etale} the sheaf
$\mathcal{O}_X$ corresponds to a system of \'etale sheaves $(\mathcal{O}_X)_U$
for $U$ ranging through the objects of $X_\etale$. It is clear from
the proof of that lemma and our definition that we have simply
$(\mathcal{O}_X)_U = \mathcal{O}_U$ where $\mathcal{O}_U$ is the structure
sheaf of $U_\etale$ as introduced in
Descent, Definition \ref{descent-definition-structure-sheaf}.
In particular, if $X$ is a scheme we recover the sheaf $\mathcal{O}_X$
on the small \'etale site of $X$.

\medskip\noindent
Via the equivalence
$\Sh(X_\etale) = \Sh(X_{spaces, \etale})$
of Lemma \ref{lemma-compare-etale-sites} we may also think of $\mathcal{O}_X$
as a sheaf of rings on $X_{spaces, \etale}$. It is explained in
Remark \ref{remark-explain-equivalence}
how to compute $\mathcal{O}_X(Y)$, and in particular $\mathcal{O}_X(X)$, when
$Y \to X$ is an object of $X_{spaces, \etale}$.

\begin{lemma}
\label{lemma-morphism-ringed-topoi}
Let $S$ be a scheme.
Let $f : X \to Y$ be a morphism of algebraic spaces over $S$.
Then there is a canonical map
$f^\sharp : f_{small}^{-1}\mathcal{O}_Y \to \mathcal{O}_X$ such that
$$
(f_{small}, f^\sharp) :
(\Sh(X_\etale), \mathcal{O}_X)
\longrightarrow
(\Sh(Y_\etale), \mathcal{O}_Y)
$$
is a morphism of ringed topoi. Furthermore,
\begin{enumerate}
\item The construction $f \mapsto (f_{small}, f^\sharp)$ is compatible with
compositions.
\item If $f$ is a morphism of schemes, then $f^\sharp$ is the map described in
Descent, Remark \ref{descent-remark-change-topologies-ringed}.
\end{enumerate}
\end{lemma}

\begin{proof}
By Lemma \ref{lemma-f-map} it suffices to give an $f$-map from
$\mathcal{O}_Y$ to $\mathcal{O}_X$. In other words, for every
commutative diagram
$$
\xymatrix{
U \ar[d]_g \ar[r] & X \ar[d]^f \\
V \ar[r] & Y
}
$$
where $U \in X_\etale$, $V \in Y_\etale$ we have to give a
map of rings
$
(f^\sharp)_{(U, V, g)} :
\Gamma(V, \mathcal{O}_V)
\to
\Gamma(U, \mathcal{O}_U).
$
Of course we just take $(f^\sharp)_{(U, V, g)} = g^\sharp$.
It is clear that this is compatible with restriction mappings
and hence indeed gives an $f$-map.
We omit checking compatibility with compositions and agreement with the
construction in
Descent, Remark \ref{descent-remark-change-topologies-ringed}.
\end{proof}

\begin{lemma}
\label{lemma-reduced-space}
Let $S$ be a scheme. Let $X$ be an algebraic space over $S$.
The following are equivalent
\begin{enumerate}
\item $X$ is reduced,
\item for every $x \in |X|$ the local ring of $X$ at $x$ is
reduced (Remark \ref{remark-list-properties-local-ring-local-etale-topology}).
\end{enumerate}
In this case $\Gamma(X, \mathcal{O}_X)$ is a reduced ring and
if $f \in \Gamma(X, \mathcal{O}_X)$ has $X = V(f)$, then $f = 0$.
\end{lemma}

\begin{proof}
The equivalence of (1) and (2) follows from
Properties, Lemma \ref{properties-lemma-characterize-reduced}
applied to affine schemes \'etale over $X$. The final statements
follow the cited lemma and fact that $\Gamma(X, \mathcal{O}_X)$ is
a subring of $\Gamma(U, \mathcal{O}_U)$ for some
reduced scheme $U$ \'etale over $X$.
\end{proof}







\section{Stalks of the structure sheaf}
\label{section-stalks-structure-sheaf}

\noindent
This section is the analogue of
\'Etale Cohomology, Section
\ref{etale-cohomology-section-stalks-structure-sheaf}.

\begin{lemma}
\label{lemma-describe-etale-local-ring}
Let $S$ be a scheme.
Let $X$ be an algebraic space over $S$.
Let $\overline{x}$ be a geometric point of $X$.
Let $(U, \overline{u})$ be an \'etale neighbourhood of $\overline{x}$
where $U$ is a scheme. Then we have
$$
\mathcal{O}_{X, \overline{x}} =
\mathcal{O}_{U, \overline{u}} =
\mathcal{O}_{U, u}^{sh}
$$
where the left hand side is the stalk of the structure sheaf of $X$,
and the right hand side is the strict henselization of the local ring
of $U$ at the point $u$ at which $\overline{u}$ is centered.
\end{lemma}

\begin{proof}
We know that the structure sheaf $\mathcal{O}_U$ on
$U_\etale$ is the restriction of the structure sheaf of $X$.
Hence the first equality follows from
Lemma \ref{lemma-stalk-pullback} part (4).
The second equality is explained in
\'Etale Cohomology,
Lemma \ref{etale-cohomology-lemma-describe-etale-local-ring}.
\end{proof}

\begin{definition}
\label{definition-etale-local-rings}
Let $S$ be a scheme.
Let $X$ be an algebraic space over $S$.
Let $\overline{x}$ be a geometric point of $X$ lying over the point
$x \in |X|$.
\begin{enumerate}
\item The {\it \'etale local ring of $X$ at $\overline{x}$}
is the stalk of the structure sheaf $\mathcal{O}_X$ on $X_\etale$
at $\overline{x}$.
Notation: $\mathcal{O}_{X, \overline{x}}$.
\item The {\it strict henselization of $X$ at $\overline{x}$}
is the scheme $\Spec(\mathcal{O}_{X, \overline{x}})$.
\end{enumerate}
\end{definition}

\noindent
The isomorphism type of the strict henselization of $X$ at $\overline{x}$
(as a scheme over $X$) depends only on the point $x \in |X|$ and not on
the choice of the geometric point lying over $x$, see
Remark \ref{remark-map-stalks}.

\begin{lemma}
\label{lemma-etale-site-locally-ringed}
Let $S$ be a scheme.
Let $X$ be an algebraic space over $S$.
The small \'etale site $X_\etale$ endowed with its
structure sheaf $\mathcal{O}_X$ is a locally ringed site, see
Modules on Sites, Definition \ref{sites-modules-definition-locally-ringed}.
\end{lemma}

\begin{proof}
This follows because the stalks
$\mathcal{O}_{X, \overline{x}}$ are
local, and because $S_\etale$ has enough points, see
Lemmas \ref{lemma-describe-etale-local-ring} and
Theorem \ref{theorem-exactness-stalks}.
See
Modules on Sites, Lemma \ref{sites-modules-lemma-locally-ringed-stalk} and
\ref{sites-modules-lemma-ringed-stalk-not-zero}
for the fact that this implies the small \'etale site is locally ringed.
\end{proof}

\begin{lemma}
\label{lemma-dimension-local-ring}
Let $S$ be a scheme. Let $X$ be an algebraic space over $S$.
Let $x \in |X|$ be a point. Let $d \in \{0, 1, 2, \ldots, \infty\}$.
The following are equivalent
\begin{enumerate}
\item the dimension of the local ring of $X$ at $x$
(Definition \ref{definition-dimension-local-ring}) is $d$,
\item $\dim(\mathcal{O}_{X, \overline{x}}) = d$ for some geometric
point $\overline{x}$ lying over $x$, and
\item $\dim(\mathcal{O}_{X, \overline{x}}) = d$ for any geometric
point $\overline{x}$ lying over $x$.
\end{enumerate}
\end{lemma}

\begin{proof}
The equivalence of (2) and (3) follows from the fact that the
isomorphism type of $\mathcal{O}_{X, \overline{x}}$ only depends
on $x \in |X|$, see Remark \ref{remark-map-stalks}.
Using Lemma \ref{lemma-describe-etale-local-ring}
the equivalence of (1) and (2)$+$(3) comes down to the
following statement: Given any local ring $R$ we have
$\dim(R) = \dim(R^{sh})$. This is
More on Algebra, Lemma \ref{more-algebra-lemma-henselization-dimension}.
\end{proof}

\begin{lemma}
\label{lemma-dimension-decent-invariant-under-etale}
Let $S$ be a scheme. Let $f : X \to Y$ be an \'etale morphism
of algebraic spaces over $S$. Let $x \in X$. Then
(1) $\dim_x(X) = \dim_{f(x)}(Y)$ and (2) the dimension of
the local ring of $X$ at $x$ equals the dimension of
the local ring of $Y$ at $f(x)$. If $f$ is surjective, then
(3) $\dim(X) = \dim(Y)$.
\end{lemma}

\begin{proof}
Choose a scheme $U$ and a point $u \in U$ and an \'etale morphism
$U \to X$ which maps $u$ to $x$. Then the composition $U \to Y$
is also \'etale and maps $u$ to $f(x)$. Thus the statements (1) and (2)
follow as the relevant integers are defined in terms of the behaviour
of the scheme $U$ at $u$. See
Definition \ref{definition-dimension-at-point} for (1). Part (3) is
an immediate consequence of (1), see Definition \ref{definition-dimension}.
\end{proof}

\begin{lemma}
\label{lemma-reduced-local-ring}
Let $S$ be a scheme. Let $X$ be an algebraic space over $S$.
Let $x \in |X|$ be a point. The following are equivalent
\begin{enumerate}
\item the local ring of $X$ at $x$ is reduced
(Remark \ref{remark-list-properties-local-ring-local-etale-topology}),
\item $\mathcal{O}_{X, \overline{x}}$ is reduced for some geometric
point $\overline{x}$ lying over $x$, and
\item $\mathcal{O}_{X, \overline{x}}$ is reduced for any geometric
point $\overline{x}$ lying over $x$.
\end{enumerate}
\end{lemma}

\begin{proof}
The equivalence of (2) and (3) follows from the fact that the
isomorphism type of $\mathcal{O}_{X, \overline{x}}$ only depends
on $x \in |X|$, see Remark \ref{remark-map-stalks}.
Using Lemma \ref{lemma-describe-etale-local-ring}
the equivalence of (1) and (2)$+$(3) comes down to the
following statement: a local ring is reduced if and only if
its strict henselization is reduced. This is
More on Algebra, Lemma \ref{more-algebra-lemma-henselization-reduced}.
\end{proof}











\section{Local irreducibility}
\label{section-irreducible-local-ring}

\noindent
A point on an algebraic space has a well defined \'etale local ring, which
corresponds to the strict henselization of the local ring in the case of a
scheme. In general we cannot see how many irreducible components of a
scheme or an algebraic space pass through the given point from the
\'etale local ring. We can only count the number of geometric branches.

\begin{lemma}
\label{lemma-irreducible-local-ring}
Let $S$ be a scheme.
Let $X$ be an algebraic space over $S$.
Let $x \in |X|$ be a point.
The following are equivalent
\begin{enumerate}
\item for any scheme $U$ and \'etale morphism $a : U \to X$ and
$u \in U$ with $a(u) = x$ the local ring $\mathcal{O}_{U, u}$ has a
unique minimal prime,
\item for any scheme $U$ and \'etale morphism $a : U \to X$ and
$u \in U$ with $a(u) = x$ there is a unique irreducible component of $U$
through $u$,
\item for any scheme $U$ and \'etale morphism $a : U \to X$ and
$u \in U$ with $a(u) = x$ the local ring $\mathcal{O}_{U, u}$
is unibranch,
\item for any scheme $U$ and \'etale morphism $a : U \to X$ and
$u \in U$ with $a(u) = x$ the local ring $\mathcal{O}_{U, u}$
is geometrically unibranch,
\item $\mathcal{O}_{X, \overline{x}}$ has a unique minimal prime
for any geometric point $\overline{x}$ lying over $x$.
\end{enumerate}
\end{lemma}

\begin{proof}
The equivalence of (1) and (2) follows from the fact that irreducible
components of $U$ passing through $u$ are in $1$-$1$ correspondence with
minimal primes of the local ring of $U$ at $u$. Let $a : U \to X$ and
$u \in U$ be as in (1). Then $\mathcal{O}_{X, \overline{x}}$ is
the strict henselization of $\mathcal{O}_{U, u}$ by
Lemma \ref{lemma-describe-etale-local-ring}.
In particular (4) and (5) are equivalent by
More on Algebra, Lemma \ref{more-algebra-lemma-geometrically-unibranch}.
The equivalence of (2), (3), and (4) follows from
More on Morphisms, Lemma \ref{more-morphisms-lemma-nr-branches}.
\end{proof}

\begin{definition}
\label{definition-unibranch}
Let $S$ be a scheme. Let $X$ be an algebraic space over $S$.
Let $x \in |X|$. We say that $X$ is {\it geometrically unibranch
at $x$} if the equivalent conditions of
Lemma \ref{lemma-irreducible-local-ring}
hold. We say that $X$ is {\it geometrically unibranch} if $X$ is
geometrically unibranch at every $x \in |X|$.
\end{definition}

\noindent
This is consistent with the definition for schemes
(Properties, Definition \ref{properties-definition-unibranch}).

\begin{lemma}
\label{lemma-nr-branches-local-ring}
Let $S$ be a scheme. Let $X$ be an algebraic space over $S$.
Let $x \in |X|$ be a point. Let $n \in \{1, 2, \ldots\}$ be an integer.
The following are equivalent
\begin{enumerate}
\item for any scheme $U$ and \'etale morphism $a : U \to X$ and
$u \in U$ with $a(u) = x$ the number of minimal primes
of the local ring $\mathcal{O}_{U, u}$ is $\leq n$
and for at least one choice of $U, a, u$ it is $n$,
\item for any scheme $U$ and \'etale morphism $a : U \to X$ and
$u \in U$ with $a(u) = x$ the number irreducible components of
$U$ passing through $u$ is $\leq n$ and for at least one choice
of $U, a, u$ it is $n$,
\item for any scheme $U$ and \'etale morphism $a : U \to X$ and $u \in U$
with $a(u) = x$ the number of branches of $U$ at $u$ is $\leq n$
and for at least one choice of $U, a, u$ it is $n$,
\item for any scheme $U$ and \'etale morphism $a : U \to X$ and $u \in U$
with $a(u) = x$ the number of geometric branches of $U$ at $u$ is $n$, and
\item the number of minimal prime ideals of
$\mathcal{O}_{X, \overline{x}}$ is $n$.
\end{enumerate}
\end{lemma}

\begin{proof}
The equivalence of (1) and (2) follows from the fact that irreducible
components of $U$ passing through $u$ are in $1$-$1$ correspondence with
minimal primes of the local ring of $U$ at $u$. Let $a : U \to X$ and
$u \in U$ be as in (1). Then $\mathcal{O}_{X, \overline{x}}$ is the
strict henselization of $\mathcal{O}_{U, u}$ by
Lemma \ref{lemma-describe-etale-local-ring}. Recall that
the (geometric) number of branches of $U$ at $u$ is the number
of minimal prime ideals of the (strict) henselization of $\mathcal{O}_{U, u}$.
In particular (4) and (5) are equivalent.
The equivalence of (2), (3), and (4) follows from
More on Morphisms, Lemma \ref{more-morphisms-lemma-nr-branches}.
\end{proof}

\begin{definition}
\label{definition-number-of-geometric-branches}
Let $S$ be a scheme.
Let $X$ be an algebraic space over $S$.
Let $x \in |X|$. The {\it number of geometric branches of $X$ at $x$} is
either $n \in \mathbf{N}$ if the equivalent conditions
of Lemma \ref{lemma-nr-branches-local-ring}
hold, or else $\infty$.
\end{definition}








\section{Noetherian spaces}
\label{section-noetherian}

\noindent
We have already defined locally Noetherian algebraic spaces in
Section \ref{section-types-properties}.

\begin{definition}
\label{definition-noetherian}
Let $S$ be a scheme. Let $X$ be an algebraic space over $S$.
We say $X$ is {\it Noetherian} if $X$ is quasi-compact, quasi-separated
and locally Noetherian.
\end{definition}

\noindent
Note that a Noetherian algebraic space $X$ is not just quasi-compact
and locally Noetherian, but also quasi-separated. This does not conflict
with the definition of a Noetherian scheme, as a locally Noetherian
scheme is quasi-separated, see
Properties, Lemma \ref{properties-lemma-locally-Noetherian-quasi-separated}.
This does not hold for algebraic spaces. Namely,
$X = \mathbf{A}^1_k/\mathbf{Z}$, see
Spaces, Example \ref{spaces-example-affine-line-translation}
is locally Noetherian and quasi-compact but not quasi-separated
(hence not Noetherian according to our definitions).

\medskip\noindent
A consequence of the choice made above is that an algebraic space
of finite type over a Noetherian algebraic space is not automatically
Noetherian, i.e., the analogue of
Morphisms, Lemma \ref{morphisms-lemma-finite-type-noetherian}
does not hold. The correct statement is that an algebraic space of
finite presentation over a Noetherian algebraic space is Noetherian
(see
Morphisms of Spaces,
Lemma \ref{spaces-morphisms-lemma-finite-presentation-noetherian}).

\medskip\noindent
A Noetherian algebraic space $X$ is very close to being a scheme.
In the rest of this section we collect some lemmas to illustrate this.

\begin{lemma}
\label{lemma-Noetherian-topology}
Let $S$ be a scheme. Let $X$ be an algebraic space over $S$.
\begin{enumerate}
\item If $X$ is locally Noetherian then $|X|$ is a locally Noetherian
topological space.
\item If $X$ is quasi-compact and locally Noetherian, then $|X|$
is a Noetherian topological space.
\end{enumerate}
\end{lemma}

\begin{proof}
Assume $X$ is locally Noetherian.
Choose a scheme $U$ and a surjective \'etale morphism
$U \to X$. As $X$ is locally Noetherian we see that $U$ is locally
Noetherian. By
Properties, Lemma \ref{properties-lemma-Noetherian-topology}
this means that $|U|$ is a locally Noetherian topological space.
Since $|U| \to |X|$ is open and surjective we conclude that
$|X|$ is locally Noetherian by
Topology, Lemma \ref{topology-lemma-image-Noetherian}.
This proves (1). If $X$ is quasi-compact and locally Noetherian,
then $|X|$ is quasi-compact and locally Noetherian. Hence $|X|$
is Noetherian by
Topology,
Lemma \ref{topology-lemma-quasi-compact-locally-Noetherian-Noetherian}.
\end{proof}

\begin{lemma}
\label{lemma-Noetherian-sober}
Let $S$ be a scheme. Let $X$ be an algebraic space over $S$.
If $X$ is Noetherian, then $|X|$ is a sober Noetherian topological space.
\end{lemma}

\begin{proof}
A quasi-separated algebraic space has an underlying sober topological
space, see
Lemma \ref{lemma-quasi-separated-sober}.
It is Noetherian by
Lemma \ref{lemma-Noetherian-topology}.
\end{proof}

\begin{lemma}
\label{lemma-Noetherian-local-ring-Noetherian}
Let $S$ be a scheme. Let $X$ be a Noetherian algebraic space over $S$.
Let $\overline{x}$ be a geometric point of $X$. Then
$\mathcal{O}_{X, \overline{x}}$ is a Noetherian local ring.
\end{lemma}

\begin{proof}
Choose an \'etale neighbourhood $(U, \overline{u})$ of $\overline{x}$
where $U$ is a scheme. Then $\mathcal{O}_{X, \overline{x}}$ is the
strict henselization of the local ring of $U$ at $u$, see
Lemma \ref{lemma-describe-etale-local-ring}.
By our definition of Noetherian spaces the scheme $U$ is locally Noetherian.
Hence we conclude by
More on Algebra, Lemma \ref{more-algebra-lemma-henselization-noetherian}.
\end{proof}








\section{Regular algebraic spaces}
\label{section-regular}

\noindent
We have already defined regular algebraic spaces in
Section \ref{section-types-properties}.

\begin{lemma}
\label{lemma-regular}
Let $S$ be a scheme.
Let $X$ be a locally Noetherian algebraic space over $S$.
The following are equivalent
\begin{enumerate}
\item $X$ is regular, and
\item every \'etale local ring $\mathcal{O}_{X, \overline{x}}$ is
regular.
\end{enumerate}
\end{lemma}

\begin{proof}
Let $U$ be a scheme and let $U \to X$ be a surjective \'etale morphism.
By assumption $U$ is locally Noetherian. Moreover, every \'etale local
ring $\mathcal{O}_{X, \overline{x}}$ is the strict henselization of
a local ring on $U$ and conversely, see
Lemma \ref{lemma-describe-etale-local-ring}.
Thus by
More on Algebra, Lemma \ref{more-algebra-lemma-henselization-regular}
we see that (2) is equivalent to every local ring of $U$ being
regular, i.e., $U$ being a regular scheme (see
Properties, Lemma \ref{properties-lemma-characterize-regular}).
This equivalent to (1) by
Definition \ref{definition-type-property}.
\end{proof}

\noindent
We can use Descent, Lemma \ref{descent-lemma-regular-local-ring-local}
to define what it means for an algebraic space $X$ to be regular at a
point $x$.

\begin{definition}
\label{definition-regular-at-point}
Let $S$ be a scheme. Let $X$ be an algebraic space over $S$.
Let $x \in |X|$ be a point. We say {\it $X$ is regular at $x$}
if $\mathcal{O}_{U, u}$ is a regular local ring for any
(equivalently some) pair $(a : U \to X, u)$ consisting of an
\'etale morphism $a : U \to X$ from a scheme to $X$ and a point
$u \in U$ with $a(u) = x$.
\end{definition}

\noindent
See Definition \ref{definition-property-at-point},
Lemma \ref{lemma-local-source-target-at-point}, and
Descent, Lemma \ref{descent-lemma-regular-local-ring-local}.

\begin{lemma}
\label{lemma-regular-at-x}
Let $S$ be a scheme. Let $X$ be an algebraic space over $S$.
Let $x \in |X|$ be a point. The following are equivalent
\begin{enumerate}
\item $X$ is regular at $x$, and
\item the \'etale local ring $\mathcal{O}_{X, \overline{x}}$ is
regular for any (equivalently some) geometric point $\overline{x}$
lying over $x$.
\end{enumerate}
\end{lemma}

\begin{proof}
Let $U$ be a scheme, $u \in U$ a point, and let $a : U \to X$ be an
\'etale morphism mapping $u$ to $x$. For any geometric point
$\overline{x}$ of $X$ lying over $x$, the \'etale local
ring $\mathcal{O}_{X, \overline{x}}$ is the strict henselization of
a local ring on $U$ at $u$, see
Lemma \ref{lemma-describe-etale-local-ring}.
Thus we conclude by
More on Algebra, Lemma \ref{more-algebra-lemma-henselization-regular}.
\end{proof}

\begin{lemma}
\label{lemma-regular-normal}
A regular algebraic space is normal.
\end{lemma}

\begin{proof}
This follows from the definitions and the case of schemes
See Properties, Lemma \ref{properties-lemma-regular-normal}.
\end{proof}

 








\section{Sheaves of modules on algebraic spaces}
\label{section-modules}

\noindent
If $X$ is an algebraic space, then a sheaf of modules on $X$ is
a sheaf of $\mathcal{O}_X$-modules on the small \'etale site of $X$
where $\mathcal{O}_X$ is the structure sheaf of $X$. The category
of sheaves of modules is denoted $\textit{Mod}(\mathcal{O}_X)$.

\medskip\noindent
Given a morphism $f : X \to Y$ of algebraic spaces, by
Lemma \ref{lemma-morphism-ringed-topoi}
we get a morphism of ringed topoi and hence by
Modules on Sites, Definition \ref{sites-modules-definition-pushforward}
we get well defined pullback and direct image functors
\begin{equation}
\label{equation-push-pull}
f^* :
\textit{Mod}(\mathcal{O}_Y)
\longrightarrow
\textit{Mod}(\mathcal{O}_X), \quad
f_* :
\textit{Mod}(\mathcal{O}_X)
\longrightarrow
\textit{Mod}(\mathcal{O}_Y)
\end{equation}
which are adjoint in the usual way. If $g : Y \to Z$ is another morphism
of algebraic spaces over $S$, then we have
$(g \circ f)^* = f^* \circ g^*$ and $(g \circ f)_* = g_* \circ f_*$
simply because the morphisms of ringed topoi compose in the corresponding
way (by the lemma).

\begin{lemma}
\label{lemma-etale-exact-pullback}
Let $S$ be a scheme.
Let $f : X \to Y$ be an \'etale morphism of algebraic spaces over $S$.
Then $f^{-1}\mathcal{O}_Y = \mathcal{O}_X$, and
$f^*\mathcal{G} = f_{small}^{-1}\mathcal{G}$ for any sheaf of
$\mathcal{O}_Y$-modules $\mathcal{G}$. In particular,
$f^* : \textit{Mod}(\mathcal{O}_Y) \to \textit{Mod}(\mathcal{O}_X)$
is exact.
\end{lemma}

\begin{proof}
By the description of inverse image in Lemma \ref{lemma-etale-morphism-topoi}
and the definition of the structure sheaves it is clear that
$f_{small}^{-1}\mathcal{O}_Y = \mathcal{O}_X$. Since the pullback
$$
f^*\mathcal{G} =
f_{small}^{-1}\mathcal{G} \otimes_{f_{small}^{-1}\mathcal{O}_Y}
\mathcal{O}_X
$$
by definition we conclude that $f^*\mathcal{G} = f_{small}^{-1}\mathcal{G}$.
The exactness is clear because $f_{small}^{-1}$ is exact, as $f_{small}$
is a morphism of topoi.
\end{proof}

\noindent
We continue our abuse of notation introduced in
Equation (\ref{equation-restrict})
by writing
\begin{equation}
\label{equation-restrict-modules}
\mathcal{G}|_{X_\etale}
= f^*\mathcal{G}
= f_{small}^{-1}\mathcal{G}
\end{equation}
in the situation of the lemma above. We will discuss this in a more
technical fashion in
Section \ref{section-localize}.

\begin{lemma}
\label{lemma-pushforward-etale-base-change-modules}
Let $S$ be a scheme. Let
$$
\xymatrix{
X' \ar[r] \ar[d]_{f'} & X \ar[d]^f \\
Y' \ar[r]^g & Y
}
$$
be a cartesian square of algebraic spaces over $S$. Let
$\mathcal{F} \in \textit{Mod}(\mathcal{O}_X)$. If $g$ is \'etale, then
$f'_*(\mathcal{F}|_{X'}) = (f_*\mathcal{F})|_{Y'}$\footnote{Also
$(f')^*(\mathcal{G}|_{Y'}) = (f^*\mathcal{G})|_{X'}$
by commutativity of the diagram and (\ref{equation-restrict-modules})} and
$R^if'_*(\mathcal{F}|_{X'}) = (R^if_*\mathcal{F})|_{Y'}$ in
$\textit{Mod}(\mathcal{O}_{Y'})$.
\end{lemma}

\begin{proof}
This is a reformulation of
Lemma \ref{lemma-pushforward-etale-base-change}
in the case of modules.
\end{proof}

\begin{lemma}
\label{lemma-characterize-module-small-etale}
Let $S$ be a scheme. Let $X$ be an algebraic space over $S$.
A sheaf $\mathcal{F}$ of $\mathcal{O}_X$-modules is given by the following
data:
\begin{enumerate}
\item for every $U \in \Ob(X_\etale)$ a sheaf
$\mathcal{F}_U$ of $\mathcal{O}_U$-modules on $U_\etale$,
\item for every $f : U' \to U$ in $X_\etale$ an isomorphism
$c_f : f_{small}^*\mathcal{F}_U \to \mathcal{F}_{U'}$.
\end{enumerate}
These data are subject to the condition that given any $f : U' \to U$
and $g : U'' \to U'$ in $X_\etale$ the composition
$c_g \circ g_{small}^*c_f$ is equal to $c_{f \circ g}$.
\end{lemma}

\begin{proof}
Combine Lemmas \ref{lemma-etale-exact-pullback}
and \ref{lemma-characterize-sheaf-small-etale}, and use the fact that
any morphism between objects of $X_\etale$ is an \'etale morphism
of schemes.
\end{proof}







\section{\'Etale localization}
\label{section-localize}

\noindent
Reading this section should be avoided at all cost.

\medskip\noindent
Let $X \to Y$ be an \'etale morphism of algebraic spaces.
Then $X$ is an object of $Y_{spaces, \etale}$ and it is
immediate from the definitions, see also the proof of
Lemma \ref{lemma-etale-morphism-topoi},
that
\begin{equation}
\label{equation-localize}
X_{spaces, \etale} = Y_{spaces, \etale}/X
\end{equation}
where the right hand side is the localization of the site
$Y_{spaces, \etale}$ at the object $X$, see
Sites, Definition \ref{sites-definition-localize}.
Moreover, this identification is compatible with the structure sheaves by
Lemma \ref{lemma-etale-exact-pullback}.
Hence the ringed site $(X_{spaces, \etale}, \mathcal{O}_X)$
is identified with the localization of the ringed site
$(Y_{spaces, \etale}, \mathcal{O}_Y)$ at the object $X$:
\begin{equation}
\label{equation-localize-ringed}
(X_{spaces, \etale}, \mathcal{O}_X) =
(Y_{spaces, \etale}/X, \mathcal{O}_Y|_{Y_{spaces, \etale}/X})
\end{equation}
The localization of a ringed site used on the right hand side is defined in
Modules on Sites,
Definition \ref{sites-modules-definition-localize-ringed-site}.

\medskip\noindent
Assume now $X \to Y$ is an \'etale morphism of algebraic spaces and $X$ is
a scheme. Then $X$ is an object of $Y_\etale$ and it follows that
\begin{equation}
\label{equation-localize-at-scheme}
X_\etale = Y_\etale/X
\end{equation}
and
\begin{equation}
\label{equation-localize-at-scheme-ringed}
(X_\etale, \mathcal{O}_X) =
(Y_\etale/X, \mathcal{O}_Y|_{Y_\etale/X})
\end{equation}
as above.

\medskip\noindent
Finally, if $X \to Y$ is an \'etale morphism of algebraic spaces and $X$ is
an affine scheme, then $X$ is an object of $Y_{affine, \etale}$ and
\begin{equation}
\label{equation-localize-at-affine}
X_{affine, \etale} = Y_{affine, \etale}/X
\end{equation}
and
\begin{equation}
\label{equation-localize-at-affine-ringed}
(X_{affine, \etale}, \mathcal{O}_X) =
(Y_{affine, \etale}/X, \mathcal{O}_Y|_{Y_{affine, \etale}/X})
\end{equation}
as above.

\medskip\noindent
Next, we show that these localizations are compatible with morphisms.

\begin{lemma}
\label{lemma-relocalize-morphism}
Let $S$ be a scheme. Let
$$
\xymatrix{
U \ar[d]_p \ar[r]_g & V \ar[d]^q \\
X \ar[r]^f & Y
}
$$
be a commutative diagram of algebraic spaces over $S$ with $p$ and $q$ \'etale.
Via the identifications
(\ref{equation-localize-ringed}) for $U \to X$ and $V \to Y$
the morphism of ringed topoi
$$
(g_{spaces, \etale}, g^\sharp) :
(\Sh(U_{spaces, \etale}), \mathcal{O}_U)
\longrightarrow
(\Sh(V_{spaces, \etale}), \mathcal{O}_V)
$$
is $2$-isomorphic to the morphism $(f_{spaces, \etale, c}, f_c^\sharp)$
constructed in
Modules on Sites,
Lemma \ref{sites-modules-lemma-relocalize-morphism-ringed-sites}
starting with the morphism of ringed sites
$(f_{spaces, \etale}, f^\sharp)$ and
the map $c : U \to V \times_Y X$ corresponding to $g$.
\end{lemma}

\begin{proof}
The morphism $(f_{spaces, \etale, c}, f_c^\sharp)$ is defined as a
composition $f' \circ j$
of a localization and a base change map. Similarly $g$ is a composition
$U \to V \times_Y X \to V$. Hence it suffices to prove
the lemma in the following two cases: (1) $f = \text{id}$, and
(2) $U = X \times_Y V$. In case (1) the morphism $g : U \to V$ is
\'etale, see
Lemma \ref{lemma-etale-permanence}.
Hence $(g_{spaces, \etale}, g^\sharp)$ is a localization morphism
by the discussion surrounding
Equations (\ref{equation-localize}) and
(\ref{equation-localize-ringed})
which is exactly the content of the lemma in this case.
In case (2) the morphism $g_{spaces, \etale}$
comes from the morphism of ringed sites given by the functor
$V_{spaces, \etale} \to U_{spaces, \etale}$,
$V'/V \mapsto V' \times_V U/U$
which is also what the morphism $f'$ is defined by, see
Sites, Lemma \ref{sites-lemma-localize-morphism}.
We omit the verification that $(f')^\sharp = g^\sharp$
in this case (both are the restriction of $f^\sharp$
to $U_{spaces, \etale}$).
\end{proof}

\begin{lemma}
\label{lemma-relocalize-morphism-at-schemes}
Same notation and assumptions as in
Lemma \ref{lemma-relocalize-morphism}
except that we also assume $U$ and $V$ are schemes.
Via the identifications
(\ref{equation-localize-at-scheme-ringed})
for $U \to X$ and $V \to Y$ the morphism of ringed topoi
$$
(g_{small}, g^\sharp) :
(\Sh(U_\etale), \mathcal{O}_U)
\longrightarrow
(\Sh(V_\etale), \mathcal{O}_V)
$$
is $2$-isomorphic to the morphism $(f_{small, s}, f_s^\sharp)$
constructed in
Modules on Sites,
Lemma \ref{sites-modules-lemma-relocalize-morphism-ringed-topoi}
starting with $(f_{small}, f^\sharp)$ and
the map $s : h_U \to f_{small}^{-1}h_V$ corresponding to $g$.
\end{lemma}

\begin{proof}
Note that $(g_{small}, g^\sharp)$ is $2$-isomorphic as a
morphism of ringed topoi to the morphism of ringed topoi
associated to the morphism of ringed sites
$(g_{spaces, \etale}, g^\sharp)$. Hence we conclude by
Lemma \ref{lemma-relocalize-morphism}
and
Modules on Sites,
Lemma \ref{sites-modules-lemma-relocalize-morphism-compare}.
\end{proof}

\noindent
Finally, we discuss the relationship between sheaves of sets on the small
\'etale site $Y_\etale$ of an algebraic space $Y$ and algebraic spaces
\'etale over $Y$. Let $S$ be a scheme and let $Y$ be an algebraic space
over $S$. Let $\mathcal{F}$ be an object of $\Sh(Y_\etale)$. Consider
the functor
$$
X : (\Sch/S)_{fppf}^{opp} \longrightarrow \textit{Sets}
$$
defined by the rule
$$
X(T) = \{(y, s) \mid y : T \to Y\text{ is a morphism over }S\text{ and }
s \in \Gamma(T, y_{small}^{-1}\mathcal{F})\}
$$
Given a morphism $g : T' \to T$ the restriction map sends
$(y, s)$ to $(y \circ g, g_{small}^{-1}s)$. This makes sense
as $y_{small} \circ g_{small} = (y \circ g)_{small}$ by
Lemma \ref{lemma-functoriality-etale-site}. There is a canonical map
$X \to Y$ sending the pair $(y, s)$ to $y$.

\begin{lemma}
\label{lemma-sheaf-gives-space}
Let $S$ be a scheme and let $Y$ be an algebraic space over $S$.
Let $\mathcal{F}$ be a sheaf of sets on $Y_\etale$.
Provided a set theoretic condition is satisfied (see proof) we have
\begin{enumerate}
\item the functor $X$ associated to $\mathcal{F}$ above is an
algebraic space,
\item the map $X \to Y$ is an \'etale morphism of algebraic spaces,
\item via the identification $\Sh(Y_\etale) = \Sh(Y_{spaces, \etale})$
we have $\mathcal{F} \cong h_X$,
\item we have $\mathcal{F} \cong f_{small, !}*$. Here $*$ is the final object
of the category $\Sh(X_\etale)$ and $f_{small, !}$ exists by
Lemma \ref{lemma-etale-morphism-topoi}.
\end{enumerate}
\end{lemma}

\begin{proof}
Let us prove that $X$ is a sheaf for the fppf topology. Namely, suppose
that $\{g_i : T_i \to T\}$ is a covering of $(\Sch/S)_{fppf}$ and
$(y_i, s_i) \in X(T_i)$ satisfy the glueing condition, i.e., the
restriction of $(y_i, s_i)$ and $(y_j, s_j)$ to $T_i \times_T T_j$ agree.
Then since $Y$ is a sheaf for the fppf topology, we see that
the $y_i$ give rise to a unique morphism $y : T \to Y$ such that
$y_i = y \circ g_i$. Then we see that
$y_{i, small}^{-1}\mathcal{F} = g_{i, small}^{-1}y_{small}^{-1}\mathcal{F}$.
Hence the sections $s_i$ glue uniquely to a section of
$y_{small}^{-1}\mathcal{F}$ by
\'Etale Cohomology, Lemma \ref{etale-cohomology-lemma-describe-pullback}.

\medskip\noindent
The construction that sends $\mathcal{F} \in \Ob(\Sh(Y_\etale))$
to $X \in \Ob((\Sch/S)_{fppf})$ preserves finite limits and
all colimits since each of the functors $y_{small}^{-1}$ have
this property. Of course, if $V \in \Ob(Y_\etale)$, then
the construction sends the representable sheaf $h_V$ on $Y_\etale$
to the representable functor represented by $V$.

\medskip\noindent
By Sites, Lemma \ref{sites-lemma-sheaf-coequalizer-representable}
we can find a set $I$, for each $i \in I$ an object $V_i$ of $Y_\etale$
and a surjective map of sheaves
$$
\coprod h_{V_i} \longrightarrow \mathcal{F}
$$
on $Y_\etale$. The set theoretic condition we need is that the index
set $I$ is not too large\footnote{It suffices if the supremum of the
cardinalities of the stalks of $\mathcal{F}$ at geometric points of $Y$
is bounded by the size of some object of $(\Sch/S)_{fppf}$.}.
Then $V = \coprod V_i$ is an object of $(\Sch/S)_{fppf}$ and
therefore an object of $Y_\etale$ and we have a surjective map
$h_V \to \mathcal{F}$.

\medskip\noindent
Observe that the product of $h_V$ with itself in $\Sh(Y_\etale)$
is $h_{V \times_Y V}$. Consider the fibre product
$$
h_V \times_\mathcal{F} h_V \subset h_{V \times_Y V}
$$
There is an open subscheme $R$ of $V \times_Y V$ such that
$h_V \times_\mathcal{F} h_V = h_R$, see
Lemma \ref{lemma-support-subsheaf-final} (small detail omitted).
By the Yoneda lemma we obtain two morphisms $s, t : R \to V$ in $Y_\etale$
and we find a coequalizer diagram
$$
\xymatrix{
h_R \ar@<1ex>[r] \ar@<-1ex>[r] &
h_V \ar[r] &
\mathcal{F}
}
$$
in $\Sh(Y_\etale)$. Of course the morphisms $s, t$ are \'etale
and define an \'etale equivalence relation $(t, s) : R \to V \times_S V$.

\medskip\noindent
By the discussion in the preceding two paragraphs we find a
coequalizer diagram
$$
\xymatrix{
R \ar@<1ex>[r] \ar@<-1ex>[r] &
V \ar[r] &
X
}
$$
in $(\Sch/S)_{fppf}$. Thus $X = V/R$ is an algebraic space by
Spaces, Theorem \ref{spaces-theorem-presentation}. This proves (1).
Part (2) follows because $V \to Y$ is \'etale. Part (3) is immediate
from the definition of $X$ and $h_X$. We omit the proof of part (4);
it follows by matching the morphism associated to the cocontinuous
functor $j$ of Lemma \ref{lemma-etale-morphism-topoi}
with the description of $X_{spaces, \etale}$ as the localization
of $Y_{spaces, \etale}$ at $X$ discussed above and
Sites, Lemma \ref{sites-lemma-describe-j-shriek-representable}.
\end{proof}







\section{Recovering morphisms}
\label{section-morphisms}

\noindent
In this section we prove that the rule which associates to an algebraic space
its locally ringed small \'etale topos is fully faithful in a suitable
sense, see
Theorem \ref{theorem-fully-faithful}.

\begin{lemma}
\label{lemma-morphism-locally-ringed}
Let $S$ be a scheme.
Let $f : X \to Y$ be a morphism of algebraic spaces over $S$.
The morphism of ringed topoi $(f_{small}, f^\sharp)$
associated to $f$ is a morphism of locally ringed topoi, see
Modules on Sites,
Definition \ref{sites-modules-definition-morphism-locally-ringed-topoi}.
\end{lemma}

\begin{proof}
Note that the assertion makes sense since we have seen that
$(X_\etale, \mathcal{O}_{X_\etale})$ and
$(Y_\etale, \mathcal{O}_{Y_\etale})$
are locally ringed sites, see
Lemma \ref{lemma-etale-site-locally-ringed}.
Moreover, we know that $X_\etale$ has enough points, see
Theorem \ref{theorem-exactness-stalks}.
Hence it suffices to prove that $(f_{small}, f^\sharp)$
satisfies condition (3) of
Modules on Sites,
Lemma \ref{sites-modules-lemma-locally-ringed-morphism}.
To see this take a point $p$ of $X_\etale$. By
Lemma \ref{lemma-points-small-etale-site}
$p$ corresponds to a geometric point $\overline{x}$ of $X$.
By
Lemma \ref{lemma-stalk-pullback}
the point $q = f_{small} \circ p$ corresponds to the
geometric point $\overline{y} = f \circ \overline{x}$ of $Y$.
Hence the assertion we have to prove is that the induced map
of \'etale local rings
$$
\mathcal{O}_{Y, \overline{y}} \longrightarrow \mathcal{O}_{X, \overline{x}}
$$
is a local ring map. You can prove this directly, but instead we deduce it
from the corresponding result for schemes. To do this choose a commutative
diagram
$$
\xymatrix{
U \ar[d] \ar[r]_\psi & V \ar[d] \\
X \ar[r] & Y
}
$$
where $U$ and $V$ are schemes, and the vertical arrows are surjective
\'etale (see
Spaces, Lemma \ref{spaces-lemma-lift-morphism-presentations}).
Choose a lift $\overline{u} : \overline{x} \to U$ (possible by
Lemma \ref{lemma-geometric-lift-to-cover}).
Set $\overline{v} = \psi \circ \overline{u}$. We obtain a commutative
diagram of \'etale local rings
$$
\xymatrix{
\mathcal{O}_{U, \overline{u}} &
\mathcal{O}_{V, \overline{v}} \ar[l] \\
\mathcal{O}_{X, \overline{x}} \ar[u] &
\mathcal{O}_{Y, \overline{y}}. \ar[l] \ar[u]
}
$$
By
\'Etale Cohomology, Lemma \ref{etale-cohomology-lemma-morphism-locally-ringed}
the top horizontal arrow is a local ring map. Finally by
Lemma \ref{lemma-describe-etale-local-ring}
the vertical arrows are isomorphisms. Hence we win.
\end{proof}

\begin{lemma}
\label{lemma-2-morphism}
Let $S$ be a scheme.
Let $X$, $Y$ be algebraic spaces over $S$.
Let $f : X \to Y$ be a morphism of algebraic spaces over $S$.
Let $t$ be a $2$-morphism from $(f_{small}, f^\sharp)$ to itself, see
Modules on Sites,
Definition \ref{sites-modules-definition-2-morphism-ringed-topoi}.
Then $t = \text{id}$.
\end{lemma}

\begin{proof}
Let $X'$, resp.\ $Y'$ be $X$ viewed as an algebraic space over
$\Spec(\mathbf{Z})$, see
Spaces, Definition \ref{spaces-definition-base-change}.
It is clear from the construction that $(X_{small}, \mathcal{O})$
is equal to $(X'_{small}, \mathcal{O})$ and similarly for $Y$.
Hence we may work with $X'$ and $Y'$. In other words we may
assume that $S = \Spec(\mathbf{Z})$.

\medskip\noindent
Assume $S = \Spec(\mathbf{Z})$, $f : X \to Y$ and $t$ are as in
the lemma. This means that $t : f^{-1}_{small} \to f^{-1}_{small}$
is a transformation of functors such that the diagram
$$
\xymatrix{
f_{small}^{-1}\mathcal{O}_Y
\ar[rd]_{f^\sharp}  & &
f_{small}^{-1}\mathcal{O}_Y \ar[ll]^t \ar[ld]^{f^\sharp} \\
& \mathcal{O}_X
}
$$
is commutative. Suppose $V \to Y$ is \'etale with $V$ affine.
Write $V = \Spec(B)$. Choose generators $b_j \in B$, $j \in J$
for $B$ as a $\mathbf{Z}$-algebra. Set
$T = \Spec(\mathbf{Z}[\{x_j\}_{j \in J}])$.
In the following we will use that
$\Mor_{\Sch}(U, T) = \prod_{j \in J} \Gamma(U, \mathcal{O}_U)$
for any scheme $U$ without further mention.
The surjective ring map $\mathbf{Z}[x_j] \to B$, $x_j \mapsto b_j$
corresponds to a closed immersion $V \to T$.
We obtain a monomorphism
$$
i : V \longrightarrow T_Y = T \times Y
$$
of algebraic spaces over $Y$. In terms of sheaves on $Y_\etale$
the morphism $i$ induces an injection
$h_i : h_V \to \prod_{j \in J} \mathcal{O}_Y$ of sheaves.
The base change $i' : X \times_Y V \to T_X$ of $i$ to $X$
is a monomorphism too
(Spaces,
Lemma \ref{spaces-lemma-base-change-representable-transformations-property}).
Hence $i' : X \times_Y V \to T_X$ is a monomorphism, which
in turn means that
$h_{i'} : h_{X \times_Y V} \to \prod_{j \in J} \mathcal{O}_X$
is an injection of sheaves.
Via the identification $f_{small}^{-1}h_V = h_{X \times_Y V}$ of
Lemma \ref{lemma-stalk-pullback}
the map $h_{i'}$ is equal to
$$
\xymatrix{
f_{small}^{-1}h_V \ar[r]^-{f^{-1}h_i} &
\prod_{j \in J} f_{small}^{-1}\mathcal{O}_Y
\ar[r]^{\prod f^\sharp} &
\prod_{j \in J} \mathcal{O}_X
}
$$
(verification omitted). This means that the map
$t : f_{small}^{-1}h_V \to f_{small}^{-1}h_V$
fits into the commutative diagram
$$
\xymatrix{
f_{small}^{-1}h_V \ar[r]^-{f^{-1}h_i} \ar[d]^t &
\prod_{j \in J} f_{small}^{-1}\mathcal{O}_Y
\ar[r]^-{\prod f^\sharp} \ar[d]^{\prod t} &
\prod_{j \in J} \mathcal{O}_X \ar[d]^{\text{id}}\\
f_{small}^{-1}h_V \ar[r]^-{f^{-1}h_i} &
\prod_{j \in J} f_{small}^{-1}\mathcal{O}_Y
\ar[r]^-{\prod f^\sharp} &
\prod_{j \in J} \mathcal{O}_X
}
$$
The commutativity of the right square holds by our assumption on $t$
explained above.
Since the composition of the horizontal arrows is injective
by the discussion above we conclude that the left vertical arrow
is the identity map as well. Any sheaf of sets on
$Y_\etale$ admits a surjection from a (huge) coproduct of sheaves
of the form $h_V$ with $V$ affine (combine
Lemma \ref{lemma-alternative}
with
Sites, Lemma \ref{sites-lemma-sheaf-coequalizer-representable}).
Thus we conclude that $t : f_{small}^{-1} \to f_{small}^{-1}$
is the identity transformation as desired.
\end{proof}

\begin{lemma}
\label{lemma-faithful}
Let $S$ be a scheme.
Let $X$, $Y$ be algebraic spaces over $S$.
Any two morphisms $a, b : X \to Y$ of algebraic spaces over $S$
for which there exists a $2$-isomorphism
$(a_{small}, a^\sharp) \cong (b_{small}, b^\sharp)$
in the $2$-category of ringed topoi are equal.
\end{lemma}

\begin{proof}
Let $t : a_{small}^{-1} \to b_{small}^{-1}$ be the $2$-isomorphism.
We may equivalently think of $t$ as a transformation
$t : a_{spaces, \etale}^{-1} \to b_{spaces, \etale}^{-1}$
since there is not difference between sheaves on $X_\etale$
and sheaves on $X_{spaces, \etale}$.
Choose a commutative diagram
$$
\xymatrix{
U \ar[d]_p \ar[r]_\alpha  & V \ar[d]^q \\
X \ar[r]^a & Y
}
$$
where $U$ and $V$ are schemes, and $p$ and $q$ are surjective \'etale.
Consider the diagram
$$
\xymatrix{
h_U \ar[r]_-\alpha \ar@{=}[d] & a_{spaces, \etale}^{-1}h_V \ar[d]^t \\
h_U \ar@{..>}[r] & b_{spaces, \etale}^{-1}h_V
}
$$
Since the sheaf $b_{spaces, \etale}^{-1}h_V$ is isomorphic to
$h_{V \times_{Y, b} X}$ we see that the dotted arrow comes from a
morphism of schemes
$\beta : U \to V$ fitting into a commutative diagram
$$
\xymatrix{
U \ar[d]_p \ar[r]_\beta  & V \ar[d]^q \\
X \ar[r]^b & Y
}
$$
We claim that there exists a sequence of $2$-isomorphisms
\begin{align*}
(\alpha_{small}, \alpha^\sharp)
& \cong
(\alpha_{spaces, \etale}, \alpha^\sharp) \\
& \cong
(a_{spaces, \etale, c}, a_c^\sharp) \\
& \cong
(b_{spaces, \etale, d}, b_d^\sharp) \\
& \cong
(\beta_{spaces, \etale}, \beta^\sharp) \\
& \cong
(\beta_{small}, \beta^\sharp)
\end{align*}
The first and the last $2$-isomorphisms come from the identifications
between sheaves on $U_{spaces, \etale}$ and sheaves on
$U_\etale$ and similarly for $V$. The second and fourth
$2$-isomorphisms are those of
Lemma \ref{lemma-relocalize-morphism}
with $c : U \to X \times_{a, Y} V$ induced by $\alpha$ and
$d : U \to X \times_{b, Y} V$ induced by $\beta$.
The middle $2$-isomorphism comes from the transformation $t$.
Namely, the functor $a_{spaces, \etale, c}^{-1}$ corresponds
to the functor
$$
(\mathcal{H} \to h_V) \longmapsto
(a_{spaces, \etale}^{-1}\mathcal{H}
\times_{a_{spaces, \etale}^{-1}h_V, \alpha}
h_U \to
h_U)
$$
and similarly for $b_{spaces, \etale, d}^{-1}$, see
Sites, Lemma \ref{sites-lemma-relocalize-morphism}.
This uses the identification of sheaves on $Y_{spaces, \etale}/V$
as arrows $(\mathcal{H} \to h_V)$ in $\Sh(Y_{spaces, \etale})$
and similarly for $U/X$, see
Sites, Lemma \ref{sites-lemma-essential-image-j-shriek}.
Via this identification the structure sheaf $\mathcal{O}_V$ corresponds to the
pair $(\mathcal{O}_Y \times h_V \to h_V)$ and similarly
for $\mathcal{O}_U$, see
Modules on Sites,
Lemma \ref{sites-modules-lemma-localize-compare}.
Since $t$ switches $\alpha$ and $\beta$ we see that $t$ induces an isomorphism
$$
t :
a_{spaces, \etale}^{-1}\mathcal{H}
\times_{a_{spaces, \etale}^{-1}h_V, \alpha}
h_U
\longrightarrow
b_{spaces, \etale}^{-1}\mathcal{H}
\times_{b_{spaces, \etale}^{-1}h_V, \beta}
h_U
$$
over $h_U$ functorially in $(\mathcal{H} \to h_V)$. Also, $t$ is compatible
with $a_c^\sharp$ and $b_d^\sharp$ as $t$ is
compatible with $a^\sharp$ and $b^\sharp$ by our description
of the structure sheaves $\mathcal{O}_U$ and $\mathcal{O}_V$
above. Hence, the morphisms of ringed topoi
$(\alpha_{small}, \alpha^\sharp)$ and $(\beta_{small}, \beta^\sharp)$
are $2$-isomorphic. By
\'Etale Cohomology, Lemma \ref{etale-cohomology-lemma-faithful}
we conclude $\alpha = \beta$! Since $p : U \to X$ is a surjection
of sheaves it follows that $a = b$.
\end{proof}

\noindent
Here is the main result of this section.

\begin{theorem}
\label{theorem-fully-faithful}
Let $X$, $Y$ be algebraic spaces over $\Spec(\mathbf{Z})$.
Let
$$
(g, g^\sharp) :
(\Sh(X_\etale), \mathcal{O}_X)
\longrightarrow
(\Sh(Y_\etale), \mathcal{O}_Y)
$$
be a morphism of locally ringed topoi. Then there exists a
unique morphism of algebraic spaces $f : X \to Y$ such that
$(g, g^\sharp)$ is isomorphic to $(f_{small}, f^\sharp)$.
In other words, the construction
$$
\textit{Spaces}/\Spec(\mathbf{Z})
\longrightarrow \textit{Locally ringed topoi},
\quad
X \longrightarrow (X_\etale, \mathcal{O}_X)
$$
is fully faithful (morphisms up to $2$-isomorphisms on the right hand side).
\end{theorem}

\begin{proof}
The uniqueness we have seen in
Lemma \ref{lemma-faithful}.
Thus it suffices to prove existence.
In this proof we will freely use the identifications of
Equation (\ref{equation-localize-at-scheme-ringed})
as well as the result of
Lemma \ref{lemma-relocalize-morphism-at-schemes}.

\medskip\noindent
Let $U \in \Ob(X_\etale)$, let
$V \in \Ob(Y_\etale)$
and let $s \in g^{-1}h_V(U)$ be a section. We may think of
$s$ as a map of sheaves $s : h_U \to g^{-1}h_V$. By
Modules on Sites,
Lemma \ref{sites-modules-lemma-relocalize-morphism-ringed-topoi}
we obtain a commutative diagram of morphisms of ringed topoi
$$
\xymatrix{
(\Sh(X_\etale/U), \mathcal{O}_U)
\ar[rr]_-{(j, j^\sharp)} \ar[d]_{(g_s, g_s^\sharp)} & &
(\Sh(X_\etale), \mathcal{O}_X) \ar[d]^{(g, g^\sharp)} \\
(\Sh(V_\etale), \mathcal{O}_V) \ar[rr] & &
(\Sh(Y_\etale), \mathcal{O}_Y).
}
$$
By
\'Etale Cohomology, Theorem \ref{etale-cohomology-theorem-fully-faithful}
we obtain a unique morphism of schemes $f_s : U \to V$ such that
$(g_s, g_s^\sharp)$ is $2$-isomorphic to $(f_{s, small}, f_s^\sharp)$.
The construction $(U, V, s) \leadsto f_s$ just explained satisfies
the following functoriality property: Suppose given morphisms
$a : U' \to U$ in $X_\etale$ and $b : V' \to V$ in $Y_\etale$
and a map $s' : h_{U'} \to g^{-1}h_{V'}$ such that the diagram
$$
\xymatrix{
h_{U'} \ar[d]_a \ar[r]_{s'} & g^{-1}h_{V'} \ar[d]^{g^{-1}b} \\
h_U \ar[r]^s & g^{-1}h_V
}
$$
commutes. Then the diagram
$$
\xymatrix{
U' \ar[r]_-{f_{s'}} \ar[d]_a & u(V') \ar[d]^{u(b)} \\
U \ar[r]^-{f_s} & u(V)
}
$$
of schemes commutes. The reason this is true is that the same condition
holds for the morphisms $(g_s, g_s^\sharp)$ constructed in
Modules on Sites,
Lemma \ref{sites-modules-lemma-relocalize-morphism-ringed-topoi}
and the uniqueness in
\'Etale Cohomology, Theorem \ref{etale-cohomology-theorem-fully-faithful}.

\medskip\noindent
The problem is to glue the morphisms $f_s$ to a morphism of algebraic
spaces. To do this first choose a scheme $V$ and a surjective \'etale
morphism $V \to Y$. This means that $h_V \to *$ is surjective and hence
$g^{-1}h_V \to *$ is surjective too. This means there exists a scheme $U$
and a surjective \'etale morphism $U \to X$ and a morphism
$s : h_U \to g^{-1}h_V$. Next, set $R = V \times_Y V$ and
$R' = U \times_X U$. Then we get
$g^{-1}h_R = g^{-1}h_V \times g^{-1}h_V$ as $g^{-1}$ is exact.
Thus $s$ induces a morphism $s \times s : h_{R'} \to g^{-1}h_R$.
Applying the constructions above we see that we get a
commutative diagram of morphisms of schemes
$$
\xymatrix{
R' \ar@<1ex>[d] \ar@<-1ex>[d] \ar[rr]_{f_{s \times s}} & &
R \ar@<1ex>[d] \ar@<-1ex>[d] \\
U \ar[rr]^{f_s} & &
V
}
$$
Since we have $X = U/R'$ and $Y = V/R$ (see
Spaces, Lemma \ref{spaces-lemma-space-presentation})
we conclude that this diagram
defines a morphism of algebraic spaces $f : X \to Y$ fitting
into an obvious commutative diagram.
Now we still have to show that $(f_{small}, f^\sharp)$ is
$2$-isomorphic to $(g, g^\sharp)$.
Let $t_V : f_{s, small}^{-1} \to g_s^{-1}$ and
$t_R : f_{s \times s, small}^{-1} \to g_{s \times s}^{-1}$ be
the $2$-isomorphisms which are given to us by the construction above.
Let $\mathcal{G}$ be a sheaf on $Y_\etale$. Then we see that
$t_V$ defines an isomorphism
$$
f_{small}^{-1}\mathcal{G}|_{U_\etale}
=
f_{s, small}^{-1}\mathcal{G}|_{V_\etale}
\xrightarrow{t_V}
g_s^{-1}\mathcal{G}|_{V_\etale}
=
g^{-1}\mathcal{G}|_{U_\etale}.
$$
Moreover, this isomorphism pulled back to $R'$ via either projection
$R' \to U$ is the isomorphism
$$
f_{small}^{-1}\mathcal{G}|_{R'_\etale}
=
f_{s \times s, small}^{-1}\mathcal{G}|_{R_\etale}
\xrightarrow{t_R}
g_{s \times s}^{-1}\mathcal{G}|_{R_\etale}
=
g^{-1}\mathcal{G}|_{R'_\etale}.
$$
Since $\{U \to X\}$ is a covering in the site $X_{spaces, \etale}$
this means the first displayed isomorphism descends to an isomorphism
$t : f_{small}^{-1}\mathcal{G} \to g^{-1}\mathcal{G}$
of sheaves (small detail omitted). The isomorphism is functorial
in $\mathcal{G}$ since $t_V$ and $t_R$ are transformations of functors.
Finally, $t$ is compatible with $f^\sharp$ and $g^\sharp$ as
$t_V$ and $t_R$ are (some details omitted).
This finishes the proof of the theorem.
\end{proof}

\begin{lemma}
\label{lemma-isomorphism-ringed-topoi}
Let $X$, $Y$ be algebraic spaces over $\mathbf{Z}$. If
$$
(g, g^\sharp) :
(\Sh(X_\etale), \mathcal{O}_X)
\longrightarrow
(\Sh(Y_\etale), \mathcal{O}_Y)
$$
is an isomorphism of ringed topoi, then there exists a unique
morphism $f : X \to Y$ of algebraic spaces such that
$(g, g^\sharp)$ is isomorphic to $(f_{small}, f^\sharp)$
and moreover $f$ is an isomorphism of algebraic spaces.
\end{lemma}

\begin{proof}
By
Theorem \ref{theorem-fully-faithful}
it suffices to show that $(g, g^\sharp)$ is a morphism of
locally ringed topoi. By
Modules on Sites, Lemma \ref{sites-modules-lemma-locally-ringed-morphism}
(and since the site $X_\etale$ has enough points)
it suffices to check that the map
$\mathcal{O}_{Y, q} \to \mathcal{O}_{X, p}$ induced by $g^\sharp$
is a local ring map where $q = f \circ p$ and $p$ is any point of
$X_\etale$. As it is an isomorphism this is clear.
\end{proof}














\section{Quasi-coherent sheaves on algebraic spaces}
\label{section-quasi-coherent}

\noindent
In
Descent, Sections \ref{descent-section-quasi-coherent-sheaves},
\ref{descent-section-quasi-coherent-cohomology}, and
\ref{descent-section-quasi-coherent-sheaves-bis}
we have seen that for a scheme $U$, there is no difference between a
quasi-coherent $\mathcal{O}_U$-module on $U$, or a quasi-coherent
$\mathcal{O}$-module on the small \'etale site of $U$. Hence the following
definition is compatible with our original notion of a quasi-coherent sheaf
on a scheme
(Schemes, Section \ref{schemes-section-quasi-coherent}),
when applied to a representable algebraic space.

\begin{definition}
\label{definition-quasi-coherent}
Let $S$ be a scheme. Let $X$ be an algebraic space over $S$.
A {\it quasi-coherent} $\mathcal{O}_X$-module
is a quasi-coherent module on the ringed site
$(X_\etale, \mathcal{O}_X)$ in the sense of
Modules on Sites,
Definition \ref{sites-modules-definition-site-local}.
The category of quasi-coherent sheaves on $X$ is denoted
$\QCoh(\mathcal{O}_X)$.
\end{definition}

\noindent
Note that as being quasi-coherent is an intrinsic notion (see
Modules on Sites, Lemma \ref{sites-modules-lemma-special-locally-free})
this is equivalent to saying that the corresponding $\mathcal{O}_X$-module
on $X_{spaces, \etale}$ is quasi-coherent.

\medskip\noindent
As usual, quasi-coherent sheaves behave well with respect to pullback.

\begin{lemma}
\label{lemma-pullback-quasi-coherent}
Let $S$ be a scheme.
Let $f : X \to Y$ be a morphism of algebraic spaces over $S$.
The pullback functor
$f^* : \textit{Mod}(\mathcal{O}_Y) \to \textit{Mod}(\mathcal{O}_X)$
preserves quasi-coherent sheaves.
\end{lemma}

\begin{proof}
This is a general fact, see
Modules on Sites, Lemma \ref{sites-modules-lemma-local-pullback}.
\end{proof}

\noindent
Note that this pullback functor agrees with the usual pullback functor
between quasi-coherent sheaves of modules if $X$ and $Y$ happen to be
schemes, see
Descent, Proposition
\ref{descent-proposition-equivalence-quasi-coherent-functorial}.
Here is the obligatory lemma comparing this with quasi-coherent sheaves
on the objects of the small \'etale site of $X$.

\begin{lemma}
\label{lemma-characterize-quasi-coherent-small-etale}
Let $S$ be a scheme. Let $X$ be an algebraic space over $S$.
A quasi-coherent $\mathcal{O}_X$-module $\mathcal{F}$
is given by the following data:
\begin{enumerate}
\item for every $U \in \Ob(X_\etale)$ a quasi-coherent
$\mathcal{O}_U$-module $\mathcal{F}_U$ on $U_\etale$,
\item for every $f : U' \to U$ in $X_\etale$ an isomorphism
$c_f : f_{small}^*\mathcal{F}_U \to \mathcal{F}_{U'}$.
\end{enumerate}
These data are subject to the condition that given any $f : U' \to U$
and $g : U'' \to U'$ in $X_\etale$ the composition
$c_g \circ g_{small}^*c_f$ is equal to $c_{f \circ g}$.
\end{lemma}

\begin{proof}
Combine Lemmas \ref{lemma-pullback-quasi-coherent} and
\ref{lemma-characterize-module-small-etale}.
\end{proof}

\begin{lemma}
\label{lemma-stalk-quasi-coherent}
Let $S$ be a scheme.
Let $X$ be an algebraic space over $S$.
Let $\mathcal{F}$ be a quasi-coherent $\mathcal{O}_X$-module.
Let $x \in |X|$ be a point and let $\overline{x}$ be a geometric
point lying over $x$. Finally, let
$\varphi : (U, \overline{u}) \to (X, \overline{x})$
be an \'etale neighbourhood where $U$ is a scheme.
Then
$$
(\varphi^*\mathcal{F})_u \otimes_{\mathcal{O}_{U, u}}
\mathcal{O}_{X, \overline{x}} =
\mathcal{F}_{\overline{x}}
$$
where $u \in U$ is the image of $\overline{u}$.
\end{lemma}

\begin{proof}
Note that $\mathcal{O}_{X, \overline{x}} = \mathcal{O}_{U, u}^{sh}$ by
Lemma \ref{lemma-describe-etale-local-ring}
hence the tensor product makes sense. Moreover, from
Definition \ref{definition-stalk}
it is clear that
$$
\mathcal{F}_{\overline{u}} = \colim (\varphi^*\mathcal{F})_u
$$
where the colimit is over $\varphi : (U, \overline{u}) \to (X, \overline{x})$
as in the lemma. Hence there is a canonical map from left to right in
the statement of the lemma. We have a similar colimit description for
$\mathcal{O}_{X, \overline{x}}$
and by
Lemma \ref{lemma-characterize-quasi-coherent-small-etale}
we have
$$
((\varphi')^*\mathcal{F})_{u'} =
(\varphi^*\mathcal{F})_u \otimes_{\mathcal{O}_{U, u}} \mathcal{O}_{U', u'}
$$
whenever $(U', \overline{u}') \to (U, \overline{u})$ is a morphism of
\'etale neighbourhoods. To complete the proof we use that
$\otimes$ commutes with colimits.
\end{proof}

\begin{lemma}
\label{lemma-stalk-pullback-quasi-coherent}
Let $S$ be a scheme. Let $f : X \to Y$ be a morphism of algebraic spaces
over $S$. Let $\mathcal{G}$ be a quasi-coherent $\mathcal{O}_Y$-module.
Let $\overline{x}$ be a geometric point of $X$ and let
$\overline{y} = f \circ \overline{x}$ be the image in $Y$.
Then there is a canonical isomorphism
$$
(f^*\mathcal{G})_{\overline{x}} =
\mathcal{G}_{\overline{y}} \otimes_{\mathcal{O}_{Y, \overline{y}}}
\mathcal{O}_{X, \overline{x}}
$$
of the stalk of the pullback with the tensor product of the stalk
with the local ring of $X$ at $\overline{x}$.
\end{lemma}

\begin{proof}
Since $f^*\mathcal{G} =
f_{small}^{-1}\mathcal{G} \otimes_{f_{small}^{-1}\mathcal{O}_Y} \mathcal{O}_X$
this follows from the description of stalks of pullbacks in
Lemma \ref{lemma-stalk-pullback}
and the fact that taking stalks commutes with tensor products.
A more direct way to see this is as follows.
Choose a commutative diagram
$$
\xymatrix{
U \ar[d]_p \ar[r]_\alpha  & V \ar[d]^q \\
X \ar[r]^a & Y
}
$$
where $U$ and $V$ are schemes, and $p$ and $q$ are surjective \'etale.
By
Lemma \ref{lemma-geometric-lift-to-usual}
we can choose a geometric point $\overline{u}$ of $U$ such that
$\overline{x} = p \circ \overline{u}$. Set
$\overline{v} = \alpha \circ \overline{u}$.
Then we see that
\begin{align*}
(f^*\mathcal{G})_{\overline{x}} & =
(p^*f^*\mathcal{G})_u \otimes_{\mathcal{O}_{U, u}}
\mathcal{O}_{X, \overline{x}} \\
& = (\alpha^*q^*\mathcal{G})_u \otimes_{\mathcal{O}_{U, u}}
\mathcal{O}_{X, \overline{x}} \\
& = (q^*\mathcal{G})_v \otimes_{\mathcal{O}_{V, v}}
\mathcal{O}_{U, u} \otimes_{\mathcal{O}_{U, u}}
\mathcal{O}_{X, \overline{x}} \\
& = (q^*\mathcal{G})_v \otimes_{\mathcal{O}_{V, v}}
\mathcal{O}_{X, \overline{x}} \\
& = (q^*\mathcal{G})_v \otimes_{\mathcal{O}_{V, v}}
\mathcal{O}_{Y, \overline{y}} \otimes_{\mathcal{O}_{Y, \overline{y}}}
\mathcal{O}_{X, \overline{x}} \\
& = \mathcal{G}_{\overline{y}} \otimes_{\mathcal{O}_{Y, \overline{y}}}
\mathcal{O}_{X, \overline{x}}
\end{align*}
Here we have used
Lemma \ref{lemma-stalk-quasi-coherent} (twice)
and the corresponding result for pullbacks of quasi-coherent sheaves
on schemes, see
Sheaves, Lemma \ref{sheaves-lemma-stalk-pullback-modules}.
\end{proof}

\begin{lemma}
\label{lemma-characterize-quasi-coherent}
Let $S$ be a scheme. Let $X$ be an algebraic space over $S$.
Let $\mathcal{F}$ be a sheaf of $\mathcal{O}_X$-modules.
The following are equivalent
\begin{enumerate}
\item $\mathcal{F}$ is a quasi-coherent $\mathcal{O}_X$-module,
\item there exists an \'etale morphism $f : Y \to X$ of
algebraic spaces over $S$ with $|f| : |Y| \to |X|$ surjective
such that $f^*\mathcal{F}$ is quasi-coherent on $Y$,
\item there exists a scheme $U$ and a surjective \'etale morphism
$\varphi : U \to X$ such that $\varphi^*\mathcal{F}$ is a quasi-coherent
$\mathcal{O}_U$-module, and
\item for every affine scheme $U$ and \'etale morphism $\varphi : U \to X$ the
restriction $\varphi^*\mathcal{F}$ is a quasi-coherent $\mathcal{O}_U$-module.
\end{enumerate}
\end{lemma}

\begin{proof}
It is clear that (1) implies (2) by considering $\text{id}_X$.
Assume $f : Y \to X$ is as in (2), and let $V \to Y$ be a surjective
\'etale morphism from a scheme towards $Y$. Then the composition $V \to X$ is
surjective \'etale as well
and by Lemma \ref{lemma-pullback-quasi-coherent} the pullback of $\mathcal{F}$
to $V$ is quasi-coherent as well. Hence we see that (2) implies (3).

\medskip\noindent
Let $U \to X$ be as in (3). Let us use the abuse of notation introduced
in Equation (\ref{equation-restrict-modules}).
As $\mathcal{F}|_{U_\etale}$ is quasi-coherent there exists an
\'etale covering $\{U_i \to U\}$ such that
$\mathcal{F}|_{U_{i, \etale}}$ has a global presentation, see
Modules on Sites, Definition \ref{sites-modules-definition-global} and
Lemma \ref{sites-modules-lemma-local-final-object}.
Let $V \to X$ be an object of $X_\etale$. Since $U \to X$ is
surjective and \'etale, the family of maps $\{U_i \times_X V \to V\}$ is an
\'etale covering
of $V$. Via the morphisms $U_i \times_X V \to U_i$ we can restrict the
global presentations of $\mathcal{F}|_{U_{i, \etale}}$ to get a global
presentation of $\mathcal{F}|_{(U_i \times_X V)_\etale}$
Hence the sheaf $\mathcal{F}$ on $X_\etale$ satisfies the condition of
Modules on Sites, Definition \ref{sites-modules-definition-site-local}
and hence is quasi-coherent.

\medskip\noindent
The equivalence of (3) and (4) comes from the fact that any scheme has
an affine open covering.
\end{proof}

\begin{lemma}
\label{lemma-properties-quasi-coherent}
Let $S$ be a scheme. Let $X$ be an algebraic space over $S$.
The category $\QCoh(\mathcal{O}_X)$ of quasi-coherent sheaves on $X$
has the following properties:
\begin{enumerate}
\item Any direct sum of quasi-coherent sheaves is quasi-coherent.
\item Any colimit of quasi-coherent sheaves is quasi-coherent.
\item The kernel and cokernel of a morphism of quasi-coherent sheaves
is quasi-coherent.
\item Given a short exact sequence of $\mathcal{O}_X$-modules
$0 \to \mathcal{F}_1 \to \mathcal{F}_2 \to \mathcal{F}_3 \to 0$
if two out of three are quasi-coherent so is the third.
\item Given two quasi-coherent $\mathcal{O}_X$-modules
the tensor product is quasi-coherent.
\item Given two quasi-coherent $\mathcal{O}_X$-modules
$\mathcal{F}$, $\mathcal{G}$ such that $\mathcal{F}$
is of finite presentation (see
Section \ref{section-properties-modules}),
then the internal hom
$\SheafHom_{\mathcal{O}_X}(\mathcal{F}, \mathcal{G})$
is quasi-coherent.
\end{enumerate}
\end{lemma}

\begin{proof}
If $X$ is a scheme, then this is
Descent, Lemma \ref{descent-lemma-properties-quasi-coherent}.
We will reduce the lemma to this case by \'etale localization.

\medskip\noindent
Choose a scheme $U$ and a surjective \'etale morphism $\varphi : U \to X$.
Our notation will be that $\textit{Mod}(\mathcal{O}_U) =
\textit{Mod}(U_\etale, \mathcal{O}_U)$ and
$\QCoh(\mathcal{O}_U) = \QCoh(U_\etale, \mathcal{O}_U)$; in other
words, even though $U$ is a scheme we think of quasi-coherent
modules on $U$ as modules on the small \'etale site of $U$.
By Lemma \ref{lemma-pullback-quasi-coherent} we have a commutative diagram
$$
\xymatrix{
\QCoh(\mathcal{O}_X) \ar[r]_{\varphi^*} \ar[d] &
\QCoh(\mathcal{O}_U) \ar[d] \\
\textit{Mod}(\mathcal{O}_X) \ar[r]^{\varphi^*} &
\textit{Mod}(\mathcal{O}_U)
}
$$
The bottom horizontal arrow is the restriction functor
(\ref{equation-restrict-modules})
$\mathcal{G} \mapsto \mathcal{G}|_{U_\etale}$.
This functor has both a left adjoint and a right adjoint, see
Modules on Sites, Section \ref{sites-modules-section-localize},
hence commutes with all limits and colimits.
Moreover, we know that an object of $\textit{Mod}(\mathcal{O}_X)$ is in
$\QCoh(\mathcal{O}_X)$ if and only if its restriction to $U$ is in
$\QCoh(\mathcal{O}_U)$, see
Lemma \ref{lemma-characterize-quasi-coherent}.
With these preliminaries out of the way we can start the proof.

\medskip\noindent
Proof of (1). Let $\mathcal{F}_i$, $i \in I$ be a family of quasi-coherent
$\mathcal{O}_X$-modules. By the discussion above we have
$$
\Big(\bigoplus \mathcal{F}_i\Big)|_{U_\etale} =
\bigoplus \mathcal{F}_i|_{U_\etale}
$$
Each of the modules $\mathcal{F}_i|_{U_\etale}$ is quasi-coherent.
Hence the direct sum is quasi-coherent by the case of schemes.
Hence $\bigoplus \mathcal{F}_i$ is quasi-coherent as a module
restricting to a quasi-coherent module on $U$.

\medskip\noindent
Proof of (2). Let $\mathcal{I} \to \QCoh(\mathcal{O}_X)$,
$i \mapsto \mathcal{F}_i$ be a diagram. Then
$$
(\colim \mathcal{F}_i)|_{U_\etale} = \colim \mathcal{F}_i|_{U_\etale}
$$
by the discussion above and we conclude in the same manner.

\medskip\noindent
Proof of (3). Let $a : \mathcal{F} \to \mathcal{F}'$
be an arrow of $\QCoh(\mathcal{O}_X)$. Then
we have $\Ker(a)|_{U_\etale} = \Ker(a|_{U_\etale})$
and $\Coker(a)|_{U_\etale} = \Coker(a|_{U_\etale})$
and we conclude in the same manner.

\medskip\noindent
Proof of (4). The restriction
$0 \to \mathcal{F}_1|_{U_\etale} \to \mathcal{F}_2|_{U_\etale}
\to \mathcal{F}_3|_{U_\etale} \to 0$
is short exact. Hence we have the 2-out-of-3 property
for this sequence and we conclude as before.

\medskip\noindent
Proof of (5). Let $\mathcal{F}$ and $\mathcal{G}$ be in $\QCoh(\mathcal{O}_X)$.
Then we have
$$
(\mathcal{F} \otimes_{\mathcal{O}_X} \mathcal{G})_{U_\etale} =
\mathcal{F}|_{U_\etale} \otimes_{\mathcal{O}_U} \mathcal{G}|_{U_\etale}
$$
and we conclude as before.

\medskip\noindent
Proof of (6). Let $\mathcal{F}$ and $\mathcal{G}$
be in $\QCoh(\mathcal{O}_X)$ with $\mathcal{F}$
of finite presentation. We have
$$
\SheafHom_{\mathcal{O}_X}(\mathcal{F}, \mathcal{G})|_{U_\etale} =
\SheafHom_{\mathcal{O}_U}(\mathcal{F}|_{U_\etale}, \mathcal{G}|_{U_\etale})
$$
Namely, restriction is a localization, see
Section \ref{section-localize}, especially formula
(\ref{equation-localize-at-scheme-ringed})) and formation of internal
hom commutes with localization, see
Modules on Sites, Lemma \ref{sites-modules-lemma-internal-hom-restriction}.
Thus we conclude as before.
\end{proof}

\noindent
It is in general not the case that the pushforward of a quasi-coherent sheaf
along a morphism of algebraic spaces is quasi-coherent. We will return to this
issue in
Morphisms of Spaces, Section \ref{spaces-morphisms-section-pushforward}.




\section{Properties of modules}
\label{section-properties-modules}

\noindent
In
Modules on Sites, Sections
\ref{sites-modules-section-global},
\ref{sites-modules-section-local}, and
Definition \ref{sites-modules-definition-flat}
we have defined a number of intrinsic properties of modules of
$\mathcal{O}$-module on any ringed topos. If $X$ is an algebraic
space, we will apply these notions freely to modules on the ringed
site $(X_\etale, \mathcal{O}_X)$, or equivalently on the ringed site
$(X_{spaces, \etale}, \mathcal{O}_X)$.

\medskip\noindent
Global properties $\mathcal{P}$:
\begin{enumerate}
\item[(a)] {\it free},
\item[(b)] {\it finite free},
\item[(c)] {\it generated by global sections},
\item[(d)] {\it generated by finitely many global sections},
\item[(e)] having a {\it global presentation}, and
\item[(f)] having a {\it global finite presentation}.
\end{enumerate}
Local properties $\mathcal{P}$:
\begin{enumerate}
\item[(g)] {\it locally free},
\item[(f)] {\it finite locally free},
\item[(h)] {\it locally generated by sections},
\item[(i)] {\it locally generated by $r$ sections},
\item[(j)] {\it finite type},
\item[(k)] {\it quasi-coherent} (see Section \ref{section-quasi-coherent}),
\item[(l)] {\it of finite presentation},
\item[(m)] {\it coherent}, and
\item[(n)] {\it flat}.
\end{enumerate}
Here are some results which follow immediately from the definitions:
\begin{enumerate}
\item In each case, except for $\mathcal{P}=$``coherent'', the property
is preserved under pullback, see
Modules on Sites, Lemmas \ref{sites-modules-lemma-global-pullback},
\ref{sites-modules-lemma-local-pullback}, and
\ref{sites-modules-lemma-pullback-flat}.
\item Each of the properties above (including coherent) are preserved under
pullbacks by \'etale morphisms of algebraic spaces (because in this
case pullback is given by restriction, see
Lemma \ref{lemma-etale-morphism-topoi}).
\item Assume $f : Y \to X$ is a surjective \'etale morphism of algebraic
spaces. For each of the local properties (g) -- (m), the fact that
$f^*\mathcal{F}$ has $\mathcal{P}$ implies that $\mathcal{F}$ has
$\mathcal{P}$. This follows as $\{Y \to X\}$ is a covering in
$X_{spaces, \etale}$ and
Modules on Sites, Lemma \ref{sites-modules-lemma-local-final-object}.
\item If $X$ is a scheme, $\mathcal{F}$ is a quasi-coherent module
on $X_\etale$, and $\mathcal{P}$ any property except ``coherent'' or
``locally free'', then $\mathcal{P}$ for $\mathcal{F}$ on $X_\etale$
is equivalent to the corresponding property for
$\mathcal{F}|_{X_{Zar}}$, i.e., it corresponds to $\mathcal{P}$
for $\mathcal{F}$ when we think of it as a quasi-coherent sheaf
on the scheme $X$. See
Descent, Lemma \ref{descent-lemma-equivalence-quasi-coherent-properties}.
\item If $X$ is a locally Noetherian scheme, $\mathcal{F}$ is a
quasi-coherent module on $X_\etale$, then $\mathcal{F}$
is coherent on $X_\etale$ if and only if
$\mathcal{F}|_{X_{Zar}}$ is coherent, i.e., it corresponds to
the usual notion of a coherent sheaf on the scheme $X$ being
coherent. See
Descent, Lemma \ref{descent-lemma-equivalence-quasi-coherent-properties}.
\end{enumerate}



\section{Locally projective modules}
\label{section-locally-projective}

\noindent
Recall that in
Properties, Section \ref{properties-section-locally-projective}
we defined the notion of a locally projective
quasi-coherent module.

\begin{lemma}
\label{lemma-locally-projective}
Let $S$ be a scheme. Let $X$ be an algebraic space over $S$.
Let $\mathcal{F}$ be a quasi-coherent $\mathcal{O}_X$-module.
The following are equivalent
\begin{enumerate}
\item for some scheme $U$ and surjective \'etale morphism
$U \to X$ the restriction $\mathcal{F}|_U$ is locally projective
on $U$, and
\item for any scheme $U$ and any \'etale morphism
$U \to X$ the restriction $\mathcal{F}|_U$ is locally projective
on $U$.
\end{enumerate}
\end{lemma}

\begin{proof}
Let $U \to X$ be as in (1) and let $V \to X$ be \'etale where
$V$ is a scheme. Then $\{U \times_X V \to V\}$ is an fppf covering
of schemes. Hence if $\mathcal{F}|_U$ is locally projective, then
$\mathcal{F}|_{U \times_X V}$ is locally projective (see
Properties, Lemma \ref{properties-lemma-locally-projective-pullback})
and hence $\mathcal{F}|_V$ is locally projective, see
Descent, Lemma \ref{descent-lemma-locally-projective-descends}.
\end{proof}

\begin{definition}
\label{definition-locally-projective}
Let $S$ be a scheme. Let $X$ be an algebraic space over $S$.
Let $\mathcal{F}$ be a quasi-coherent $\mathcal{O}_X$-module.
We say $\mathcal{F}$ is {\it locally projective}
if the equivalent conditions of
Lemma \ref{lemma-locally-projective}
are satisfied.
\end{definition}

\begin{lemma}
\label{lemma-locally-projective-pullback}
Let $S$ be a scheme.
Let $f : X \to Y$ be a morphism of algebraic spaces over $S$.
Let $\mathcal{G}$ be a quasi-coherent $\mathcal{O}_Y$-module.
If $\mathcal{G}$ is locally projective on $Y$, then $f^*\mathcal{G}$
is locally projective on $X$.
\end{lemma}

\begin{proof}
Choose a surjective \'etale morphism $V \to Y$ with $V$ a scheme.
Choose a surjective \'etale morphism $U \to V \times_Y X$ with
$U$ a scheme. Denote $\psi : U \to V$ the induced morphism.
Then
$$
f^*\mathcal{G}|_U = \psi^*(\mathcal{G}|_V)
$$
Hence the lemma follows from the definition and the result in the
case of schemes, see
Properties, Lemma \ref{properties-lemma-locally-projective-pullback}.
\end{proof}





\section{Quasi-coherent sheaves and presentations}
\label{section-quasi-coherent-presentation}

\noindent
Let $S$ be a scheme. Let $X$ be an algebraic space over $S$.
Let $X = U/R$ be a presentation of $X$ coming from any surjective
\'etale morphism $\varphi : U \to X$, see
Spaces, Definition \ref{spaces-definition-presentation}.
In particular, we obtain a groupoid $(U, R, s, t, c)$, such that
$j = (t, s) : R \to U \times_S U$, see
Groupoids, Lemma \ref{groupoids-lemma-equivalence-groupoid}.
In
Groupoids, Definition \ref{groupoids-definition-groupoid-module}
we have the defined the notion of a quasi-coherent sheaf
on an arbitrary groupoid. With these notions in place we have
the following observation.

\begin{proposition}
\label{proposition-quasi-coherent}
With $S$, $\varphi : U \to X$, and $(U, R, s, t, c)$ as above.
For any quasi-coherent $\mathcal{O}_X$-module $\mathcal{F}$ the
sheaf $\varphi^*\mathcal{F}$ comes equipped with a canonical
isomorphism
$$
\alpha : t^*\varphi^*\mathcal{F} \longrightarrow s^*\varphi^*\mathcal{F}
$$
which satisfies the conditions of
Groupoids, Definition \ref{groupoids-definition-groupoid-module}
and therefore defines a quasi-coherent sheaf on $(U, R, s, t, c)$.
The functor $\mathcal{F} \mapsto (\varphi^*\mathcal{F}, \alpha)$
defines an equivalence of categories
$$
\begin{matrix}
\text{Quasi-coherent} \\
\mathcal{O}_X\text{-modules}
\end{matrix}
\longleftrightarrow
\begin{matrix}
\text{Quasi-coherent modules}\\
\text{on }(U, R, s, t, c)
\end{matrix}
$$
\end{proposition}

\begin{proof}
In the statement of the proposition, and in this proof we think of a
quasi-coherent sheaf on a scheme as a quasi-coherent sheaf on the small
\'etale site of that scheme. This is permissible by the results of
Descent, Sections \ref{descent-section-quasi-coherent-sheaves},
\ref{descent-section-quasi-coherent-cohomology}, and
\ref{descent-section-quasi-coherent-sheaves-bis}.

\medskip\noindent
The existence of $\alpha$ comes from the fact that
$\varphi \circ t = \varphi \circ s$ and that pullback is
functorial in the morphism, see discussion surrounding
Equation (\ref{equation-push-pull}). In exactly the same way, i.e., by
functoriality of pullback, we see that the isomorphism $\alpha$ satisfies
condition (1) of
Groupoids, Definition \ref{groupoids-definition-groupoid-module}.
To see condition (2) of the definition it suffices to see that $\alpha$
is an isomorphism which is clear. The construction
$\mathcal{F} \mapsto (\varphi^*\mathcal{F}, \alpha)$
is clearly functorial in the quasi-coherent sheaf $\mathcal{F}$.
Hence we obtain the functor from left to right in the displayed
formula of the lemma.

\medskip\noindent
Conversely, suppose that $(\mathcal{F}, \alpha)$ is a quasi-coherent
sheaf on $(U, R, s, t, c)$. Let $V \to X$ be an object of $X_\etale$.
In this case the morphism $V' = U \times_X V \to V$ is a surjective \'etale
morphism of schemes, and hence $\{V' \to V\}$ is an \'etale
covering of $V$. Moreover, the quasi-coherent sheaf $\mathcal{F}$
pulls back to a quasi-coherent sheaf $\mathcal{F}'$ on $V'$.
Since $R = U \times_X U$ with $t = \text{pr}_0$ and $s = \text{pr}_0$
we see that $V' \times_V V' = R \times_X V$ with projection maps
$V' \times_V V' \to V'$ equal to the pullbacks of $t$ and $s$. Hence
$\alpha$ pulls back to an isomorphism
$\alpha' : \text{pr}_0^*\mathcal{F}' \to \text{pr}_1^*\mathcal{F}'$, and
the pair $(\mathcal{F}', \alpha')$ is a descend datum for quasi-coherent
sheaves with respect to $\{V' \to V\}$. By
Descent, Proposition
\ref{descent-proposition-fpqc-descent-quasi-coherent}
this descent datum is effective, and we obtain a quasi-coherent
$\mathcal{O}_V$-module $\mathcal{F}_V$ on $V_\etale$.
To see that this gives a quasi-coherent sheaf on $X_\etale$ we have
to show (by
Lemma \ref{lemma-characterize-quasi-coherent-small-etale})
that for any morphism $f : V_1 \to V_2$ in $X_\etale$
there is a canonical isomorphism
$c_f : \mathcal{F}_{V_1} \to \mathcal{F}_{V_2}$
compatible with compositions of morphisms. We omit the verification.
We also omit the verification that this defines a functor from the
category on the right to the category on the left which is inverse
to the functor described above.
\end{proof}

\begin{proposition}
\label{proposition-coherator}
Let $S$ be a scheme. Let $X$ be an algebraic space over $S$.
\begin{enumerate}
\item The category $\QCoh(\mathcal{O}_X)$ is a Grothendieck
abelian category. Consequently, $\QCoh(\mathcal{O}_X)$
has enough injectives and all limits.
\item The inclusion functor
$\QCoh(\mathcal{O}_X) \to \textit{Mod}(\mathcal{O}_X)$
has a right adjoint\footnote{This functor is sometimes called
the {\it coherator}.}
$$
Q : \textit{Mod}(\mathcal{O}_X) \longrightarrow \QCoh(\mathcal{O}_X)
$$
such that for every quasi-coherent sheaf $\mathcal{F}$ the adjunction mapping
$Q(\mathcal{F}) \to \mathcal{F}$ is an isomorphism.
\end{enumerate}
\end{proposition}

\begin{proof}
This proof is a repeat of the proof in the case of schemes, see
Properties, Proposition \ref{properties-proposition-coherator}.
We advise the reader to read that proof first.

\medskip\noindent
Part (1) means $\QCoh(\mathcal{O}_X)$ (a) has all colimits,
(b) filtered colimits are exact, and (c) has a generator, see
Injectives, Section \ref{injectives-section-grothendieck-conditions}.
By Lemma \ref{lemma-properties-quasi-coherent}
colimits in $\QCoh(\mathcal{O}_X)$ exist and agree
with colimits in $\textit{Mod}(\mathcal{O}_X)$. By
Modules on Sites, Lemma \ref{sites-modules-lemma-limits-colimits}
filtered colimits are exact. Hence (a) and (b) hold.

\medskip\noindent
To construct a generator, choose a presentation $X = U/R$ so that
$(U, R, s, t, c)$ is an
\'etale groupoid scheme and in particular $s$ and $t$ are flat morphisms
of schemes. Pick a cardinal $\kappa$ as in
Groupoids, Lemma \ref{groupoids-lemma-colimit-kappa}.
Pick a collection $(\mathcal{E}_t, \alpha_t)_{t \in T}$ of
$\kappa$-generated quasi-coherent modules on
$(U, R, s, t, c)$ as in
Groupoids, Lemma \ref{groupoids-lemma-set-of-iso-classes}.
Let $\mathcal{F}_t$ be the quasi-coherent module on $X$ which
corresponds to the quasi-coherent module $(\mathcal{E}_t, \alpha_t)$ via
the equivalence of categories of
Proposition \ref{proposition-quasi-coherent}.
Then we see that every quasi-coherent module $\mathcal{H}$ is the
directed colimit of its quasi-coherent submodules which are isomorphic
to one of the $\mathcal{F}_t$. Thus $\bigoplus_t \mathcal{F}_t$ is
a generator of $\QCoh(\mathcal{O}_X)$ and we conclude that (c) holds.
The assertions on limits and injectives hold in any
Grothendieck abelian category, see
Injectives, Theorem
\ref{injectives-theorem-injective-embedding-grothendieck} and
Lemma \ref{injectives-lemma-grothendieck-products}.

\medskip\noindent
Proof of (2). To construct $Q$ we use the following general procedure.
Given an object $\mathcal{F}$ of $\textit{Mod}(\mathcal{O}_X)$
we consider the functor
$$
\QCoh(\mathcal{O}_X)^{opp} \longrightarrow \textit{Sets},\quad
\mathcal{G} \longmapsto \Hom_X(\mathcal{G}, \mathcal{F})
$$
This functor transforms colimits into limits,
hence is representable, see
Injectives, Lemma \ref{injectives-lemma-grothendieck-brown}.
Thus there exists a quasi-coherent sheaf $Q(\mathcal{F})$
and a functorial isomorphism
$\Hom_X(\mathcal{G}, \mathcal{F}) = \Hom_X(\mathcal{G}, Q(\mathcal{F}))$
for $\mathcal{G}$ in $\QCoh(\mathcal{O}_X)$. By the Yoneda lemma
(Categories, Lemma \ref{categories-lemma-yoneda})
the construction $\mathcal{F} \leadsto Q(\mathcal{F})$ is
functorial in $\mathcal{F}$. By construction $Q$ is a right
adjoint to the inclusion functor.
The fact that $Q(\mathcal{F}) \to \mathcal{F}$ is an isomorphism
when $\mathcal{F}$ is quasi-coherent is a formal consequence of the fact
that the inclusion functor
$\QCoh(\mathcal{O}_X) \to \textit{Mod}(\mathcal{O}_X)$
is fully faithful.
\end{proof}






\section{Morphisms towards schemes}
\label{section-morphisms-to-schemes}

\noindent
Here is the analogue of
Schemes, Lemma \ref{schemes-lemma-morphism-into-affine}.

\begin{lemma}
\label{lemma-morphism-to-affine-scheme}
Let $X$ be an algebraic space over $\mathbf{Z}$.
Let $T$ be an affine scheme.
The map
$$
\Mor(X, T)
\longrightarrow
\Hom(\Gamma(T, \mathcal{O}_T), \Gamma(X, \mathcal{O}_X))
$$
which maps $f$ to $f^\sharp$ (on global sections) is bijective.
\end{lemma}

\begin{proof}
We construct the inverse of the map.
Let $\varphi : \Gamma(T, \mathcal{O}_T) \to \Gamma(X, \mathcal{O}_X)$
be a ring map. Choose a presentation $X = U/R$, see
Spaces, Definition \ref{spaces-definition-presentation}.
By
Schemes, Lemma \ref{schemes-lemma-morphism-into-affine}
the composition
$$
\Gamma(T, \mathcal{O}_T) \to \Gamma(X, \mathcal{O}_X) \to
\Gamma(U, \mathcal{O}_U)
$$
corresponds to a unique morphism of schemes $g : U \to T$. By the same lemma
the two compositions $R \to U \to T$ are equal. Hence we obtain a morphism
$f : X = U/R \to T$ such that $U \to X \to T$ equals $g$. By construction
the diagram
$$
\xymatrix{
\Gamma(U, \mathcal{O}_U) & \Gamma(X, \mathcal{O}_X) \ar[l] \\
& \Gamma(T, \mathcal{O}_T) \ar[lu]^{g^\sharp} \ar[u]^{\varphi}_{f^\sharp}
}
$$
commutes. Hence $f^\sharp$ equals $\varphi$ because $U \to X$ is an
\'etale covering and $\mathcal{O}_X$ is a sheaf on $X_\etale$.
The uniqueness of $f$ follows from the uniqueness of $g$.
\end{proof}





\section{Quotients by free actions}
\label{section-quotient-by-free}

\noindent
Let $S$ be a scheme.
Let $X$ be an algebraic space over $S$.
Let $G$ be an abstract group.
Let $a : G \to \text{Aut}(X)$ be a homomorphism, i.e., $a$ is an
{\it action} of $G$ on $X$. We will say the action is {\it free}
if for every scheme $T$ over $S$ the map
$$
G \times X(T) \longrightarrow X(T)
$$
is free. (We cannot use a criterion as in
Spaces, Lemma \ref{spaces-lemma-quotient}
because points may not have well defined residue fields.)
In case the action is free we're going to construct the quotient $X/G$
as an algebraic space. This is a special case of the general
Bootstrap, Lemma \ref{bootstrap-lemma-quotient-free-action}
that we will prove later.

\begin{lemma}
\label{lemma-quotient}
Let $S$ be a scheme.
Let $X$ be an algebraic space over $S$.
Let $G$ be an abstract group with a free action on $X$.
Then the quotient sheaf $X/G$ is an algebraic space.
\end{lemma}

\begin{proof}
The statement means that the sheaf $F$ associated to the presheaf
$$
T \longmapsto X(T)/G
$$
is an algebraic space. To see this we will construct a presentation.
Namely, choose a scheme $U$ and a surjective \'etale morphism
$\varphi : U \to X$. Set $V = \coprod_{g \in G} U$ and set
$\psi : V \to X$ equal to $a(g) \circ \varphi$ on the component corresponding
to $g \in G$. Let $G$ act on $V$ by permuting the components, i.e.,
$g_0 \in G$ maps the component corresponding to $g$ to the component
corresponding to $g_0g$ via the identity morphism of $U$.
Then $\psi$ is a $G$-equivariant morphism, i.e., we reduce to the
case dealt with in the next paragraph.

\medskip\noindent
Assume that there exists a $G$-action on $U$ and that $U \to X$ is surjective,
\'etale and $G$-equivariant. In this case there is an induced
action of $G$ on $R = U \times_X U$ compatible with the projection
mappings $t, s : R \to U$. Now we claim that
$$
X/G = U/\coprod\nolimits_{g \in G} R
$$
where the map
$$
j : \coprod\nolimits_{g \in G} R
\longrightarrow
U \times_S U
$$
is given by $(r, g) \mapsto (t(r), g(s(r)))$. Note that $j$ is a monomorphism:
If $(t(r), g(s(r))) = (t(r'), g'(s(r')))$, then
$t(r) = t(r')$, hence $r$ and $r'$ have the same image in $X$ under
both $s$ and $t$, hence $g = g'$ (as $G$ acts freely on $X$), hence
$s(r) = s(r')$, hence $r = r'$ (as $R$ is an equivalence relation on $U$).
Moreover $j$ is an equivalence relation (details omitted).
Both projections $\coprod\nolimits_{g \in G} R \to U$ are \'etale, as
$s$ and $t$ are \'etale. Thus $j$ is an \'etale equivalence relation
and $U/\coprod\nolimits_{g \in G} R$ is an algebraic space by
Spaces, Theorem \ref{spaces-theorem-presentation}.
There is a map
$$
U/\coprod\nolimits_{g \in G} R \longrightarrow X/G
$$
induced by the map $U \to X$. We omit the proof that it is an
isomorphism of sheaves.
\end{proof}






\begin{multicols}{2}[\section{Autres chapitres}]
\noindent
Préliminaires
\begin{enumerate}
\item \hyperref[introduction-section-phantom]{Introduction}
\item \hyperref[conventions-section-phantom]{Conventions}
\item \hyperref[sets-section-phantom]{Théorie des ensembles}
\item \hyperref[categories-section-phantom]{Catégories}
\item \hyperref[topology-section-phantom]{Topologie}
\item \hyperref[sheaves-section-phantom]{Faisceaux sur les espaces}
\item \hyperref[sites-section-phantom]{Sites et faisceaux}
\item \hyperref[stacks-section-phantom]{Champs}
\item \hyperref[fields-section-phantom]{Corps}
\item \hyperref[algebra-section-phantom]{Algèbre commutative}
\item \hyperref[brauer-section-phantom]{Groupes de Brauer}
\item \hyperref[homology-section-phantom]{Algèbre homologique}
\item \hyperref[derived-section-phantom]{Catégories dérivées}
\item \hyperref[simplicial-section-phantom]{Méthodes simpliciales}
\item \hyperref[more-algebra-section-phantom]{Compléments d'algèbre}
\item \hyperref[smoothing-section-phantom]{Lissage des morphismes d'anneaux}
\item \hyperref[modules-section-phantom]{Faisceaux de modules}
\item \hyperref[sites-modules-section-phantom]{Modules sur les sites}
\item \hyperref[injectives-section-phantom]{Objets injectifs}
\item \hyperref[cohomology-section-phantom]{Cohomologie des faisceaux}
\item \hyperref[sites-cohomology-section-phantom]{Cohomologie sur les sites}
\item \hyperref[dga-section-phantom]{Algèbre différentielle graduée}
\item \hyperref[dpa-section-phantom]{Algèbre à puissances divisées}
\item \hyperref[sdga-section-phantom]{Faisceaux différentiels gradués}
\item \hyperref[hypercovering-section-phantom]{Hyperrecouvrements}
\end{enumerate}
Schémas
\begin{enumerate}
\setcounter{enumi}{25}
\item \hyperref[schemes-section-phantom]{Schémas}
\item \hyperref[constructions-section-phantom]{Constructions de schémas}
\item \hyperref[properties-section-phantom]{Propriétés des schémas}
\item \hyperref[morphisms-section-phantom]{Morphismes de schémas}
\item \hyperref[coherent-section-phantom]{Cohomologie des schémas}
\item \hyperref[divisors-section-phantom]{Diviseurs}
\item \hyperref[limits-section-phantom]{Limites de schémas}
\item \hyperref[varieties-section-phantom]{Variétés}
\item \hyperref[topologies-section-phantom]{Topologies sur les schémas}
\item \hyperref[descent-section-phantom]{Descente}
\item \hyperref[perfect-section-phantom]{Catégories dérivées de schémas}
\item \hyperref[more-morphisms-section-phantom]{Compléments sur les morphismes}
\item \hyperref[flat-section-phantom]{Compléments sur la platitude}
\item \hyperref[groupoids-section-phantom]{Schémas en groupoïdes}
\item \hyperref[more-groupoids-section-phantom]{Compléments sur les schémas en groupoïdes}
\item \hyperref[etale-section-phantom]{Morphismes étales de schémas}
\end{enumerate}
Thèmes de théorie des schémas
\begin{enumerate}
\setcounter{enumi}{41}
\item \hyperref[chow-section-phantom]{Homologie de Chow}
\item \hyperref[intersection-section-phantom]{Théorie de l'intersection}
\item \hyperref[pic-section-phantom]{Schémas de Picard des courbes}
\item \hyperref[weil-section-phantom]{Théories cohomologiques de Weil}
\item \hyperref[adequate-section-phantom]{Modules adéquats}
\item \hyperref[dualizing-section-phantom]{Complexes dualisants}
\item \hyperref[duality-section-phantom]{Dualité pour les schémas}
\item \hyperref[discriminant-section-phantom]{Discriminants et différentes}
\item \hyperref[derham-section-phantom]{Cohomologie de de Rham}
\item \hyperref[local-cohomology-section-phantom]{Cohomologie locale}
\item \hyperref[algebraization-section-phantom]{Géométrie algébrique et formelle}
\item \hyperref[curves-section-phantom]{Courbes algébriques}
\item \hyperref[resolve-section-phantom]{Résolution des surfaces}
\item \hyperref[models-section-phantom]{Réduction semi-stable}
\item \hyperref[functors-section-phantom]{Foncteurs et morphismes}
\item \hyperref[equiv-section-phantom]{Catégories dérivées de variétés}
\item \hyperref[pione-section-phantom]{Groupes fondamentaux de schémas}
\item \hyperref[etale-cohomology-section-phantom]{Cohomologie étale}
\item \hyperref[crystalline-section-phantom]{Cohomologie cristalline}
\item \hyperref[proetale-section-phantom]{Cohomologie pro-étale}
\item \hyperref[relative-cycles-section-phantom]{Cycles relatifs}
\item \hyperref[more-etale-section-phantom]{Compléments de cohomologie étale}
\item \hyperref[trace-section-phantom]{La formule des traces}
\end{enumerate}
Espaces algébriques
\begin{enumerate}
\setcounter{enumi}{64}
\item \hyperref[spaces-section-phantom]{Espaces algébriques}
\item \hyperref[spaces-properties-section-phantom]{Propriétés des espaces algébriques}
\item \hyperref[spaces-morphisms-section-phantom]{Morphismes d'espaces algébriques}
\item \hyperref[decent-spaces-section-phantom]{Espaces algébriques décents}
\item \hyperref[spaces-cohomology-section-phantom]{Cohomologie des espaces algébriques}
\item \hyperref[spaces-limits-section-phantom]{Limites d'espaces algébriques}
\item \hyperref[spaces-divisors-section-phantom]{Diviseurs sur les espaces algébriques}
\item \hyperref[spaces-over-fields-section-phantom]{Espaces algébriques sur des corps}
\item \hyperref[spaces-topologies-section-phantom]{Topologies sur les espaces algébriques}
\item \hyperref[spaces-descent-section-phantom]{Descente et espaces algébriques}
\item \hyperref[spaces-perfect-section-phantom]{Catégories dérivées d'espaces}
\item \hyperref[spaces-more-morphisms-section-phantom]{Compléments sur les morphismes d'espaces}
\item \hyperref[spaces-flat-section-phantom]{Platitude sur les espaces algébriques}
\item \hyperref[spaces-groupoids-section-phantom]{Groupoïdes dans les espaces algébriques}
\item \hyperref[spaces-more-groupoids-section-phantom]{Compléments sur les groupoïdes dans les espaces}
\item \hyperref[bootstrap-section-phantom]{Amorçage}
\item \hyperref[spaces-pushouts-section-phantom]{Sommes amalgamées d'espaces algébriques}
\end{enumerate}
Thèmes de géométrie
\begin{enumerate}
\setcounter{enumi}{81}
\item \hyperref[spaces-chow-section-phantom]{Groupes de Chow des espaces}
\item \hyperref[groupoids-quotients-section-phantom]{Quotients de groupoïdes}
\item \hyperref[spaces-more-cohomology-section-phantom]{Compléments sur la cohomologie des espaces}
\item \hyperref[spaces-simplicial-section-phantom]{Espaces simpliciaux}
\item \hyperref[spaces-duality-section-phantom]{Dualité pour les espaces}
\item \hyperref[formal-spaces-section-phantom]{Espaces algébriques formels}
\item \hyperref[restricted-section-phantom]{Algébrisation des espaces formels}
\item \hyperref[spaces-resolve-section-phantom]{Résolution des surfaces revisitée}
\end{enumerate}
Théorie des déformations
\begin{enumerate}
\setcounter{enumi}{89}
\item \hyperref[formal-defos-section-phantom]{Théorie formelle des déformations}
\item \hyperref[defos-section-phantom]{Théorie des déformations}
\item \hyperref[cotangent-section-phantom]{Le complexe cotangent}
\item \hyperref[examples-defos-section-phantom]{Problèmes de déformation}
\end{enumerate}
Champs algébriques
\begin{enumerate}
\setcounter{enumi}{93}
\item \hyperref[algebraic-section-phantom]{Champs algébriques}
\item \hyperref[examples-stacks-section-phantom]{Exemples de champs}
\item \hyperref[stacks-sheaves-section-phantom]{Faisceaux sur les champs algébriques}
\item \hyperref[criteria-section-phantom]{Critères de représentabilité}
\item \hyperref[artin-section-phantom]{Axiomes d'Artin}
\item \hyperref[quot-section-phantom]{Espaces de Quot et de Hilbert}
\item \hyperref[stacks-properties-section-phantom]{Propriétés des champs algébriques}
\item \hyperref[stacks-morphisms-section-phantom]{Morphismes de champs algébriques}
\item \hyperref[stacks-limits-section-phantom]{Limites de champs algébriques}
\item \hyperref[stacks-cohomology-section-phantom]{Cohomologie des champs algébriques}
\item \hyperref[stacks-perfect-section-phantom]{Catégories dérivées de champs}
\item \hyperref[stacks-introduction-section-phantom]{Introduction aux champs algébriques}
\item \hyperref[stacks-more-morphisms-section-phantom]{Compléments sur les morphismes de champs}
\item \hyperref[stacks-geometry-section-phantom]{La géométrie des champs}
\end{enumerate}
Thèmes de théorie des espaces de modules
\begin{enumerate}
\setcounter{enumi}{107}
\item \hyperref[moduli-section-phantom]{Champs de modules}
\item \hyperref[moduli-curves-section-phantom]{Espaces de modules de courbes}
\end{enumerate}
Divers
\begin{enumerate}
\setcounter{enumi}{109}
\item \hyperref[examples-section-phantom]{Exemples}
\item \hyperref[exercises-section-phantom]{Exercices}
\item \hyperref[guide-section-phantom]{Guide bibliographique}
\item \hyperref[desirables-section-phantom]{Desiderata}
\item \hyperref[coding-section-phantom]{Style de programmation}
\item \hyperref[obsolete-section-phantom]{Obsolète}
\item \hyperref[fdl-section-phantom]{Licence de documentation libre GNU}
\item \hyperref[index-section-phantom]{Index généré automatiquement}
\end{enumerate}
\end{multicols}


\bibliography{my}
\bibliographystyle{amsalpha}

\end{document}
